%% Generated by Sphinx.
\def\sphinxdocclass{jupyterBook}
\documentclass[letterpaper,10pt,english]{jupyterBook}
\ifdefined\pdfpxdimen
   \let\sphinxpxdimen\pdfpxdimen\else\newdimen\sphinxpxdimen
\fi \sphinxpxdimen=.75bp\relax
\ifdefined\pdfimageresolution
    \pdfimageresolution= \numexpr \dimexpr1in\relax/\sphinxpxdimen\relax
\fi
%% let collapsible pdf bookmarks panel have high depth per default
\PassOptionsToPackage{bookmarksdepth=5}{hyperref}
%% turn off hyperref patch of \index as sphinx.xdy xindy module takes care of
%% suitable \hyperpage mark-up, working around hyperref-xindy incompatibility
\PassOptionsToPackage{hyperindex=false}{hyperref}
%% memoir class requires extra handling
\makeatletter\@ifclassloaded{memoir}
{\ifdefined\memhyperindexfalse\memhyperindexfalse\fi}{}\makeatother

\PassOptionsToPackage{booktabs}{sphinx}
\PassOptionsToPackage{colorrows}{sphinx}

\PassOptionsToPackage{warn}{textcomp}

\catcode`^^^^00a0\active\protected\def^^^^00a0{\leavevmode\nobreak\ }
\usepackage{cmap}
\usepackage{fontspec}
\defaultfontfeatures[\rmfamily,\sffamily,\ttfamily]{}
\usepackage{amsmath,amssymb,amstext}
\usepackage{polyglossia}
\setmainlanguage{english}



\setmainfont{FreeSerif}[
  Extension      = .otf,
  UprightFont    = *,
  ItalicFont     = *Italic,
  BoldFont       = *Bold,
  BoldItalicFont = *BoldItalic
]
\setsansfont{FreeSans}[
  Extension      = .otf,
  UprightFont    = *,
  ItalicFont     = *Oblique,
  BoldFont       = *Bold,
  BoldItalicFont = *BoldOblique,
]
\setmonofont{FreeMono}[
  Extension      = .otf,
  UprightFont    = *,
  ItalicFont     = *Oblique,
  BoldFont       = *Bold,
  BoldItalicFont = *BoldOblique,
]



\usepackage[Bjarne]{fncychap}
\usepackage[,numfigreset=1,mathnumfig]{sphinx}

\fvset{fontsize=\small}
\usepackage{geometry}


% Include hyperref last.
\usepackage{hyperref}
% Fix anchor placement for figures with captions.
\usepackage{hypcap}% it must be loaded after hyperref.
% Set up styles of URL: it should be placed after hyperref.
\urlstyle{same}


\usepackage{sphinxmessages}



        % Start of preamble defined in sphinx-jupyterbook-latex %
         \usepackage[Latin,Greek]{ucharclasses}
        \usepackage{unicode-math}
        % fixing title of the toc
        \addto\captionsenglish{\renewcommand{\contentsname}{Contents}}
        \hypersetup{
            pdfencoding=auto,
            psdextra
        }
        % End of preamble defined in sphinx-jupyterbook-latex %
        

\title{Fluid Mechanics Class Notes}
\date{Oct 22, 2024}
\release{}
\author{Carlos D.\@{} Hoyos}
\newcommand{\sphinxlogo}{\vbox{}}
\renewcommand{\releasename}{}
\makeindex
\begin{document}

\pagestyle{empty}
\sphinxmaketitle
\pagestyle{plain}
\sphinxtableofcontents
\pagestyle{normal}
\phantomsection\label{\detokenize{intro::doc}}


\sphinxAtStartPar
My name is \sphinxstylestrong{Carlos D. Hoyos}, and I am one of the professors at the \sphinxstylestrong{Universidad Nacional de Colombia, Facultad de Minas, Sede Medellín}, guiding you through the principles and applications of fluid mechanics. As part of this journey, I am committed to building a set of online course materials to complement our in\sphinxhyphen{}class lectures. This content is made accessible through Jupyter Books hosted on GitHub, providing you with dynamic, interactive, and up\sphinxhyphen{}to\sphinxhyphen{}date resources for your studies.

\begin{DUlineblock}{0em}
\item[] \sphinxstylestrong{\Large The Importance of Fluid Mechanics}
\end{DUlineblock}

\sphinxAtStartPar
Fluid mechanics is not just a core subject within engineering; it’s the backbone of many fields that influence our daily lives. Whether you’re dealing with the flow of air over an aircraft wing, the circulation of blood in the human body, or the dynamics of ocean currents driving global climate patterns, fluid mechanics is involved. The discipline provides the tools to understand how fluids (liquids and gases) behave, whether they are at rest or in motion, and how these behaviors interact with the forces and energy around them.

\sphinxAtStartPar
From an engineering perspective, fluid mechanics plays a crucial role in \sphinxstylestrong{energy systems}, where it helps us design more efficient turbines, optimize hydraulic systems, and improve aerodynamics in transportation. \sphinxstylestrong{Environmental applications} are just as critical. For instance, fluid mechanics is at the heart of \sphinxstylestrong{water resource management}—whether we are designing irrigation systems, controlling floodwaters, or managing the quality of water bodies like rivers and lakes.

\sphinxAtStartPar
One of the most exciting areas where fluid mechanics comes into play is in \sphinxstylestrong{climate science}. As an engineer with a background in both climate and fluid dynamics, I can attest that understanding how fluids move in the atmosphere and oceans is fundamental to making accurate climate predictions. Large\sphinxhyphen{}scale systems like \sphinxstylestrong{jet streams, ocean currents, and monsoons} are governed by the principles of fluid mechanics. These systems, in turn, influence everything from \sphinxstylestrong{local weather patterns} to \sphinxstylestrong{long\sphinxhyphen{}term climate trends} that affect agriculture, ecosystems, and even energy generation.

\sphinxAtStartPar
At the other end of the spectrum, fluid mechanics also helps solve \sphinxstylestrong{small\sphinxhyphen{}scale, everyday problems}. Think about how fluids are used in \sphinxstylestrong{biomedical engineering}—from the way drugs are delivered through microfluidic devices to how engineers design heart pumps and blood flow models. Even at this smaller scale, the principles we study in this course apply and enable breakthroughs in technology and medicine.

\sphinxAtStartPar
Understanding fluid mechanics empowers us to address these diverse challenges, giving engineers a universal toolkit to approach problems ranging from global to local scales.

\begin{DUlineblock}{0em}
\item[] \sphinxstylestrong{\large Fluid Mechanics and Climate Science}
\end{DUlineblock}

\sphinxAtStartPar
While we often associate fluid mechanics with mechanical and civil engineering, its relevance to \sphinxstylestrong{climate science} cannot be overstated. Fluid behavior in the atmosphere and oceans dictates weather patterns, the distribution of heat and moisture across the planet, and even how pollutants disperse. As the climate changes, fluid mechanics becomes crucial in modeling these evolving dynamics and predicting their impacts on both natural and human systems.

\sphinxAtStartPar
For example, \sphinxstylestrong{wind energy}, one of the most promising renewable energy sources, relies heavily on our understanding of fluid flow. Engineers design wind turbines and wind farms by modeling how air moves across landscapes. Similarly, in \sphinxstylestrong{solar energy}, fluid dynamics helps in improving the cooling systems of solar panels and optimizing the flow of air around large solar fields.

\sphinxAtStartPar
Moreover, \sphinxstylestrong{climate resilience}—a topic that is increasingly critical as we face more extreme weather events—benefits directly from fluid mechanics. Flood control systems, stormwater management, and predicting the behavior of hurricanes or typhoons all require a deep understanding of how fluids move and interact with their surroundings.

\begin{DUlineblock}{0em}
\item[] \sphinxstylestrong{\large Small\sphinxhyphen{}Scale Applications: From Biomedical Engineering to Household Technologies}
\end{DUlineblock}

\sphinxAtStartPar
While climate science deals with global fluid systems, the principles of fluid mechanics also apply at much smaller scales. \sphinxstylestrong{Biomedical engineering} is a field that deeply integrates fluid dynamics to design medical devices, simulate blood flow, and develop life\sphinxhyphen{}saving technologies. Engineers use fluid mechanics to understand how the heart pumps blood, how artificial valves can mimic natural flow, and even how we might deliver drugs through targeted microfluidic channels.

\sphinxAtStartPar
Even in everyday household technologies, fluid mechanics plays a role—whether it’s the way your \sphinxstylestrong{showerhead} disperses water efficiently or how a \sphinxstylestrong{vacuum cleaner} is designed to create suction through fluid flow. These seemingly simple systems are based on complex principles that we will explore in this course.

\begin{DUlineblock}{0em}
\item[] \sphinxstylestrong{\large Why Jupyter Books?}
\end{DUlineblock}

\sphinxAtStartPar
Jupyter Books provide a unique platform for integrating live code, dynamic visualizations, and interactive examples directly into our course materials. This allows us to break down complex concepts and equations in real\sphinxhyphen{}time, giving you the opportunity to not only read but also engage with the material. Imagine being able to \sphinxstylestrong{run a simulation} of a fluid system while learning the theory behind it, or exploring \sphinxstylestrong{interactive visualizations} that help you better grasp the forces at play in different fluid scenarios.

\sphinxAtStartPar
The flexibility of Jupyter Books makes them an ideal medium for fluid mechanics, a subject that often requires \sphinxstylestrong{hands\sphinxhyphen{}on analysis} and \sphinxstylestrong{computational tools} to fully understand. As we progress through the course, I encourage you to experiment with the code and simulations provided in these materials to deepen your learning. You can learn about Jupyter Books at \sphinxhref{https://jupyterbook.org}{jupyterbook.org}.

\begin{DUlineblock}{0em}
\item[] \sphinxstylestrong{\large Why GitHub?}
\end{DUlineblock}

\sphinxAtStartPar
By hosting these materials on GitHub, I’m inviting you into a \sphinxstylestrong{collaborative learning environment}. Version control on GitHub ensures that the materials are always up to date, and it offers an opportunity for you to \sphinxstylestrong{suggest improvements}, raise issues, and even contribute to new content. This platform promotes transparency, community\sphinxhyphen{}driven learning, and continuous improvement—values that are essential in both academia and engineering practice.

\sphinxAtStartPar
If you don’t already have a GitHub account, I recommend creating one, as it will be useful not only for this course but also for future projects and collaborations in your career.

\begin{DUlineblock}{0em}
\item[] \sphinxstylestrong{\large Why the Notes in English?}
\end{DUlineblock}

\sphinxAtStartPar
Although this course is taught in Spanish, the course notes are provided in English to help you get comfortable with the \sphinxstylestrong{technical terminology} that dominates international literature, research, and professional work in fluid mechanics. Many of the key resources—such as academic papers, research articles, and advanced textbooks—are published in English, and mastering this language will help you stay connected with the \sphinxstylestrong{global scientific community}.

\sphinxAtStartPar
By engaging with the material in English, you’ll be better prepared to attend \sphinxstylestrong{international conferences}, contribute to \sphinxstylestrong{global research efforts}, and work with colleagues from around the world. It’s an investment that will expand your opportunities both academically and professionally.

\begin{DUlineblock}{0em}
\item[] \sphinxstylestrong{\large A Bit About Me}
\end{DUlineblock}

\sphinxAtStartPar
I have a \sphinxstylestrong{Ph.D. in Earth and Atmospheric Sciences from Georgia Institute of Technology}, along with a \sphinxstylestrong{Masters in Oceanic and Atmospheric Sciences}, a \sphinxstylestrong{Masters in Water Resources}, and a \sphinxstylestrong{Civil Engineering degree}. My work has taken me through diverse fields—studying the behavior of \sphinxstylestrong{atmospheric currents}, understanding \sphinxstylestrong{water resources}, and applying \sphinxstylestrong{data science} to solve problems related to agriculture, energy, risk management, and climate.

\sphinxAtStartPar
Fluid dynamics has been a common thread throughout my career, from analyzing \sphinxstylestrong{atmospheric circulations} that drive global climate patterns to modeling \sphinxstylestrong{river flows} that sustain ecosystems and cities. This course gives me the opportunity to share not only the fundamental concepts of fluid mechanics but also their real\sphinxhyphen{}world applications—whether in \sphinxstylestrong{renewable energy}, \sphinxstylestrong{environmental protection}, or \sphinxstylestrong{cutting\sphinxhyphen{}edge technologies}.

\sphinxstepscope


\chapter{Setting Up the Environment}
\label{\detokenize{environment:setting-up-the-environment}}\label{\detokenize{environment::doc}}
\sphinxAtStartPar
To replicate the exact Python environment I am using for this course, you can replicate the \sphinxcode{\sphinxupquote{environment.yml}} file or view its contents below.


\section{Environment File Contents}
\label{\detokenize{environment:environment-file-contents}}
\begin{sphinxVerbatim}[commandchars=\\\{\}]
\PYG{n+nt}{name}\PYG{p}{:}\PYG{+w}{ }\PYG{l+lScalar+lScalarPlain}{jupyterbook\PYGZus{}env}
\PYG{n+nt}{channels}\PYG{p}{:}
\PYG{+w}{  }\PYG{p+pIndicator}{\PYGZhy{}}\PYG{+w}{ }\PYG{l+lScalar+lScalarPlain}{defaults}
\PYG{n+nt}{dependencies}\PYG{p}{:}
\PYG{+w}{  }\PYG{p+pIndicator}{\PYGZhy{}}\PYG{+w}{ }\PYG{l+lScalar+lScalarPlain}{anyio=4.6.2=py311hecd8cb5\PYGZus{}0}
\PYG{+w}{  }\PYG{p+pIndicator}{\PYGZhy{}}\PYG{+w}{ }\PYG{l+lScalar+lScalarPlain}{appnope=0.1.2=py311hecd8cb5\PYGZus{}1001}
\PYG{+w}{  }\PYG{p+pIndicator}{\PYGZhy{}}\PYG{+w}{ }\PYG{l+lScalar+lScalarPlain}{argon2\PYGZhy{}cffi=21.3.0=pyhd3eb1b0\PYGZus{}0}
\PYG{+w}{  }\PYG{p+pIndicator}{\PYGZhy{}}\PYG{+w}{ }\PYG{l+lScalar+lScalarPlain}{argon2\PYGZhy{}cffi\PYGZhy{}bindings=21.2.0=py311h6c40b1e\PYGZus{}0}
\PYG{+w}{  }\PYG{p+pIndicator}{\PYGZhy{}}\PYG{+w}{ }\PYG{l+lScalar+lScalarPlain}{asttokens=2.0.5=pyhd3eb1b0\PYGZus{}0}
\PYG{+w}{  }\PYG{p+pIndicator}{\PYGZhy{}}\PYG{+w}{ }\PYG{l+lScalar+lScalarPlain}{async\PYGZhy{}lru=2.0.4=py311hecd8cb5\PYGZus{}0}
\PYG{+w}{  }\PYG{p+pIndicator}{\PYGZhy{}}\PYG{+w}{ }\PYG{l+lScalar+lScalarPlain}{attrs=24.2.0=py311hecd8cb5\PYGZus{}0}
\PYG{+w}{  }\PYG{p+pIndicator}{\PYGZhy{}}\PYG{+w}{ }\PYG{l+lScalar+lScalarPlain}{beautifulsoup4=4.12.3=py311hecd8cb5\PYGZus{}0}
\PYG{+w}{  }\PYG{p+pIndicator}{\PYGZhy{}}\PYG{+w}{ }\PYG{l+lScalar+lScalarPlain}{blas=1.0=mkl}
\PYG{+w}{  }\PYG{p+pIndicator}{\PYGZhy{}}\PYG{+w}{ }\PYG{l+lScalar+lScalarPlain}{bleach=4.1.0=pyhd3eb1b0\PYGZus{}0}
\PYG{+w}{  }\PYG{p+pIndicator}{\PYGZhy{}}\PYG{+w}{ }\PYG{l+lScalar+lScalarPlain}{bottleneck=1.3.7=py311hb3a5e46\PYGZus{}0}
\PYG{+w}{  }\PYG{p+pIndicator}{\PYGZhy{}}\PYG{+w}{ }\PYG{l+lScalar+lScalarPlain}{brotli=1.0.9=h6c40b1e\PYGZus{}8}
\PYG{+w}{  }\PYG{p+pIndicator}{\PYGZhy{}}\PYG{+w}{ }\PYG{l+lScalar+lScalarPlain}{brotli\PYGZhy{}bin=1.0.9=h6c40b1e\PYGZus{}8}
\PYG{+w}{  }\PYG{p+pIndicator}{\PYGZhy{}}\PYG{+w}{ }\PYG{l+lScalar+lScalarPlain}{brotli\PYGZhy{}python=1.0.9=py311hcec6c5f\PYGZus{}8}
\PYG{+w}{  }\PYG{p+pIndicator}{\PYGZhy{}}\PYG{+w}{ }\PYG{l+lScalar+lScalarPlain}{bzip2=1.0.8=h6c40b1e\PYGZus{}6}
\PYG{+w}{  }\PYG{p+pIndicator}{\PYGZhy{}}\PYG{+w}{ }\PYG{l+lScalar+lScalarPlain}{ca\PYGZhy{}certificates=2024.9.24=hecd8cb5\PYGZus{}0}
\PYG{+w}{  }\PYG{p+pIndicator}{\PYGZhy{}}\PYG{+w}{ }\PYG{l+lScalar+lScalarPlain}{certifi=2024.8.30=py311hecd8cb5\PYGZus{}0}
\PYG{+w}{  }\PYG{p+pIndicator}{\PYGZhy{}}\PYG{+w}{ }\PYG{l+lScalar+lScalarPlain}{cffi=1.17.1=py311h9205ec4\PYGZus{}0}
\PYG{+w}{  }\PYG{p+pIndicator}{\PYGZhy{}}\PYG{+w}{ }\PYG{l+lScalar+lScalarPlain}{charset\PYGZhy{}normalizer=3.3.2=pyhd3eb1b0\PYGZus{}0}
\PYG{+w}{  }\PYG{p+pIndicator}{\PYGZhy{}}\PYG{+w}{ }\PYG{l+lScalar+lScalarPlain}{comm=0.2.1=py311hecd8cb5\PYGZus{}0}
\PYG{+w}{  }\PYG{p+pIndicator}{\PYGZhy{}}\PYG{+w}{ }\PYG{l+lScalar+lScalarPlain}{contourpy=1.2.0=py311ha357a0b\PYGZus{}0}
\PYG{+w}{  }\PYG{p+pIndicator}{\PYGZhy{}}\PYG{+w}{ }\PYG{l+lScalar+lScalarPlain}{cycler=0.11.0=pyhd3eb1b0\PYGZus{}0}
\PYG{+w}{  }\PYG{p+pIndicator}{\PYGZhy{}}\PYG{+w}{ }\PYG{l+lScalar+lScalarPlain}{debugpy=1.6.7=py311hcec6c5f\PYGZus{}0}
\PYG{+w}{  }\PYG{p+pIndicator}{\PYGZhy{}}\PYG{+w}{ }\PYG{l+lScalar+lScalarPlain}{decorator=5.1.1=pyhd3eb1b0\PYGZus{}0}
\PYG{+w}{  }\PYG{p+pIndicator}{\PYGZhy{}}\PYG{+w}{ }\PYG{l+lScalar+lScalarPlain}{defusedxml=0.7.1=pyhd3eb1b0\PYGZus{}0}
\PYG{+w}{  }\PYG{p+pIndicator}{\PYGZhy{}}\PYG{+w}{ }\PYG{l+lScalar+lScalarPlain}{executing=0.8.3=pyhd3eb1b0\PYGZus{}0}
\PYG{+w}{  }\PYG{p+pIndicator}{\PYGZhy{}}\PYG{+w}{ }\PYG{l+lScalar+lScalarPlain}{fonttools=4.51.0=py311h6c40b1e\PYGZus{}0}
\PYG{+w}{  }\PYG{p+pIndicator}{\PYGZhy{}}\PYG{+w}{ }\PYG{l+lScalar+lScalarPlain}{freetype=2.12.1=hd8bbffd\PYGZus{}0}
\PYG{+w}{  }\PYG{p+pIndicator}{\PYGZhy{}}\PYG{+w}{ }\PYG{l+lScalar+lScalarPlain}{h11=0.14.0=py311hecd8cb5\PYGZus{}0}
\PYG{+w}{  }\PYG{p+pIndicator}{\PYGZhy{}}\PYG{+w}{ }\PYG{l+lScalar+lScalarPlain}{httpcore=1.0.2=py311hecd8cb5\PYGZus{}0}
\PYG{+w}{  }\PYG{p+pIndicator}{\PYGZhy{}}\PYG{+w}{ }\PYG{l+lScalar+lScalarPlain}{httpx=0.27.0=py311hecd8cb5\PYGZus{}0}
\PYG{+w}{  }\PYG{p+pIndicator}{\PYGZhy{}}\PYG{+w}{ }\PYG{l+lScalar+lScalarPlain}{idna=3.7=py311hecd8cb5\PYGZus{}0}
\PYG{+w}{  }\PYG{p+pIndicator}{\PYGZhy{}}\PYG{+w}{ }\PYG{l+lScalar+lScalarPlain}{intel\PYGZhy{}openmp=2023.1.0=ha357a0b\PYGZus{}43548}
\PYG{+w}{  }\PYG{p+pIndicator}{\PYGZhy{}}\PYG{+w}{ }\PYG{l+lScalar+lScalarPlain}{ipykernel=6.29.5=py311hecd8cb5\PYGZus{}0}
\PYG{+w}{  }\PYG{p+pIndicator}{\PYGZhy{}}\PYG{+w}{ }\PYG{l+lScalar+lScalarPlain}{ipython=8.27.0=py311hecd8cb5\PYGZus{}0}
\PYG{+w}{  }\PYG{p+pIndicator}{\PYGZhy{}}\PYG{+w}{ }\PYG{l+lScalar+lScalarPlain}{jedi=0.19.1=py311hecd8cb5\PYGZus{}0}
\PYG{+w}{  }\PYG{p+pIndicator}{\PYGZhy{}}\PYG{+w}{ }\PYG{l+lScalar+lScalarPlain}{jinja2=3.1.4=py311hecd8cb5\PYGZus{}0}
\PYG{+w}{  }\PYG{p+pIndicator}{\PYGZhy{}}\PYG{+w}{ }\PYG{l+lScalar+lScalarPlain}{jpeg=9e=h46256e1\PYGZus{}3}
\PYG{+w}{  }\PYG{p+pIndicator}{\PYGZhy{}}\PYG{+w}{ }\PYG{l+lScalar+lScalarPlain}{json5=0.9.6=pyhd3eb1b0\PYGZus{}0}
\PYG{+w}{  }\PYG{p+pIndicator}{\PYGZhy{}}\PYG{+w}{ }\PYG{l+lScalar+lScalarPlain}{jsonschema=4.23.0=py311hecd8cb5\PYGZus{}0}
\PYG{+w}{  }\PYG{p+pIndicator}{\PYGZhy{}}\PYG{+w}{ }\PYG{l+lScalar+lScalarPlain}{jsonschema\PYGZhy{}specifications=2023.7.1=py311hecd8cb5\PYGZus{}0}
\PYG{+w}{  }\PYG{p+pIndicator}{\PYGZhy{}}\PYG{+w}{ }\PYG{l+lScalar+lScalarPlain}{jupyter\PYGZhy{}lsp=2.2.0=py311hecd8cb5\PYGZus{}0}
\PYG{+w}{  }\PYG{p+pIndicator}{\PYGZhy{}}\PYG{+w}{ }\PYG{l+lScalar+lScalarPlain}{jupyter\PYGZus{}client=8.6.0=py311hecd8cb5\PYGZus{}0}
\PYG{+w}{  }\PYG{p+pIndicator}{\PYGZhy{}}\PYG{+w}{ }\PYG{l+lScalar+lScalarPlain}{jupyter\PYGZus{}core=5.7.2=py311hecd8cb5\PYGZus{}0}
\PYG{+w}{  }\PYG{p+pIndicator}{\PYGZhy{}}\PYG{+w}{ }\PYG{l+lScalar+lScalarPlain}{jupyter\PYGZus{}events=0.10.0=py311hecd8cb5\PYGZus{}0}
\PYG{+w}{  }\PYG{p+pIndicator}{\PYGZhy{}}\PYG{+w}{ }\PYG{l+lScalar+lScalarPlain}{jupyter\PYGZus{}server=2.14.1=py311hecd8cb5\PYGZus{}0}
\PYG{+w}{  }\PYG{p+pIndicator}{\PYGZhy{}}\PYG{+w}{ }\PYG{l+lScalar+lScalarPlain}{jupyter\PYGZus{}server\PYGZus{}terminals=0.4.4=py311hecd8cb5\PYGZus{}1}
\PYG{+w}{  }\PYG{p+pIndicator}{\PYGZhy{}}\PYG{+w}{ }\PYG{l+lScalar+lScalarPlain}{jupyterlab=4.2.5=py311hecd8cb5\PYGZus{}0}
\PYG{+w}{  }\PYG{p+pIndicator}{\PYGZhy{}}\PYG{+w}{ }\PYG{l+lScalar+lScalarPlain}{jupyterlab\PYGZus{}pygments=0.1.2=py\PYGZus{}0}
\PYG{+w}{  }\PYG{p+pIndicator}{\PYGZhy{}}\PYG{+w}{ }\PYG{l+lScalar+lScalarPlain}{jupyterlab\PYGZus{}server=2.27.3=py311hecd8cb5\PYGZus{}0}
\PYG{+w}{  }\PYG{p+pIndicator}{\PYGZhy{}}\PYG{+w}{ }\PYG{l+lScalar+lScalarPlain}{kiwisolver=1.4.4=py311hcec6c5f\PYGZus{}0}
\PYG{+w}{  }\PYG{p+pIndicator}{\PYGZhy{}}\PYG{+w}{ }\PYG{l+lScalar+lScalarPlain}{lcms2=2.12=hf1fd2bf\PYGZus{}0}
\PYG{+w}{  }\PYG{p+pIndicator}{\PYGZhy{}}\PYG{+w}{ }\PYG{l+lScalar+lScalarPlain}{lerc=3.0=he9d5cce\PYGZus{}0}
\PYG{+w}{  }\PYG{p+pIndicator}{\PYGZhy{}}\PYG{+w}{ }\PYG{l+lScalar+lScalarPlain}{libbrotlicommon=1.0.9=h6c40b1e\PYGZus{}8}
\PYG{+w}{  }\PYG{p+pIndicator}{\PYGZhy{}}\PYG{+w}{ }\PYG{l+lScalar+lScalarPlain}{libbrotlidec=1.0.9=h6c40b1e\PYGZus{}8}
\PYG{+w}{  }\PYG{p+pIndicator}{\PYGZhy{}}\PYG{+w}{ }\PYG{l+lScalar+lScalarPlain}{libbrotlienc=1.0.9=h6c40b1e\PYGZus{}8}
\PYG{+w}{  }\PYG{p+pIndicator}{\PYGZhy{}}\PYG{+w}{ }\PYG{l+lScalar+lScalarPlain}{libcxx=14.0.6=h9765a3e\PYGZus{}0}
\PYG{+w}{  }\PYG{p+pIndicator}{\PYGZhy{}}\PYG{+w}{ }\PYG{l+lScalar+lScalarPlain}{libdeflate=1.17=hb664fd8\PYGZus{}1}
\PYG{+w}{  }\PYG{p+pIndicator}{\PYGZhy{}}\PYG{+w}{ }\PYG{l+lScalar+lScalarPlain}{libffi=3.4.4=hecd8cb5\PYGZus{}1}
\PYG{+w}{  }\PYG{p+pIndicator}{\PYGZhy{}}\PYG{+w}{ }\PYG{l+lScalar+lScalarPlain}{libgfortran=5.0.0=11\PYGZus{}3\PYGZus{}0\PYGZus{}hecd8cb5\PYGZus{}28}
\PYG{+w}{  }\PYG{p+pIndicator}{\PYGZhy{}}\PYG{+w}{ }\PYG{l+lScalar+lScalarPlain}{libgfortran5=11.3.0=h9dfd629\PYGZus{}28}
\PYG{+w}{  }\PYG{p+pIndicator}{\PYGZhy{}}\PYG{+w}{ }\PYG{l+lScalar+lScalarPlain}{libpng=1.6.39=h6c40b1e\PYGZus{}0}
\PYG{+w}{  }\PYG{p+pIndicator}{\PYGZhy{}}\PYG{+w}{ }\PYG{l+lScalar+lScalarPlain}{libsodium=1.0.18=h1de35cc\PYGZus{}0}
\PYG{+w}{  }\PYG{p+pIndicator}{\PYGZhy{}}\PYG{+w}{ }\PYG{l+lScalar+lScalarPlain}{libtiff=4.5.1=hcec6c5f\PYGZus{}0}
\PYG{+w}{  }\PYG{p+pIndicator}{\PYGZhy{}}\PYG{+w}{ }\PYG{l+lScalar+lScalarPlain}{libwebp\PYGZhy{}base=1.3.2=h46256e1\PYGZus{}1}
\PYG{+w}{  }\PYG{p+pIndicator}{\PYGZhy{}}\PYG{+w}{ }\PYG{l+lScalar+lScalarPlain}{llvm\PYGZhy{}openmp=14.0.6=h0dcd299\PYGZus{}0}
\PYG{+w}{  }\PYG{p+pIndicator}{\PYGZhy{}}\PYG{+w}{ }\PYG{l+lScalar+lScalarPlain}{lz4\PYGZhy{}c=1.9.4=hcec6c5f\PYGZus{}1}
\PYG{+w}{  }\PYG{p+pIndicator}{\PYGZhy{}}\PYG{+w}{ }\PYG{l+lScalar+lScalarPlain}{markupsafe=2.1.3=py311h6c40b1e\PYGZus{}0}
\PYG{+w}{  }\PYG{p+pIndicator}{\PYGZhy{}}\PYG{+w}{ }\PYG{l+lScalar+lScalarPlain}{matplotlib=3.9.2=py311hecd8cb5\PYGZus{}0}
\PYG{+w}{  }\PYG{p+pIndicator}{\PYGZhy{}}\PYG{+w}{ }\PYG{l+lScalar+lScalarPlain}{matplotlib\PYGZhy{}base=3.9.2=py311hc25a08b\PYGZus{}0}
\PYG{+w}{  }\PYG{p+pIndicator}{\PYGZhy{}}\PYG{+w}{ }\PYG{l+lScalar+lScalarPlain}{matplotlib\PYGZhy{}inline=0.1.6=py311hecd8cb5\PYGZus{}0}
\PYG{+w}{  }\PYG{p+pIndicator}{\PYGZhy{}}\PYG{+w}{ }\PYG{l+lScalar+lScalarPlain}{mistune=2.0.4=py311hecd8cb5\PYGZus{}0}
\PYG{+w}{  }\PYG{p+pIndicator}{\PYGZhy{}}\PYG{+w}{ }\PYG{l+lScalar+lScalarPlain}{mkl=2023.1.0=h8e150cf\PYGZus{}43560}
\PYG{+w}{  }\PYG{p+pIndicator}{\PYGZhy{}}\PYG{+w}{ }\PYG{l+lScalar+lScalarPlain}{mkl\PYGZhy{}service=2.4.0=py311h6c40b1e\PYGZus{}1}
\PYG{+w}{  }\PYG{p+pIndicator}{\PYGZhy{}}\PYG{+w}{ }\PYG{l+lScalar+lScalarPlain}{mkl\PYGZus{}fft=1.3.8=py311h6c40b1e\PYGZus{}0}
\PYG{+w}{  }\PYG{p+pIndicator}{\PYGZhy{}}\PYG{+w}{ }\PYG{l+lScalar+lScalarPlain}{mkl\PYGZus{}random=1.2.4=py311ha357a0b\PYGZus{}0}
\PYG{+w}{  }\PYG{p+pIndicator}{\PYGZhy{}}\PYG{+w}{ }\PYG{l+lScalar+lScalarPlain}{nbclient=0.8.0=py311hecd8cb5\PYGZus{}0}
\PYG{+w}{  }\PYG{p+pIndicator}{\PYGZhy{}}\PYG{+w}{ }\PYG{l+lScalar+lScalarPlain}{nbconvert=7.16.4=py311hecd8cb5\PYGZus{}0}
\PYG{+w}{  }\PYG{p+pIndicator}{\PYGZhy{}}\PYG{+w}{ }\PYG{l+lScalar+lScalarPlain}{nbformat=5.10.4=py311hecd8cb5\PYGZus{}0}
\PYG{+w}{  }\PYG{p+pIndicator}{\PYGZhy{}}\PYG{+w}{ }\PYG{l+lScalar+lScalarPlain}{ncurses=6.4=hcec6c5f\PYGZus{}0}
\PYG{+w}{  }\PYG{p+pIndicator}{\PYGZhy{}}\PYG{+w}{ }\PYG{l+lScalar+lScalarPlain}{nest\PYGZhy{}asyncio=1.6.0=py311hecd8cb5\PYGZus{}0}
\PYG{+w}{  }\PYG{p+pIndicator}{\PYGZhy{}}\PYG{+w}{ }\PYG{l+lScalar+lScalarPlain}{notebook=7.2.2=py311hecd8cb5\PYGZus{}1}
\PYG{+w}{  }\PYG{p+pIndicator}{\PYGZhy{}}\PYG{+w}{ }\PYG{l+lScalar+lScalarPlain}{notebook\PYGZhy{}shim=0.2.3=py311hecd8cb5\PYGZus{}0}
\PYG{+w}{  }\PYG{p+pIndicator}{\PYGZhy{}}\PYG{+w}{ }\PYG{l+lScalar+lScalarPlain}{numexpr=2.8.7=py311h728a8a3\PYGZus{}0}
\PYG{+w}{  }\PYG{p+pIndicator}{\PYGZhy{}}\PYG{+w}{ }\PYG{l+lScalar+lScalarPlain}{numpy=1.26.4=py311h728a8a3\PYGZus{}0}
\PYG{+w}{  }\PYG{p+pIndicator}{\PYGZhy{}}\PYG{+w}{ }\PYG{l+lScalar+lScalarPlain}{numpy\PYGZhy{}base=1.26.4=py311h53bf9ac\PYGZus{}0}
\PYG{+w}{  }\PYG{p+pIndicator}{\PYGZhy{}}\PYG{+w}{ }\PYG{l+lScalar+lScalarPlain}{openjpeg=2.5.2=hbf2204d\PYGZus{}0}
\PYG{+w}{  }\PYG{p+pIndicator}{\PYGZhy{}}\PYG{+w}{ }\PYG{l+lScalar+lScalarPlain}{openssl=3.0.15=h46256e1\PYGZus{}0}
\PYG{+w}{  }\PYG{p+pIndicator}{\PYGZhy{}}\PYG{+w}{ }\PYG{l+lScalar+lScalarPlain}{overrides=7.4.0=py311hecd8cb5\PYGZus{}0}
\PYG{+w}{  }\PYG{p+pIndicator}{\PYGZhy{}}\PYG{+w}{ }\PYG{l+lScalar+lScalarPlain}{packaging=24.1=py311hecd8cb5\PYGZus{}0}
\PYG{+w}{  }\PYG{p+pIndicator}{\PYGZhy{}}\PYG{+w}{ }\PYG{l+lScalar+lScalarPlain}{pandas=2.2.2=py311he327ffe\PYGZus{}0}
\PYG{+w}{  }\PYG{p+pIndicator}{\PYGZhy{}}\PYG{+w}{ }\PYG{l+lScalar+lScalarPlain}{pandocfilters=1.5.0=pyhd3eb1b0\PYGZus{}0}
\PYG{+w}{  }\PYG{p+pIndicator}{\PYGZhy{}}\PYG{+w}{ }\PYG{l+lScalar+lScalarPlain}{parso=0.8.3=pyhd3eb1b0\PYGZus{}0}
\PYG{+w}{  }\PYG{p+pIndicator}{\PYGZhy{}}\PYG{+w}{ }\PYG{l+lScalar+lScalarPlain}{pexpect=4.8.0=pyhd3eb1b0\PYGZus{}3}
\PYG{+w}{  }\PYG{p+pIndicator}{\PYGZhy{}}\PYG{+w}{ }\PYG{l+lScalar+lScalarPlain}{pillow=10.4.0=py311h46256e1\PYGZus{}0}
\PYG{+w}{  }\PYG{p+pIndicator}{\PYGZhy{}}\PYG{+w}{ }\PYG{l+lScalar+lScalarPlain}{pip=24.2=py311hecd8cb5\PYGZus{}0}
\PYG{+w}{  }\PYG{p+pIndicator}{\PYGZhy{}}\PYG{+w}{ }\PYG{l+lScalar+lScalarPlain}{platformdirs=3.10.0=py311hecd8cb5\PYGZus{}0}
\PYG{+w}{  }\PYG{p+pIndicator}{\PYGZhy{}}\PYG{+w}{ }\PYG{l+lScalar+lScalarPlain}{prometheus\PYGZus{}client=0.14.1=py311hecd8cb5\PYGZus{}0}
\PYG{+w}{  }\PYG{p+pIndicator}{\PYGZhy{}}\PYG{+w}{ }\PYG{l+lScalar+lScalarPlain}{prompt\PYGZhy{}toolkit=3.0.43=py311hecd8cb5\PYGZus{}0}
\PYG{+w}{  }\PYG{p+pIndicator}{\PYGZhy{}}\PYG{+w}{ }\PYG{l+lScalar+lScalarPlain}{prompt\PYGZus{}toolkit=3.0.43=hd3eb1b0\PYGZus{}0}
\PYG{+w}{  }\PYG{p+pIndicator}{\PYGZhy{}}\PYG{+w}{ }\PYG{l+lScalar+lScalarPlain}{psutil=5.9.0=py311h6c40b1e\PYGZus{}0}
\PYG{+w}{  }\PYG{p+pIndicator}{\PYGZhy{}}\PYG{+w}{ }\PYG{l+lScalar+lScalarPlain}{ptyprocess=0.7.0=pyhd3eb1b0\PYGZus{}2}
\PYG{+w}{  }\PYG{p+pIndicator}{\PYGZhy{}}\PYG{+w}{ }\PYG{l+lScalar+lScalarPlain}{pure\PYGZus{}eval=0.2.2=pyhd3eb1b0\PYGZus{}0}
\PYG{+w}{  }\PYG{p+pIndicator}{\PYGZhy{}}\PYG{+w}{ }\PYG{l+lScalar+lScalarPlain}{pycparser=2.21=pyhd3eb1b0\PYGZus{}0}
\PYG{+w}{  }\PYG{p+pIndicator}{\PYGZhy{}}\PYG{+w}{ }\PYG{l+lScalar+lScalarPlain}{pyparsing=3.1.2=py311hecd8cb5\PYGZus{}0}
\PYG{+w}{  }\PYG{p+pIndicator}{\PYGZhy{}}\PYG{+w}{ }\PYG{l+lScalar+lScalarPlain}{pysocks=1.7.1=py311hecd8cb5\PYGZus{}0}
\PYG{+w}{  }\PYG{p+pIndicator}{\PYGZhy{}}\PYG{+w}{ }\PYG{l+lScalar+lScalarPlain}{python=3.11.10=h4d6d9e5\PYGZus{}0}
\PYG{+w}{  }\PYG{p+pIndicator}{\PYGZhy{}}\PYG{+w}{ }\PYG{l+lScalar+lScalarPlain}{python\PYGZhy{}dateutil=2.9.0post0=py311hecd8cb5\PYGZus{}2}
\PYG{+w}{  }\PYG{p+pIndicator}{\PYGZhy{}}\PYG{+w}{ }\PYG{l+lScalar+lScalarPlain}{python\PYGZhy{}fastjsonschema=2.16.2=py311hecd8cb5\PYGZus{}0}
\PYG{+w}{  }\PYG{p+pIndicator}{\PYGZhy{}}\PYG{+w}{ }\PYG{l+lScalar+lScalarPlain}{python\PYGZhy{}json\PYGZhy{}logger=2.0.7=py311hecd8cb5\PYGZus{}0}
\PYG{+w}{  }\PYG{p+pIndicator}{\PYGZhy{}}\PYG{+w}{ }\PYG{l+lScalar+lScalarPlain}{python\PYGZhy{}tzdata=2023.3=pyhd3eb1b0\PYGZus{}0}
\PYG{+w}{  }\PYG{p+pIndicator}{\PYGZhy{}}\PYG{+w}{ }\PYG{l+lScalar+lScalarPlain}{pytz=2024.1=py311hecd8cb5\PYGZus{}0}
\PYG{+w}{  }\PYG{p+pIndicator}{\PYGZhy{}}\PYG{+w}{ }\PYG{l+lScalar+lScalarPlain}{pyyaml=6.0.2=py311h46256e1\PYGZus{}0}
\PYG{+w}{  }\PYG{p+pIndicator}{\PYGZhy{}}\PYG{+w}{ }\PYG{l+lScalar+lScalarPlain}{pyzmq=25.1.2=py311hcec6c5f\PYGZus{}0}
\PYG{+w}{  }\PYG{p+pIndicator}{\PYGZhy{}}\PYG{+w}{ }\PYG{l+lScalar+lScalarPlain}{readline=8.2=hca72f7f\PYGZus{}0}
\PYG{+w}{  }\PYG{p+pIndicator}{\PYGZhy{}}\PYG{+w}{ }\PYG{l+lScalar+lScalarPlain}{referencing=0.30.2=py311hecd8cb5\PYGZus{}0}
\PYG{+w}{  }\PYG{p+pIndicator}{\PYGZhy{}}\PYG{+w}{ }\PYG{l+lScalar+lScalarPlain}{requests=2.32.3=py311hecd8cb5\PYGZus{}0}
\PYG{+w}{  }\PYG{p+pIndicator}{\PYGZhy{}}\PYG{+w}{ }\PYG{l+lScalar+lScalarPlain}{rfc3339\PYGZhy{}validator=0.1.4=py311hecd8cb5\PYGZus{}0}
\PYG{+w}{  }\PYG{p+pIndicator}{\PYGZhy{}}\PYG{+w}{ }\PYG{l+lScalar+lScalarPlain}{rfc3986\PYGZhy{}validator=0.1.1=py311hecd8cb5\PYGZus{}0}
\PYG{+w}{  }\PYG{p+pIndicator}{\PYGZhy{}}\PYG{+w}{ }\PYG{l+lScalar+lScalarPlain}{rpds\PYGZhy{}py=0.10.6=py311hf2ad997\PYGZus{}0}
\PYG{+w}{  }\PYG{p+pIndicator}{\PYGZhy{}}\PYG{+w}{ }\PYG{l+lScalar+lScalarPlain}{scipy=1.12.0=py311h224febf\PYGZus{}0}
\PYG{+w}{  }\PYG{p+pIndicator}{\PYGZhy{}}\PYG{+w}{ }\PYG{l+lScalar+lScalarPlain}{send2trash=1.8.2=py311hecd8cb5\PYGZus{}0}
\PYG{+w}{  }\PYG{p+pIndicator}{\PYGZhy{}}\PYG{+w}{ }\PYG{l+lScalar+lScalarPlain}{setuptools=75.1.0=py311hecd8cb5\PYGZus{}0}
\PYG{+w}{  }\PYG{p+pIndicator}{\PYGZhy{}}\PYG{+w}{ }\PYG{l+lScalar+lScalarPlain}{six=1.16.0=pyhd3eb1b0\PYGZus{}1}
\PYG{+w}{  }\PYG{p+pIndicator}{\PYGZhy{}}\PYG{+w}{ }\PYG{l+lScalar+lScalarPlain}{sniffio=1.3.0=py311hecd8cb5\PYGZus{}0}
\PYG{+w}{  }\PYG{p+pIndicator}{\PYGZhy{}}\PYG{+w}{ }\PYG{l+lScalar+lScalarPlain}{soupsieve=2.5=py311hecd8cb5\PYGZus{}0}
\PYG{+w}{  }\PYG{p+pIndicator}{\PYGZhy{}}\PYG{+w}{ }\PYG{l+lScalar+lScalarPlain}{sqlite=3.45.3=h6c40b1e\PYGZus{}0}
\PYG{+w}{  }\PYG{p+pIndicator}{\PYGZhy{}}\PYG{+w}{ }\PYG{l+lScalar+lScalarPlain}{stack\PYGZus{}data=0.2.0=pyhd3eb1b0\PYGZus{}0}
\PYG{+w}{  }\PYG{p+pIndicator}{\PYGZhy{}}\PYG{+w}{ }\PYG{l+lScalar+lScalarPlain}{tbb=2021.8.0=ha357a0b\PYGZus{}0}
\PYG{+w}{  }\PYG{p+pIndicator}{\PYGZhy{}}\PYG{+w}{ }\PYG{l+lScalar+lScalarPlain}{terminado=0.17.1=py311hecd8cb5\PYGZus{}0}
\PYG{+w}{  }\PYG{p+pIndicator}{\PYGZhy{}}\PYG{+w}{ }\PYG{l+lScalar+lScalarPlain}{tinycss2=1.2.1=py311hecd8cb5\PYGZus{}0}
\PYG{+w}{  }\PYG{p+pIndicator}{\PYGZhy{}}\PYG{+w}{ }\PYG{l+lScalar+lScalarPlain}{tk=8.6.14=h4d00af3\PYGZus{}0}
\PYG{+w}{  }\PYG{p+pIndicator}{\PYGZhy{}}\PYG{+w}{ }\PYG{l+lScalar+lScalarPlain}{tornado=6.4.1=py311h46256e1\PYGZus{}0}
\PYG{+w}{  }\PYG{p+pIndicator}{\PYGZhy{}}\PYG{+w}{ }\PYG{l+lScalar+lScalarPlain}{traitlets=5.14.3=py311hecd8cb5\PYGZus{}0}
\PYG{+w}{  }\PYG{p+pIndicator}{\PYGZhy{}}\PYG{+w}{ }\PYG{l+lScalar+lScalarPlain}{typing\PYGZhy{}extensions=4.11.0=py311hecd8cb5\PYGZus{}0}
\PYG{+w}{  }\PYG{p+pIndicator}{\PYGZhy{}}\PYG{+w}{ }\PYG{l+lScalar+lScalarPlain}{typing\PYGZus{}extensions=4.11.0=py311hecd8cb5\PYGZus{}0}
\PYG{+w}{  }\PYG{p+pIndicator}{\PYGZhy{}}\PYG{+w}{ }\PYG{l+lScalar+lScalarPlain}{tzdata=2024b=h04d1e81\PYGZus{}0}
\PYG{+w}{  }\PYG{p+pIndicator}{\PYGZhy{}}\PYG{+w}{ }\PYG{l+lScalar+lScalarPlain}{unicodedata2=15.1.0=py311h6c40b1e\PYGZus{}0}
\PYG{+w}{  }\PYG{p+pIndicator}{\PYGZhy{}}\PYG{+w}{ }\PYG{l+lScalar+lScalarPlain}{urllib3=2.2.3=py311hecd8cb5\PYGZus{}0}
\PYG{+w}{  }\PYG{p+pIndicator}{\PYGZhy{}}\PYG{+w}{ }\PYG{l+lScalar+lScalarPlain}{wcwidth=0.2.5=pyhd3eb1b0\PYGZus{}0}
\PYG{+w}{  }\PYG{p+pIndicator}{\PYGZhy{}}\PYG{+w}{ }\PYG{l+lScalar+lScalarPlain}{webencodings=0.5.1=py311hecd8cb5\PYGZus{}1}
\PYG{+w}{  }\PYG{p+pIndicator}{\PYGZhy{}}\PYG{+w}{ }\PYG{l+lScalar+lScalarPlain}{websocket\PYGZhy{}client=1.8.0=py311hecd8cb5\PYGZus{}0}
\PYG{+w}{  }\PYG{p+pIndicator}{\PYGZhy{}}\PYG{+w}{ }\PYG{l+lScalar+lScalarPlain}{wheel=0.44.0=py311hecd8cb5\PYGZus{}0}
\PYG{+w}{  }\PYG{p+pIndicator}{\PYGZhy{}}\PYG{+w}{ }\PYG{l+lScalar+lScalarPlain}{xz=5.4.6=h6c40b1e\PYGZus{}1}
\PYG{+w}{  }\PYG{p+pIndicator}{\PYGZhy{}}\PYG{+w}{ }\PYG{l+lScalar+lScalarPlain}{yaml=0.2.5=haf1e3a3\PYGZus{}0}
\PYG{+w}{  }\PYG{p+pIndicator}{\PYGZhy{}}\PYG{+w}{ }\PYG{l+lScalar+lScalarPlain}{zeromq=4.3.5=hcec6c5f\PYGZus{}0}
\PYG{+w}{  }\PYG{p+pIndicator}{\PYGZhy{}}\PYG{+w}{ }\PYG{l+lScalar+lScalarPlain}{zlib=1.2.13=h4b97444\PYGZus{}1}
\PYG{+w}{  }\PYG{p+pIndicator}{\PYGZhy{}}\PYG{+w}{ }\PYG{l+lScalar+lScalarPlain}{zstd=1.5.6=h138b38a\PYGZus{}0}
\PYG{+w}{  }\PYG{p+pIndicator}{\PYGZhy{}}\PYG{+w}{ }\PYG{n+nt}{pip}\PYG{p}{:}
\PYG{+w}{    }\PYG{p+pIndicator}{\PYGZhy{}}\PYG{+w}{ }\PYG{l+lScalar+lScalarPlain}{accessible\PYGZhy{}pygments==0.0.5}
\PYG{+w}{    }\PYG{p+pIndicator}{\PYGZhy{}}\PYG{+w}{ }\PYG{l+lScalar+lScalarPlain}{alabaster==0.7.16}
\PYG{+w}{    }\PYG{p+pIndicator}{\PYGZhy{}}\PYG{+w}{ }\PYG{l+lScalar+lScalarPlain}{babel==2.16.0}
\PYG{+w}{    }\PYG{p+pIndicator}{\PYGZhy{}}\PYG{+w}{ }\PYG{l+lScalar+lScalarPlain}{click==8.1.7}
\PYG{+w}{    }\PYG{p+pIndicator}{\PYGZhy{}}\PYG{+w}{ }\PYG{l+lScalar+lScalarPlain}{docutils==0.20.1}
\PYG{+w}{    }\PYG{p+pIndicator}{\PYGZhy{}}\PYG{+w}{ }\PYG{l+lScalar+lScalarPlain}{greenlet==3.1.1}
\PYG{+w}{    }\PYG{p+pIndicator}{\PYGZhy{}}\PYG{+w}{ }\PYG{l+lScalar+lScalarPlain}{imagesize==1.4.1}
\PYG{+w}{    }\PYG{p+pIndicator}{\PYGZhy{}}\PYG{+w}{ }\PYG{l+lScalar+lScalarPlain}{importlib\PYGZhy{}metadata==8.5.0}
\PYG{+w}{    }\PYG{p+pIndicator}{\PYGZhy{}}\PYG{+w}{ }\PYG{l+lScalar+lScalarPlain}{jupyter\PYGZhy{}book==1.0.3}
\PYG{+w}{    }\PYG{p+pIndicator}{\PYGZhy{}}\PYG{+w}{ }\PYG{l+lScalar+lScalarPlain}{jupyter\PYGZhy{}cache==1.0.0}
\PYG{+w}{    }\PYG{p+pIndicator}{\PYGZhy{}}\PYG{+w}{ }\PYG{l+lScalar+lScalarPlain}{latexcodec==3.0.0}
\PYG{+w}{    }\PYG{p+pIndicator}{\PYGZhy{}}\PYG{+w}{ }\PYG{l+lScalar+lScalarPlain}{linkify\PYGZhy{}it\PYGZhy{}py==2.0.3}
\PYG{+w}{    }\PYG{p+pIndicator}{\PYGZhy{}}\PYG{+w}{ }\PYG{l+lScalar+lScalarPlain}{markdown\PYGZhy{}it\PYGZhy{}py==3.0.0}
\PYG{+w}{    }\PYG{p+pIndicator}{\PYGZhy{}}\PYG{+w}{ }\PYG{l+lScalar+lScalarPlain}{mdit\PYGZhy{}py\PYGZhy{}plugins==0.4.2}
\PYG{+w}{    }\PYG{p+pIndicator}{\PYGZhy{}}\PYG{+w}{ }\PYG{l+lScalar+lScalarPlain}{mdurl==0.1.2}
\PYG{+w}{    }\PYG{p+pIndicator}{\PYGZhy{}}\PYG{+w}{ }\PYG{l+lScalar+lScalarPlain}{myst\PYGZhy{}nb==1.1.2}
\PYG{+w}{    }\PYG{p+pIndicator}{\PYGZhy{}}\PYG{+w}{ }\PYG{l+lScalar+lScalarPlain}{myst\PYGZhy{}parser==2.0.0}
\PYG{+w}{    }\PYG{p+pIndicator}{\PYGZhy{}}\PYG{+w}{ }\PYG{l+lScalar+lScalarPlain}{pybtex==0.24.0}
\PYG{+w}{    }\PYG{p+pIndicator}{\PYGZhy{}}\PYG{+w}{ }\PYG{l+lScalar+lScalarPlain}{pybtex\PYGZhy{}docutils==1.0.3}
\PYG{+w}{    }\PYG{p+pIndicator}{\PYGZhy{}}\PYG{+w}{ }\PYG{l+lScalar+lScalarPlain}{pydata\PYGZhy{}sphinx\PYGZhy{}theme==0.15.4}
\PYG{+w}{    }\PYG{p+pIndicator}{\PYGZhy{}}\PYG{+w}{ }\PYG{l+lScalar+lScalarPlain}{pygments==2.18.0}
\PYG{+w}{    }\PYG{p+pIndicator}{\PYGZhy{}}\PYG{+w}{ }\PYG{l+lScalar+lScalarPlain}{snowballstemmer==2.2.0}
\PYG{+w}{    }\PYG{p+pIndicator}{\PYGZhy{}}\PYG{+w}{ }\PYG{l+lScalar+lScalarPlain}{sphinx==7.4.7}
\PYG{+w}{    }\PYG{p+pIndicator}{\PYGZhy{}}\PYG{+w}{ }\PYG{l+lScalar+lScalarPlain}{sphinx\PYGZhy{}book\PYGZhy{}theme==1.1.3}
\PYG{+w}{    }\PYG{p+pIndicator}{\PYGZhy{}}\PYG{+w}{ }\PYG{l+lScalar+lScalarPlain}{sphinx\PYGZhy{}comments==0.0.3}
\PYG{+w}{    }\PYG{p+pIndicator}{\PYGZhy{}}\PYG{+w}{ }\PYG{l+lScalar+lScalarPlain}{sphinx\PYGZhy{}copybutton==0.5.2}
\PYG{+w}{    }\PYG{p+pIndicator}{\PYGZhy{}}\PYG{+w}{ }\PYG{l+lScalar+lScalarPlain}{sphinx\PYGZhy{}design==0.6.1}
\PYG{+w}{    }\PYG{p+pIndicator}{\PYGZhy{}}\PYG{+w}{ }\PYG{l+lScalar+lScalarPlain}{sphinx\PYGZhy{}external\PYGZhy{}toc==1.0.1}
\PYG{+w}{    }\PYG{p+pIndicator}{\PYGZhy{}}\PYG{+w}{ }\PYG{l+lScalar+lScalarPlain}{sphinx\PYGZhy{}jupyterbook\PYGZhy{}latex==1.0.0}
\PYG{+w}{    }\PYG{p+pIndicator}{\PYGZhy{}}\PYG{+w}{ }\PYG{l+lScalar+lScalarPlain}{sphinx\PYGZhy{}multitoc\PYGZhy{}numbering==0.1.3}
\PYG{+w}{    }\PYG{p+pIndicator}{\PYGZhy{}}\PYG{+w}{ }\PYG{l+lScalar+lScalarPlain}{sphinx\PYGZhy{}thebe==0.3.1}
\PYG{+w}{    }\PYG{p+pIndicator}{\PYGZhy{}}\PYG{+w}{ }\PYG{l+lScalar+lScalarPlain}{sphinx\PYGZhy{}togglebutton==0.3.2}
\PYG{+w}{    }\PYG{p+pIndicator}{\PYGZhy{}}\PYG{+w}{ }\PYG{l+lScalar+lScalarPlain}{sphinxcontrib\PYGZhy{}applehelp==2.0.0}
\PYG{+w}{    }\PYG{p+pIndicator}{\PYGZhy{}}\PYG{+w}{ }\PYG{l+lScalar+lScalarPlain}{sphinxcontrib\PYGZhy{}bibtex==2.6.3}
\PYG{+w}{    }\PYG{p+pIndicator}{\PYGZhy{}}\PYG{+w}{ }\PYG{l+lScalar+lScalarPlain}{sphinxcontrib\PYGZhy{}devhelp==2.0.0}
\PYG{+w}{    }\PYG{p+pIndicator}{\PYGZhy{}}\PYG{+w}{ }\PYG{l+lScalar+lScalarPlain}{sphinxcontrib\PYGZhy{}htmlhelp==2.1.0}
\PYG{+w}{    }\PYG{p+pIndicator}{\PYGZhy{}}\PYG{+w}{ }\PYG{l+lScalar+lScalarPlain}{sphinxcontrib\PYGZhy{}jsmath==1.0.1}
\PYG{+w}{    }\PYG{p+pIndicator}{\PYGZhy{}}\PYG{+w}{ }\PYG{l+lScalar+lScalarPlain}{sphinxcontrib\PYGZhy{}qthelp==2.0.0}
\PYG{+w}{    }\PYG{p+pIndicator}{\PYGZhy{}}\PYG{+w}{ }\PYG{l+lScalar+lScalarPlain}{sphinxcontrib\PYGZhy{}serializinghtml==2.0.0}
\PYG{+w}{    }\PYG{p+pIndicator}{\PYGZhy{}}\PYG{+w}{ }\PYG{l+lScalar+lScalarPlain}{sqlalchemy==2.0.36}
\PYG{+w}{    }\PYG{p+pIndicator}{\PYGZhy{}}\PYG{+w}{ }\PYG{l+lScalar+lScalarPlain}{tabulate==0.9.0}
\PYG{+w}{    }\PYG{p+pIndicator}{\PYGZhy{}}\PYG{+w}{ }\PYG{l+lScalar+lScalarPlain}{uc\PYGZhy{}micro\PYGZhy{}py==1.0.3}
\PYG{+w}{    }\PYG{p+pIndicator}{\PYGZhy{}}\PYG{+w}{ }\PYG{l+lScalar+lScalarPlain}{zipp==3.20.2}
\end{sphinxVerbatim}

\sphinxAtStartPar
To set up the environment, follow these steps:
\begin{enumerate}
\sphinxsetlistlabels{\arabic}{enumi}{enumii}{}{.}%
\item {} 
\sphinxAtStartPar
Create the \sphinxcode{\sphinxupquote{environment.yml}} file from the content above.

\item {} 
\sphinxAtStartPar
Run the following command in your terminal:

\begin{sphinxVerbatim}[commandchars=\\\{\}]
conda\PYG{+w}{ }env\PYG{+w}{ }create\PYG{+w}{ }\PYGZhy{}f\PYG{+w}{ }environment.yml
\end{sphinxVerbatim}

\end{enumerate}

\sphinxstepscope


\chapter{Syllabus}
\label{\detokenize{syllabus:syllabus}}\label{\detokenize{syllabus::doc}}
\sphinxAtStartPar
\sphinxstylestrong{Course Duration}: 16 Weeks


\section{Course Description}
\label{\detokenize{syllabus:course-description}}
\sphinxAtStartPar
This course introduces the fundamental principles of fluid mechanics, focusing on the behavior of incompressible fluids. Topics include fluid statics, kinematics, and dynamics, with an emphasis on differential analysis, including the Navier\sphinxhyphen{}Stokes equations and Bernoulli’s equation. Students will apply these principles to solve practical engineering problems.


\section{Week\sphinxhyphen{}by\sphinxhyphen{}Week Outline}
\label{\detokenize{syllabus:week-by-week-outline}}

\subsection{\sphinxstylestrong{Weeks 1\sphinxhyphen{}2: Introduction to Fluid Mechanics and Properties}}
\label{\detokenize{syllabus:weeks-1-2-introduction-to-fluid-mechanics-and-properties}}\begin{itemize}
\item {} 
\sphinxAtStartPar
Prerequisites: Math and physics

\item {} 
\sphinxAtStartPar
Definition of fluids and distinction between solids and fluids.

\item {} 
\sphinxAtStartPar
Fluid properties: density, viscosity, surface tension.

\item {} 
\sphinxAtStartPar
Continuum hypothesis: fluid as a continuum.

\item {} 
\sphinxAtStartPar
Basic concepts: fluid pressure, temperature, and velocity.

\end{itemize}

\sphinxAtStartPar
\sphinxstylestrong{Lab}: Measurement of basic fluid properties (density, viscosity, surface tension).


\bigskip\hrule\bigskip



\subsection{\sphinxstylestrong{Weeks 3\sphinxhyphen{}4: Differential Analysis and Navier\sphinxhyphen{}Stokes Equations}}
\label{\detokenize{syllabus:weeks-3-4-differential-analysis-and-navier-stokes-equations}}\begin{itemize}
\item {} 
\sphinxAtStartPar
Review of Newton’s second law and its application to fluids.

\item {} 
\sphinxAtStartPar
Differential control volumes and the derivation of the \sphinxstylestrong{Navier\sphinxhyphen{}Stokes equations}.

\item {} 
\sphinxAtStartPar
Forces in fluid flow: pressure, body forces, and viscous forces.

\item {} 
\sphinxAtStartPar
Introduction to boundary conditions.

\item {} 
\item {} 
\item {} 
\item {} 
\end{itemize}


\bigskip\hrule\bigskip



\subsection{\sphinxstylestrong{Weeks 7\sphinxhyphen{}8: Fluid Statics as an Application of Navier\sphinxhyphen{}Stokes}}
\label{\detokenize{syllabus:weeks-7-8-fluid-statics-as-an-application-of-navier-stokes}}\begin{itemize}
\item {} 
\sphinxAtStartPar
\sphinxstylestrong{Fluid statics} as a special case of the Navier\sphinxhyphen{}Stokes equations (no velocity, steady\sphinxhyphen{}state).

\item {} 
\sphinxAtStartPar
Hydrostatic equation: (\textbackslash{}frac\{dp\}\{dz\} = \sphinxhyphen{}\textbackslash{}rho g).

\item {} 
\sphinxAtStartPar
Pressure variation in incompressible and compressible static fluids.

\item {} 
\sphinxAtStartPar
Applications: pressure forces on submerged surfaces, buoyancy, and stability.

\end{itemize}

\sphinxAtStartPar
\sphinxstylestrong{Lab}: Verification of hydrostatic pressure distribution using manometers and hydrostatic force on submerged surfaces.

\sphinxAtStartPar
\sphinxstylestrong{Lab}: Flow visualization and numerical simulation of simple flow fields using NS equations.


\bigskip\hrule\bigskip



\subsection{\sphinxstylestrong{Weeks 5\sphinxhyphen{}6: Fluid Kinematics (Description of Fluid Motion)}}
\label{\detokenize{syllabus:weeks-5-6-fluid-kinematics-description-of-fluid-motion}}\begin{itemize}
\item {} 
\sphinxAtStartPar
\sphinxstylestrong{Introduction to Fluid Kinematics}:
\begin{itemize}
\item {} 
\sphinxAtStartPar
Lagrangian vs. Eulerian descriptions of fluid motion.

\item {} 
\sphinxAtStartPar
Velocity and acceleration fields.

\item {} 
\sphinxAtStartPar
Flow classification: steady vs. unsteady, uniform vs. non\sphinxhyphen{}uniform, rotational vs. irrotational.

\item {} 
\sphinxAtStartPar
Streamlines, streaklines, and pathlines.

\end{itemize}

\item {} 
\sphinxAtStartPar
\sphinxstylestrong{Mathematical Description of Flow}:
\begin{itemize}
\item {} 
\sphinxAtStartPar
Velocity potential and stream function.

\item {} 
\sphinxAtStartPar
Concept of circulation and vorticity.

\item {} 
\sphinxAtStartPar
Application to incompressible and inviscid flows.

\end{itemize}

\end{itemize}

\sphinxAtStartPar
\sphinxstylestrong{Lab}: Visualization of streamlines and velocity fields using flow visualization techniques (dye injection, PIV).


\bigskip\hrule\bigskip



\bigskip\hrule\bigskip



\subsection{\sphinxstylestrong{Weeks 9\sphinxhyphen{}10: Bernoulli’s Equation and Energy Conservation}}
\label{\detokenize{syllabus:weeks-9-10-bernoullis-equation-and-energy-conservation}}\begin{itemize}
\item {} 
\sphinxAtStartPar
Derivation of \sphinxstylestrong{Bernoulli’s equation} from the Euler equations for steady, incompressible, inviscid flow.

\item {} 
\sphinxAtStartPar
Physical interpretation of Bernoulli’s equation as energy conservation.

\item {} 
\sphinxAtStartPar
Practical applications of Bernoulli’s equation (nozzles, flow meters, venturi tubes).

\end{itemize}

\sphinxAtStartPar
\sphinxstylestrong{Lab}: Flow measurement using Bernoulli’s principle in pipe flow and nozzle experiments.


\bigskip\hrule\bigskip



\subsection{\sphinxstylestrong{Weeks 11\sphinxhyphen{}12: Control Volume Analysis and Integral Conservation Laws}}
\label{\detokenize{syllabus:weeks-11-12-control-volume-analysis-and-integral-conservation-laws}}\begin{itemize}
\item {} 
\sphinxAtStartPar
Control volume approach for mass, momentum, and energy conservation.

\item {} 
\sphinxAtStartPar
Applications to real\sphinxhyphen{}world systems: jet propulsion, pipe bends, turbines, and pumps.

\item {} 
\sphinxAtStartPar
Introduction to the \sphinxstylestrong{continuity equation} in integral form.

\end{itemize}

\sphinxAtStartPar
\sphinxstylestrong{Lab}: Application of the control volume approach to analyze jet propulsion and momentum conservation in pipe bends.


\subsection{\sphinxstylestrong{Weeks 13\sphinxhyphen{}14: Viscous Flow and Internal Flow in Pipes}}
\label{\detokenize{syllabus:weeks-13-14-viscous-flow-and-internal-flow-in-pipes}}\begin{itemize}
\item {} 
\sphinxAtStartPar
Laminar vs. turbulent flow: \sphinxstylestrong{Hagen\sphinxhyphen{}Poiseuille equation} for laminar flow.

\item {} 
\sphinxAtStartPar
Head loss in pipes: \sphinxstylestrong{Moody chart} and friction factor.

\item {} 
\sphinxAtStartPar
Minor losses and pipe networks.

\end{itemize}

\sphinxAtStartPar
\sphinxstylestrong{Lab}: Measurement of head loss in pipe flow and analysis of friction losses.


\subsection{\sphinxstylestrong{Weeks 15\sphinxhyphen{}16: Dimensional Analysis, External Flow, and Turbomachinery}}
\label{\detokenize{syllabus:weeks-15-16-dimensional-analysis-external-flow-and-turbomachinery}}\begin{itemize}
\item {} 
\sphinxAtStartPar
\sphinxstylestrong{Dimensional Analysis}:
\begin{itemize}
\item {} 
\sphinxAtStartPar
\sphinxstylestrong{Buckingham Pi theorem} and dimensional homogeneity.

\item {} 
\sphinxAtStartPar
Important dimensionless numbers (Reynolds number, Froude number, Mach number).

\item {} 
\sphinxAtStartPar
Model testing and similarity laws.

\end{itemize}

\item {} 
\sphinxAtStartPar
\sphinxstylestrong{Boundary Layers and External Flow}:
\begin{itemize}
\item {} 
\sphinxAtStartPar
Introduction to \sphinxstylestrong{boundary layer theory}: laminar and turbulent boundary layers.

\item {} 
\sphinxAtStartPar
Flow separation and wake formation.

\item {} 
\sphinxAtStartPar
Drag and lift on immersed bodies: flow around cylinders, airfoils, and spheres.

\end{itemize}

\item {} 
\sphinxAtStartPar
\sphinxstylestrong{Turbomachinery}:
\begin{itemize}
\item {} 
\sphinxAtStartPar
Basics of compressible flow: speed of sound and Mach number.

\item {} 
\sphinxAtStartPar
Compressible flow regimes: subsonic, transonic, supersonic.

\item {} 
\sphinxAtStartPar
Performance of pumps, turbines, and compressors.

\item {} 
\sphinxAtStartPar
Energy transfer in turbomachinery, cavitation, and performance curves.

\end{itemize}

\end{itemize}

\sphinxAtStartPar
\sphinxstylestrong{Lab}: Wind tunnel experiments measuring drag and lift on simple objects (cylinders, airfoils). Performance testing of pumps and analysis of cavitation effects.


\section{Assessments}
\label{\detokenize{syllabus:assessments}}\begin{itemize}
\item {} 
\sphinxAtStartPar
Laboratory Reports (in\sphinxhyphen{}lab experiments and computational): 20\%

\item {} 
\sphinxAtStartPar
Midterm Exams (4): 80\%

\end{itemize}


\section{Material}
\label{\detokenize{syllabus:material}}\begin{itemize}
\item {} 
\sphinxAtStartPar
\sphinxstylestrong{Cengel, Yunus A., and John M. Cimbala.} Fluid Mechanics: Fundamentals and Applications. McGraw\sphinxhyphen{}Hill, 2018. {[}\hyperlink{cite._sources/markdown:id3}{CC18}{]}

\item {} 
\sphinxAtStartPar
\sphinxstylestrong{White, Frank M.} Fluid Mechanics. McGraw\sphinxhyphen{}Hill, 2016. {[}\hyperlink{cite._sources/markdown:id4}{Whi16}{]}

\item {} 
\sphinxAtStartPar
\sphinxstylestrong{Munson, Bruce R., Donald F. Young, Theodore H. Okiishi, and Wade W. Huebsch.} Fundamentals of Fluid Mechanics. Wiley, 2012. {[}\hyperlink{cite._sources/markdown:id5}{MYOH12}{]}

\item {} 
\sphinxAtStartPar
\sphinxstylestrong{Shames, Irving H.} Mechanics of Fluids. McGraw\sphinxhyphen{}Hill, 2003. {[}\hyperlink{cite._sources/markdown:id6}{Sha03}{]}

\item {} 
\sphinxAtStartPar
\sphinxstylestrong{Fox, Robert W., Alan T. McDonald, and Philip J. Pritchard.} Introduction to Fluid Mechanics. Wiley, 2011. {[}\hyperlink{cite._sources/markdown:id7}{FMP11}{]}

\end{itemize}

\sphinxstepscope


\chapter{Prerequisites}
\label{\detokenize{prerequisites:prerequisites}}\label{\detokenize{prerequisites::doc}}

\section{Mathematics Prerequisites}
\label{\detokenize{prerequisites:mathematics-prerequisites}}\begin{enumerate}
\sphinxsetlistlabels{\arabic}{enumi}{enumii}{}{.}%
\item {} 
\sphinxAtStartPar
\sphinxstylestrong{Calculus}:
\begin{itemize}
\item {} 
\sphinxAtStartPar
Understanding of derivatives and integrals, especially related to functions of one or more variables.

\item {} 
\sphinxAtStartPar
Concepts of partial derivatives, chain rule, and multiple integrals.

\item {} 
\sphinxAtStartPar
Parametric equations and their applications.

\end{itemize}

\item {} 
\sphinxAtStartPar
\sphinxstylestrong{Differential Equations}:
\begin{itemize}
\item {} 
\sphinxAtStartPar
Ability to solve ordinary differential equations (ODEs), which are essential for understanding fluid dynamics, particularly in deriving the Navier\sphinxhyphen{}Stokes equations and other conservation laws.

\item {} 
\sphinxAtStartPar
Familiarity with linear and non\sphinxhyphen{}linear differential equations.

\item {} 
\sphinxAtStartPar
Exposure to boundary value problems and initial conditions.

\end{itemize}

\item {} 
\sphinxAtStartPar
\sphinxstylestrong{Vector Calculus}:
\begin{itemize}
\item {} 
\sphinxAtStartPar
Mastery of vector fields, gradient, divergence, and curl operations.

\item {} 
\sphinxAtStartPar
Line integrals, surface integrals, and volume integrals.

\item {} 
\sphinxAtStartPar
Understanding the physical meaning of the divergence theorem and Stokes’ theorem, which are crucial in fluid flow analysis using control volumes and differential analysis.

\end{itemize}

\item {} 
\sphinxAtStartPar
\sphinxstylestrong{Linear Algebra (optional)}:
\begin{itemize}
\item {} 
\sphinxAtStartPar
Some familiarity with matrices and solving systems of equations, particularly in computational aspects like numerical simulations or matrix\sphinxhyphen{}based methods for solving differential equations.

\end{itemize}

\end{enumerate}


\section{Physics Prerequisites}
\label{\detokenize{prerequisites:physics-prerequisites}}\begin{enumerate}
\sphinxsetlistlabels{\arabic}{enumi}{enumii}{}{.}%
\item {} 
\sphinxAtStartPar
\sphinxstylestrong{General Physics (Mechanics)}:
\begin{itemize}
\item {} 
\sphinxAtStartPar
Mastery of Newton’s laws of motion and their application to different systems.

\item {} 
\sphinxAtStartPar
Knowledge of energy conservation (kinetic and potential energy) and momentum conservation, as these concepts directly translate to Bernoulli’s equation and fluid statics.

\item {} 
\sphinxAtStartPar
Basics of forces, including gravitational forces and frictional forces.

\end{itemize}

\item {} 
\sphinxAtStartPar
\sphinxstylestrong{General Physics (Electricity and Magnetism, optional)}:
\begin{itemize}
\item {} 
\sphinxAtStartPar
Some exposure to wave propagation, which can help in understanding aspects of compressible flow and shock waves.

\item {} 
\sphinxAtStartPar
Familiarity with concepts like pressure fields, as similar ideas are used in both electromagnetism and fluid mechanics.

\end{itemize}

\end{enumerate}


\bigskip\hrule\bigskip



\section{Problem 1: Understanding Derivatives and Integrals}
\label{\detokenize{prerequisites:problem-1-understanding-derivatives-and-integrals}}
\sphinxAtStartPar
Given a function \(f(x)\), find the derivative and the integral of \(f(x)\) over a given interval.


\subsection{Problem Statement}
\label{\detokenize{prerequisites:problem-statement}}
\sphinxAtStartPar
Let \(f(x) = x^3 - 3x^2 + 2x\).
\begin{enumerate}
\sphinxsetlistlabels{\arabic}{enumi}{enumii}{}{.}%
\item {} 
\sphinxAtStartPar
Find \(f'(x)\), the derivative of the function.

\item {} 
\sphinxAtStartPar
Calculate the integral \(\int f(x) dx\) over the interval {[}1, 4{]}.

\end{enumerate}


\subsection{Solution}
\label{\detokenize{prerequisites:solution}}

\subsubsection{1. \sphinxstylestrong{Derivative}}
\label{\detokenize{prerequisites:derivative}}
\sphinxAtStartPar
Using basic differentiation rules:
\begin{equation*}
\begin{split}
f'(x) = 3x^2 - 6x + 2
\end{split}
\end{equation*}

\subsubsection{2. \sphinxstylestrong{Integral}}
\label{\detokenize{prerequisites:integral}}
\sphinxAtStartPar
The indefinite integral is:
\begin{equation*}
\begin{split}
\int (x^3 - 3x^2 + 2x) dx = \frac{x^4}{4} - x^3 + x^2 + C
\end{split}
\end{equation*}
\sphinxAtStartPar
To evaluate the definite integral over the interval {[}1, 4{]}:
\begin{equation*}
\begin{split}
\left[\frac{x^4}{4} - x^3 + x^2 \right]_{1}^{4}
\end{split}
\end{equation*}
\sphinxAtStartPar
Substituting the limits:
\begin{equation*}
\begin{split}
\left[\frac{4^4}{4} - 4^3 + 4^2\right] - \left[\frac{1^4}{4} - 1^3 + 1^2\right]
\end{split}
\end{equation*}
\sphinxAtStartPar
Evaluating both terms:
\begin{equation*}
\begin{split}
\left[\frac{256}{4} - 64 + 16\right] - \left[\frac{1}{4} - 1 + 1\right] = (64 - 64 + 16) - \left(\frac{1}{4}\right) = 16 - \frac{1}{4}
\end{split}
\end{equation*}
\sphinxAtStartPar
Thus, the definite integral evaluates to:
\begin{equation*}
\begin{split}
15.75
\end{split}
\end{equation*}

\subsection{Python Visualization and Area Calculation}
\label{\detokenize{prerequisites:python-visualization-and-area-calculation}}
\sphinxAtStartPar
Here is the Python code to visualize the function and calculate the area under the curve:

\begin{sphinxuseclass}{cell}\begin{sphinxVerbatimInput}

\begin{sphinxuseclass}{cell_input}
\begin{sphinxVerbatim}[commandchars=\\\{\}]
\PYG{k+kn}{import} \PYG{n+nn}{numpy} \PYG{k}{as} \PYG{n+nn}{np}
\PYG{k+kn}{import} \PYG{n+nn}{matplotlib}\PYG{n+nn}{.}\PYG{n+nn}{pyplot} \PYG{k}{as} \PYG{n+nn}{plt}
\PYG{k+kn}{from} \PYG{n+nn}{scipy}\PYG{n+nn}{.}\PYG{n+nn}{integrate} \PYG{k+kn}{import} \PYG{n}{quad}

\PYG{c+c1}{\PYGZsh{} Define the function f(x) = x\PYGZca{}3 \PYGZhy{} 3x\PYGZca{}2 + 2x}
\PYG{k}{def} \PYG{n+nf}{f}\PYG{p}{(}\PYG{n}{x}\PYG{p}{)}\PYG{p}{:}
    \PYG{k}{return} \PYG{n}{x}\PYG{o}{*}\PYG{o}{*}\PYG{l+m+mi}{3} \PYG{o}{\PYGZhy{}} \PYG{l+m+mi}{3}\PYG{o}{*}\PYG{n}{x}\PYG{o}{*}\PYG{o}{*}\PYG{l+m+mi}{2} \PYG{o}{+} \PYG{l+m+mi}{2}\PYG{o}{*}\PYG{n}{x}

\PYG{c+c1}{\PYGZsh{} Generate x values for plotting}
\PYG{n}{x\PYGZus{}vals} \PYG{o}{=} \PYG{n}{np}\PYG{o}{.}\PYG{n}{linspace}\PYG{p}{(}\PYG{l+m+mi}{0}\PYG{p}{,} \PYG{l+m+mi}{5}\PYG{p}{,} \PYG{l+m+mi}{100}\PYG{p}{)}
\PYG{n}{y\PYGZus{}vals} \PYG{o}{=} \PYG{n}{f}\PYG{p}{(}\PYG{n}{x\PYGZus{}vals}\PYG{p}{)}

\PYG{c+c1}{\PYGZsh{} Plot the function}
\PYG{n}{plt}\PYG{o}{.}\PYG{n}{figure}\PYG{p}{(}\PYG{n}{figsize}\PYG{o}{=}\PYG{p}{(}\PYG{l+m+mi}{8}\PYG{p}{,} \PYG{l+m+mi}{6}\PYG{p}{)}\PYG{p}{)}
\PYG{n}{plt}\PYG{o}{.}\PYG{n}{plot}\PYG{p}{(}\PYG{n}{x\PYGZus{}vals}\PYG{p}{,} \PYG{n}{y\PYGZus{}vals}\PYG{p}{,} \PYG{n}{label}\PYG{o}{=}\PYG{l+s+sa}{r}\PYG{l+s+s1}{\PYGZsq{}}\PYG{l+s+s1}{\PYGZdl{}f(x) = x\PYGZca{}3 \PYGZhy{} 3x\PYGZca{}2 + 2x\PYGZdl{}}\PYG{l+s+s1}{\PYGZsq{}}\PYG{p}{,} \PYG{n}{color}\PYG{o}{=}\PYG{l+s+s2}{\PYGZdq{}}\PYG{l+s+s2}{blue}\PYG{l+s+s2}{\PYGZdq{}}\PYG{p}{)}
\PYG{n}{plt}\PYG{o}{.}\PYG{n}{fill\PYGZus{}between}\PYG{p}{(}\PYG{n}{x\PYGZus{}vals}\PYG{p}{,} \PYG{n}{y\PYGZus{}vals}\PYG{p}{,} \PYG{n}{where}\PYG{o}{=}\PYG{p}{[}\PYG{p}{(}\PYG{l+m+mi}{1} \PYG{o}{\PYGZlt{}}\PYG{o}{=} \PYG{n}{x} \PYG{o}{\PYGZlt{}}\PYG{o}{=} \PYG{l+m+mi}{4}\PYG{p}{)} \PYG{k}{for} \PYG{n}{x} \PYG{o+ow}{in} \PYG{n}{x\PYGZus{}vals}\PYG{p}{]}\PYG{p}{,} \PYG{n}{color}\PYG{o}{=}\PYG{l+s+s1}{\PYGZsq{}}\PYG{l+s+s1}{lightblue}\PYG{l+s+s1}{\PYGZsq{}}\PYG{p}{,} \PYG{n}{alpha}\PYG{o}{=}\PYG{l+m+mf}{0.5}\PYG{p}{)}

\PYG{c+c1}{\PYGZsh{} Labels and Title}
\PYG{n}{plt}\PYG{o}{.}\PYG{n}{axhline}\PYG{p}{(}\PYG{l+m+mi}{0}\PYG{p}{,} \PYG{n}{color}\PYG{o}{=}\PYG{l+s+s1}{\PYGZsq{}}\PYG{l+s+s1}{black}\PYG{l+s+s1}{\PYGZsq{}}\PYG{p}{,}\PYG{n}{linewidth}\PYG{o}{=}\PYG{l+m+mf}{0.5}\PYG{p}{)}
\PYG{n}{plt}\PYG{o}{.}\PYG{n}{axvline}\PYG{p}{(}\PYG{l+m+mi}{0}\PYG{p}{,} \PYG{n}{color}\PYG{o}{=}\PYG{l+s+s1}{\PYGZsq{}}\PYG{l+s+s1}{black}\PYG{l+s+s1}{\PYGZsq{}}\PYG{p}{,}\PYG{n}{linewidth}\PYG{o}{=}\PYG{l+m+mf}{0.5}\PYG{p}{)}
\PYG{n}{plt}\PYG{o}{.}\PYG{n}{xlabel}\PYG{p}{(}\PYG{l+s+s1}{\PYGZsq{}}\PYG{l+s+s1}{x}\PYG{l+s+s1}{\PYGZsq{}}\PYG{p}{)}
\PYG{n}{plt}\PYG{o}{.}\PYG{n}{ylabel}\PYG{p}{(}\PYG{l+s+s1}{\PYGZsq{}}\PYG{l+s+s1}{f(x)}\PYG{l+s+s1}{\PYGZsq{}}\PYG{p}{)}
\PYG{n}{plt}\PYG{o}{.}\PYG{n}{title}\PYG{p}{(}\PYG{l+s+s1}{\PYGZsq{}}\PYG{l+s+s1}{Function \PYGZdl{}f(x) = x\PYGZca{}3 \PYGZhy{} 3x\PYGZca{}2 + 2x\PYGZdl{} and the Area Under the Curve [1, 4]}\PYG{l+s+s1}{\PYGZsq{}}\PYG{p}{)}
\PYG{n}{plt}\PYG{o}{.}\PYG{n}{legend}\PYG{p}{(}\PYG{p}{)}
\PYG{n}{plt}\PYG{o}{.}\PYG{n}{grid}\PYG{p}{(}\PYG{k+kc}{True}\PYG{p}{)}
\PYG{n}{plt}\PYG{o}{.}\PYG{n}{show}\PYG{p}{(}\PYG{p}{)}

\PYG{c+c1}{\PYGZsh{} Calculate the integral over the interval [1, 4]}
\PYG{n}{area}\PYG{p}{,} \PYG{n}{\PYGZus{}} \PYG{o}{=} \PYG{n}{quad}\PYG{p}{(}\PYG{n}{f}\PYG{p}{,} \PYG{l+m+mi}{1}\PYG{p}{,} \PYG{l+m+mi}{4}\PYG{p}{)}

\PYG{n+nb}{print}\PYG{p}{(}\PYG{l+s+sa}{f}\PYG{l+s+s2}{\PYGZdq{}}\PYG{l+s+s2}{The area under the curve from x=1 to x=4 is }\PYG{l+s+si}{\PYGZob{}}\PYG{n}{area}\PYG{l+s+si}{:}\PYG{l+s+s2}{.2f}\PYG{l+s+si}{\PYGZcb{}}\PYG{l+s+s2}{.}\PYG{l+s+s2}{\PYGZdq{}}\PYG{p}{)}
\end{sphinxVerbatim}

\end{sphinxuseclass}\end{sphinxVerbatimInput}
\begin{sphinxVerbatimOutput}

\begin{sphinxuseclass}{cell_output}
\noindent\sphinxincludegraphics{{d01ac6626333f287f9a83db74bd9509f926563aa0832f9d54a717b8698685099}.png}

\begin{sphinxVerbatim}[commandchars=\\\{\}]
The area under the curve from x=1 to x=4 is 15.75.
\end{sphinxVerbatim}

\end{sphinxuseclass}\end{sphinxVerbatimOutput}

\end{sphinxuseclass}

\bigskip\hrule\bigskip



\section{Problem 2: Partial Derivatives, Chain Rule, and Multiple Integrals}
\label{\detokenize{prerequisites:problem-2-partial-derivatives-chain-rule-and-multiple-integrals}}

\subsection{Problem Statement}
\label{\detokenize{prerequisites:id1}}
\sphinxAtStartPar
Let \(z = f(x, y)\) where \(x = u^2 + 3v\) and \(y = uv\).
\begin{enumerate}
\sphinxsetlistlabels{\arabic}{enumi}{enumii}{}{.}%
\item {} 
\sphinxAtStartPar
Find the partial derivative of \(z\) with respect to \(u\), denoted as \(\frac{\partial z}{\partial u}\), using the chain rule.

\item {} 
\sphinxAtStartPar
Calculate the double integral of \(z(x, y) = x + y\) over the region \(R\), where \(R\) is the square defined by \(0 \leq x \leq 2\) and \(0 \leq y \leq 1\).

\end{enumerate}


\subsection{Solution:}
\label{\detokenize{prerequisites:id2}}

\subsubsection{1. \sphinxstylestrong{Applying the Chain Rule}:}
\label{\detokenize{prerequisites:applying-the-chain-rule}}
\sphinxAtStartPar
Let \(f(x, y) = x + y\).

\sphinxAtStartPar
To find \(\frac{\partial z}{\partial u}\), use the chain rule:
\begin{equation*}
\begin{split}
\frac{\partial z}{\partial u} = \frac{\partial f}{\partial x} \cdot \frac{\partial x}{\partial u} + \frac{\partial f}{\partial y} \cdot \frac{\partial y}{\partial u}
\end{split}
\end{equation*}
\sphinxAtStartPar
For the function \(f(x, y) = x + y\):
\begin{equation*}
\begin{split}
\frac{\partial f}{\partial x} = 1, \quad \frac{\partial f}{\partial y} = 1
\end{split}
\end{equation*}
\sphinxAtStartPar
Next, calculate \(\frac{\partial x}{\partial u}\) and \(\frac{\partial y}{\partial u}\):

\sphinxAtStartPar
Given \(x = u^2 + 3v\) and \(y = uv\):
\begin{equation*}
\begin{split}
\frac{\partial x}{\partial u} = 2u, \quad \frac{\partial y}{\partial u} = v
\end{split}
\end{equation*}
\sphinxAtStartPar
Now substitute these into the chain rule equation:
\begin{equation*}
\begin{split}
\frac{\partial z}{\partial u} = 1 \cdot 2u + 1 \cdot v = 2u + v
\end{split}
\end{equation*}
\sphinxAtStartPar
Thus, the partial derivative of \(z\) with respect to \(u\) is:
\begin{equation*}
\begin{split}
\frac{\partial z}{\partial u} = 2u + v
\end{split}
\end{equation*}

\subsubsection{2. \sphinxstylestrong{Multiple Integrals}:}
\label{\detokenize{prerequisites:multiple-integrals}}
\sphinxAtStartPar
Now, let’s compute the double integral of \(z(x, y) = x + y\) over the region \(R\), defined by \(0 \leq x \leq 2\) and \(0 \leq y \leq 1\).

\sphinxAtStartPar
The integral to solve is:
\begin{equation*}
\begin{split}
\int_0^1 \int_0^2 (x + y) \, dx \, dy
\end{split}
\end{equation*}
\sphinxAtStartPar
First, integrate with respect to \(x\):
\begin{equation*}
\begin{split}
\int_0^2 (x + y) \, dx = \int_0^2 x \, dx + \int_0^2 y \, dx = \left[ \frac{x^2}{2} \right]_0^2 + \left[ yx \right]_0^2
\end{split}
\end{equation*}
\sphinxAtStartPar
Evaluating the integrals:
\begin{equation*}
\begin{split}
\frac{2^2}{2} - 0 + y(2 - 0) = 2 + 2y
\end{split}
\end{equation*}
\sphinxAtStartPar
Now, integrate with respect to \(y\):
\begin{equation*}
\begin{split}
\int_0^1 (2 + 2y) \, dy = \int_0^1 2 \, dy + \int_0^1 2y \, dy = \left[ 2y \right]_0^1 + \left[ y^2 \right]_0^1
\end{split}
\end{equation*}
\sphinxAtStartPar
Evaluating the integrals:
\begin{equation*}
\begin{split}
2(1 - 0) + (1^2 - 0^2) = 2 + 1 = 3
\end{split}
\end{equation*}
\sphinxAtStartPar
Thus, the value of the double integral is:
\begin{equation*}
\begin{split}
\int_0^1 \int_0^2 (x + y) \, dx \, dy = 3
\end{split}
\end{equation*}

\bigskip\hrule\bigskip



\section{Problem 3: Parametric Equations and Applications}
\label{\detokenize{prerequisites:problem-3-parametric-equations-and-applications}}

\subsection{Problem Statement}
\label{\detokenize{prerequisites:id3}}
\sphinxAtStartPar
A particle is moving along a curve defined by the following parametric equations:
\begin{equation*}
\begin{split}
x(t) = 3t^2 - 2t
\end{split}
\end{equation*}\begin{equation*}
\begin{split}
y(t) = 2t^3 - t
\end{split}
\end{equation*}
\sphinxAtStartPar
where \(x(t)\) and \(y(t)\) represent the particle’s position at time \(t\).
\begin{enumerate}
\sphinxsetlistlabels{\arabic}{enumi}{enumii}{}{.}%
\item {} 
\sphinxAtStartPar
Find the velocity and acceleration vectors of the particle at time \(t\).

\item {} 
\sphinxAtStartPar
Determine the time \(t\) when the particle’s velocity is zero.

\item {} 
\sphinxAtStartPar
Find the length of the curve traced by the particle between \(t = 0\) and \(t = 2\).

\end{enumerate}


\subsection{Solution}
\label{\detokenize{prerequisites:id4}}

\subsubsection{1. Velocity and Acceleration Vectors}
\label{\detokenize{prerequisites:velocity-and-acceleration-vectors}}
\sphinxAtStartPar
To find the velocity and acceleration, we differentiate the parametric equations with respect to time \(t\).
\begin{itemize}
\item {} 
\sphinxAtStartPar
\sphinxstylestrong{Velocity} is the derivative of position with respect to time:
\begin{equation*}
\begin{split}
  v_x(t) = \frac{dx}{dt} = 6t - 2
  \end{split}
\end{equation*}\begin{equation*}
\begin{split}
  v_y(t) = \frac{dy}{dt} = 6t^2 - 1
  \end{split}
\end{equation*}
\sphinxAtStartPar
Therefore, the velocity vector at time \(t\) is:
\begin{equation*}
\begin{split}
  \mathbf{v}(t) = (6t - 2) \hat{i} + (6t^2 - 1) \hat{j}
  \end{split}
\end{equation*}
\item {} 
\sphinxAtStartPar
\sphinxstylestrong{Acceleration} is the derivative of velocity with respect to time:
\begin{equation*}
\begin{split}
  a_x(t) = \frac{d^2x}{dt^2} = 6
  \end{split}
\end{equation*}\begin{equation*}
\begin{split}
  a_y(t) = \frac{d^2y}{dt^2} = 12t
  \end{split}
\end{equation*}
\sphinxAtStartPar
Therefore, the acceleration vector is:
\begin{equation*}
\begin{split}
  \mathbf{a}(t) = 6 \hat{i} + 12t \hat{j}
  \end{split}
\end{equation*}
\end{itemize}


\subsubsection{2. Time When the Particle’s Velocity is Zero}
\label{\detokenize{prerequisites:time-when-the-particle-s-velocity-is-zero}}
\sphinxAtStartPar
For the velocity to be zero, both components of the velocity vector must be zero simultaneously:
\begin{equation*}
\begin{split}
  6t - 2 = 0 \quad \text{and} \quad 6t^2 - 1 = 0
  \end{split}
\end{equation*}
\sphinxAtStartPar
From \(6t - 2 = 0\), solving for \(t\) gives:
\begin{equation*}
\begin{split}
  t = \frac{1}{3}
  \end{split}
\end{equation*}
\sphinxAtStartPar
Substituting \(t = \frac{1}{3}\) into the second equation:
\begin{equation*}
\begin{split}
  6 \left(\frac{1}{3}\right)^2 - 1 = \frac{2}{3} - 1 = -\frac{1}{3} \neq 0
  \end{split}
\end{equation*}
\sphinxAtStartPar
Thus, there is no time \(t\) where the particle’s velocity is exactly zero in both the \(x\) and \(y\) directions.


\subsubsection{3. Length of the Curve}
\label{\detokenize{prerequisites:length-of-the-curve}}
\sphinxAtStartPar
The length of the curve traced by the particle from \(t = 0\) to \(t = 2\) is given by the formula:
\begin{equation*}
\begin{split}
  L = \int_0^2 \sqrt{\left(\frac{dx}{dt}\right)^2 + \left(\frac{dy}{dt}\right)^2} \, dt
  \end{split}
\end{equation*}
\sphinxAtStartPar
Substitute the derivatives \(v_x(t) = 6t - 2\) and \(v_y(t) = 6t^2 - 1\):
\begin{equation*}
\begin{split}
  L = \int_0^2 \sqrt{(6t - 2)^2 + (6t^2 - 1)^2} \, dt
  \end{split}
\end{equation*}
\sphinxAtStartPar
This integral can be solved either numerically or using appropriate substitution techniques to find the length of the curve.


\subsection{Applications}
\label{\detokenize{prerequisites:applications}}
\sphinxAtStartPar
Parametric equations are widely used to describe curves and motions in physics, engineering, and computer graphics. For instance, they are used to model the trajectory of particles, motion along a path, and even design complex curves in computer\sphinxhyphen{}aided design (CAD) software.


\subsection{Python Visualization and Curve Length Calculation:}
\label{\detokenize{prerequisites:python-visualization-and-curve-length-calculation}}
\sphinxAtStartPar
Here’s the Python code to plot the parametric curve, velocity and acceleration vectors, and calculate the length of the curve numerically:

\begin{sphinxuseclass}{cell}\begin{sphinxVerbatimInput}

\begin{sphinxuseclass}{cell_input}
\begin{sphinxVerbatim}[commandchars=\\\{\}]
\PYG{k+kn}{import} \PYG{n+nn}{numpy} \PYG{k}{as} \PYG{n+nn}{np}
\PYG{k+kn}{import} \PYG{n+nn}{matplotlib}\PYG{n+nn}{.}\PYG{n+nn}{pyplot} \PYG{k}{as} \PYG{n+nn}{plt}
\PYG{k+kn}{from} \PYG{n+nn}{scipy}\PYG{n+nn}{.}\PYG{n+nn}{integrate} \PYG{k+kn}{import} \PYG{n}{quad}

\PYG{c+c1}{\PYGZsh{} Define the parametric equations for x(t) and y(t)}
\PYG{k}{def} \PYG{n+nf}{x}\PYG{p}{(}\PYG{n}{t}\PYG{p}{)}\PYG{p}{:}
    \PYG{k}{return} \PYG{l+m+mi}{3} \PYG{o}{*} \PYG{n}{t}\PYG{o}{*}\PYG{o}{*}\PYG{l+m+mi}{2} \PYG{o}{\PYGZhy{}} \PYG{l+m+mi}{2} \PYG{o}{*} \PYG{n}{t}

\PYG{k}{def} \PYG{n+nf}{y}\PYG{p}{(}\PYG{n}{t}\PYG{p}{)}\PYG{p}{:}
    \PYG{k}{return} \PYG{l+m+mi}{2} \PYG{o}{*} \PYG{n}{t}\PYG{o}{*}\PYG{o}{*}\PYG{l+m+mi}{3} \PYG{o}{\PYGZhy{}} \PYG{n}{t}

\PYG{c+c1}{\PYGZsh{} Define the velocity components}
\PYG{k}{def} \PYG{n+nf}{v\PYGZus{}x}\PYG{p}{(}\PYG{n}{t}\PYG{p}{)}\PYG{p}{:}
    \PYG{k}{return} \PYG{l+m+mi}{6} \PYG{o}{*} \PYG{n}{t} \PYG{o}{\PYGZhy{}} \PYG{l+m+mi}{2}

\PYG{k}{def} \PYG{n+nf}{v\PYGZus{}y}\PYG{p}{(}\PYG{n}{t}\PYG{p}{)}\PYG{p}{:}
    \PYG{k}{return} \PYG{l+m+mi}{6} \PYG{o}{*} \PYG{n}{t}\PYG{o}{*}\PYG{o}{*}\PYG{l+m+mi}{2} \PYG{o}{\PYGZhy{}} \PYG{l+m+mi}{1}

\PYG{c+c1}{\PYGZsh{} Define the acceleration components}
\PYG{k}{def} \PYG{n+nf}{a\PYGZus{}x}\PYG{p}{(}\PYG{n}{t}\PYG{p}{)}\PYG{p}{:}
    \PYG{k}{return} \PYG{l+m+mi}{6}

\PYG{k}{def} \PYG{n+nf}{a\PYGZus{}y}\PYG{p}{(}\PYG{n}{t}\PYG{p}{)}\PYG{p}{:}
    \PYG{k}{return} \PYG{l+m+mi}{12} \PYG{o}{*} \PYG{n}{t}

\PYG{c+c1}{\PYGZsh{} Define the integrand for the length of the curve}
\PYG{k}{def} \PYG{n+nf}{integrand}\PYG{p}{(}\PYG{n}{t}\PYG{p}{)}\PYG{p}{:}
    \PYG{k}{return} \PYG{n}{np}\PYG{o}{.}\PYG{n}{sqrt}\PYG{p}{(}\PYG{n}{v\PYGZus{}x}\PYG{p}{(}\PYG{n}{t}\PYG{p}{)}\PYG{o}{*}\PYG{o}{*}\PYG{l+m+mi}{2} \PYG{o}{+} \PYG{n}{v\PYGZus{}y}\PYG{p}{(}\PYG{n}{t}\PYG{p}{)}\PYG{o}{*}\PYG{o}{*}\PYG{l+m+mi}{2}\PYG{p}{)}

\PYG{c+c1}{\PYGZsh{} Time range for plotting}
\PYG{n}{t\PYGZus{}vals} \PYG{o}{=} \PYG{n}{np}\PYG{o}{.}\PYG{n}{linspace}\PYG{p}{(}\PYG{l+m+mi}{0}\PYG{p}{,} \PYG{l+m+mi}{2}\PYG{p}{,} \PYG{l+m+mi}{100}\PYG{p}{)}
\PYG{n}{x\PYGZus{}vals} \PYG{o}{=} \PYG{n}{x}\PYG{p}{(}\PYG{n}{t\PYGZus{}vals}\PYG{p}{)}
\PYG{n}{y\PYGZus{}vals} \PYG{o}{=} \PYG{n}{y}\PYG{p}{(}\PYG{n}{t\PYGZus{}vals}\PYG{p}{)}

\PYG{c+c1}{\PYGZsh{} Plot the parametric curve}
\PYG{n}{plt}\PYG{o}{.}\PYG{n}{figure}\PYG{p}{(}\PYG{n}{figsize}\PYG{o}{=}\PYG{p}{(}\PYG{l+m+mi}{8}\PYG{p}{,} \PYG{l+m+mi}{6}\PYG{p}{)}\PYG{p}{)}
\PYG{n}{plt}\PYG{o}{.}\PYG{n}{plot}\PYG{p}{(}\PYG{n}{x\PYGZus{}vals}\PYG{p}{,} \PYG{n}{y\PYGZus{}vals}\PYG{p}{,} \PYG{n}{label}\PYG{o}{=}\PYG{l+s+s1}{\PYGZsq{}}\PYG{l+s+s1}{Parametric Curve: Particle Path}\PYG{l+s+s1}{\PYGZsq{}}\PYG{p}{,} \PYG{n}{color}\PYG{o}{=}\PYG{l+s+s1}{\PYGZsq{}}\PYG{l+s+s1}{blue}\PYG{l+s+s1}{\PYGZsq{}}\PYG{p}{)}

\PYG{c+c1}{\PYGZsh{} Calculate and plot velocity and acceleration vectors at specific points}
\PYG{n}{t\PYGZus{}points} \PYG{o}{=} \PYG{n}{np}\PYG{o}{.}\PYG{n}{array}\PYG{p}{(}\PYG{p}{[}\PYG{l+m+mi}{0}\PYG{p}{,} \PYG{l+m+mf}{0.5}\PYG{p}{,} \PYG{l+m+mi}{1}\PYG{p}{,} \PYG{l+m+mf}{1.5}\PYG{p}{,} \PYG{l+m+mi}{2}\PYG{p}{]}\PYG{p}{)}
\PYG{n}{v\PYGZus{}x\PYGZus{}vals} \PYG{o}{=} \PYG{n}{v\PYGZus{}x}\PYG{p}{(}\PYG{n}{t\PYGZus{}points}\PYG{p}{)}
\PYG{n}{v\PYGZus{}y\PYGZus{}vals} \PYG{o}{=} \PYG{n}{v\PYGZus{}y}\PYG{p}{(}\PYG{n}{t\PYGZus{}points}\PYG{p}{)}
\PYG{n}{a\PYGZus{}x\PYGZus{}vals} \PYG{o}{=} \PYG{n}{np}\PYG{o}{.}\PYG{n}{full\PYGZus{}like}\PYG{p}{(}\PYG{n}{t\PYGZus{}points}\PYG{p}{,} \PYG{n}{a\PYGZus{}x}\PYG{p}{(}\PYG{l+m+mi}{0}\PYG{p}{)}\PYG{p}{)}  \PYG{c+c1}{\PYGZsh{} a\PYGZus{}x is constant, so create an array with the same length}
\PYG{n}{a\PYGZus{}y\PYGZus{}vals} \PYG{o}{=} \PYG{n}{a\PYGZus{}y}\PYG{p}{(}\PYG{n}{t\PYGZus{}points}\PYG{p}{)}

\PYG{c+c1}{\PYGZsh{} Plot velocity vectors}
\PYG{n}{plt}\PYG{o}{.}\PYG{n}{quiver}\PYG{p}{(}\PYG{n}{x}\PYG{p}{(}\PYG{n}{t\PYGZus{}points}\PYG{p}{)}\PYG{p}{,} \PYG{n}{y}\PYG{p}{(}\PYG{n}{t\PYGZus{}points}\PYG{p}{)}\PYG{p}{,} \PYG{n}{v\PYGZus{}x\PYGZus{}vals}\PYG{p}{,} \PYG{n}{v\PYGZus{}y\PYGZus{}vals}\PYG{p}{,} \PYG{n}{color}\PYG{o}{=}\PYG{l+s+s1}{\PYGZsq{}}\PYG{l+s+s1}{green}\PYG{l+s+s1}{\PYGZsq{}}\PYG{p}{,} \PYG{n}{angles}\PYG{o}{=}\PYG{l+s+s1}{\PYGZsq{}}\PYG{l+s+s1}{xy}\PYG{l+s+s1}{\PYGZsq{}}\PYG{p}{,} \PYG{n}{scale\PYGZus{}units}\PYG{o}{=}\PYG{l+s+s1}{\PYGZsq{}}\PYG{l+s+s1}{xy}\PYG{l+s+s1}{\PYGZsq{}}\PYG{p}{,} \PYG{n}{scale}\PYG{o}{=}\PYG{l+m+mi}{1}\PYG{p}{,} \PYG{n}{label}\PYG{o}{=}\PYG{l+s+s2}{\PYGZdq{}}\PYG{l+s+s2}{Velocity Vectors}\PYG{l+s+s2}{\PYGZdq{}}\PYG{p}{)}

\PYG{c+c1}{\PYGZsh{} Plot acceleration vectors}
\PYG{n}{plt}\PYG{o}{.}\PYG{n}{quiver}\PYG{p}{(}\PYG{n}{x}\PYG{p}{(}\PYG{n}{t\PYGZus{}points}\PYG{p}{)}\PYG{p}{,} \PYG{n}{y}\PYG{p}{(}\PYG{n}{t\PYGZus{}points}\PYG{p}{)}\PYG{p}{,} \PYG{n}{a\PYGZus{}x\PYGZus{}vals}\PYG{p}{,} \PYG{n}{a\PYGZus{}y\PYGZus{}vals}\PYG{p}{,} \PYG{n}{color}\PYG{o}{=}\PYG{l+s+s1}{\PYGZsq{}}\PYG{l+s+s1}{red}\PYG{l+s+s1}{\PYGZsq{}}\PYG{p}{,} \PYG{n}{angles}\PYG{o}{=}\PYG{l+s+s1}{\PYGZsq{}}\PYG{l+s+s1}{xy}\PYG{l+s+s1}{\PYGZsq{}}\PYG{p}{,} \PYG{n}{scale\PYGZus{}units}\PYG{o}{=}\PYG{l+s+s1}{\PYGZsq{}}\PYG{l+s+s1}{xy}\PYG{l+s+s1}{\PYGZsq{}}\PYG{p}{,} \PYG{n}{scale}\PYG{o}{=}\PYG{l+m+mi}{1}\PYG{p}{,} \PYG{n}{label}\PYG{o}{=}\PYG{l+s+s2}{\PYGZdq{}}\PYG{l+s+s2}{Acceleration Vectors}\PYG{l+s+s2}{\PYGZdq{}}\PYG{p}{)}

\PYG{c+c1}{\PYGZsh{} Labels and Title}
\PYG{n}{plt}\PYG{o}{.}\PYG{n}{xlabel}\PYG{p}{(}\PYG{l+s+s1}{\PYGZsq{}}\PYG{l+s+s1}{x}\PYG{l+s+s1}{\PYGZsq{}}\PYG{p}{)}
\PYG{n}{plt}\PYG{o}{.}\PYG{n}{ylabel}\PYG{p}{(}\PYG{l+s+s1}{\PYGZsq{}}\PYG{l+s+s1}{y}\PYG{l+s+s1}{\PYGZsq{}}\PYG{p}{)}
\PYG{n}{plt}\PYG{o}{.}\PYG{n}{title}\PYG{p}{(}\PYG{l+s+s1}{\PYGZsq{}}\PYG{l+s+s1}{Particle Motion: Parametric Curve with Velocity and Acceleration Vectors}\PYG{l+s+s1}{\PYGZsq{}}\PYG{p}{)}
\PYG{n}{plt}\PYG{o}{.}\PYG{n}{grid}\PYG{p}{(}\PYG{k+kc}{True}\PYG{p}{)}
\PYG{n}{plt}\PYG{o}{.}\PYG{n}{legend}\PYG{p}{(}\PYG{p}{)}
\PYG{n}{plt}\PYG{o}{.}\PYG{n}{show}\PYG{p}{(}\PYG{p}{)}

\PYG{c+c1}{\PYGZsh{} Calculate the length of the curve numerically using integration}
\PYG{n}{curve\PYGZus{}length}\PYG{p}{,} \PYG{n}{\PYGZus{}} \PYG{o}{=} \PYG{n}{quad}\PYG{p}{(}\PYG{n}{integrand}\PYG{p}{,} \PYG{l+m+mi}{0}\PYG{p}{,} \PYG{l+m+mi}{2}\PYG{p}{)}
\PYG{n+nb}{print}\PYG{p}{(}\PYG{l+s+sa}{f}\PYG{l+s+s2}{\PYGZdq{}}\PYG{l+s+s2}{The length of the curve traced by the particle between t = 0 and t = 2 is approximately }\PYG{l+s+si}{\PYGZob{}}\PYG{n}{curve\PYGZus{}length}\PYG{l+s+si}{:}\PYG{l+s+s2}{.2f}\PYG{l+s+si}{\PYGZcb{}}\PYG{l+s+s2}{.}\PYG{l+s+s2}{\PYGZdq{}}\PYG{p}{)}
\end{sphinxVerbatim}

\end{sphinxuseclass}\end{sphinxVerbatimInput}
\begin{sphinxVerbatimOutput}

\begin{sphinxuseclass}{cell_output}
\noindent\sphinxincludegraphics{{5c56e60418cbca51e6a8aab01052de095388c2cd89f5705da33b49b8554db597}.png}

\begin{sphinxVerbatim}[commandchars=\\\{\}]
The length of the curve traced by the particle between t = 0 and t = 2 is approximately 17.10.
\end{sphinxVerbatim}

\end{sphinxuseclass}\end{sphinxVerbatimOutput}

\end{sphinxuseclass}

\bigskip\hrule\bigskip



\section{Problem 4: Separable Differential Equation}
\label{\detokenize{prerequisites:problem-4-separable-differential-equation}}
\sphinxAtStartPar
Solve the following first\sphinxhyphen{}order separable ODE:
\begin{equation*}
\begin{split}
\frac{dy}{dx} = x^2 + x
\end{split}
\end{equation*}
\sphinxAtStartPar
with the initial condition \(y(1) = 3\).


\subsection{Solution:}
\label{\detokenize{prerequisites:id5}}
\sphinxAtStartPar
To solve, integrate both sides with respect to \(x\):
\begin{equation*}
\begin{split}
y = \int (x^2 + x) \, dx = \frac{x^3}{3} + \frac{x^2}{2} + C
\end{split}
\end{equation*}
\sphinxAtStartPar
Using the initial condition \(y(1) = 3\), substitute \(x = 1\) and \(y = 3\):
\begin{equation*}
\begin{split}
3 = \frac{1^3}{3} + \frac{1^2}{2} + C \quad \Rightarrow \quad 3 = \frac{1}{3} + \frac{1}{2} + C
\end{split}
\end{equation*}
\sphinxAtStartPar
Solving for \(C\), we get \(C = \frac{15}{6}\), so the solution is:
\begin{equation*}
\begin{split}
y = \frac{x^3}{3} + \frac{x^2}{2} + \frac{15}{6}
\end{split}
\end{equation*}

\bigskip\hrule\bigskip



\section{Problem 5: Exponential Growth}
\label{\detokenize{prerequisites:problem-5-exponential-growth}}
\sphinxAtStartPar
Solve the differential equation:
\begin{equation*}
\begin{split}
\frac{dy}{dx} = e^x + x
\end{split}
\end{equation*}
\sphinxAtStartPar
with the initial condition \(y(0) = 10\).


\subsection{Solution:}
\label{\detokenize{prerequisites:id6}}
\sphinxAtStartPar
To solve, integrate both sides with respect to \(x\):
\begin{equation*}
\begin{split}
y = \int (e^x + x) \, dx = e^x + \frac{x^2}{2} + C
\end{split}
\end{equation*}
\sphinxAtStartPar
Using the initial condition \(y(0) = 10\):
\begin{equation*}
\begin{split}
10 = e^0 + \frac{0^2}{2} + C \quad \Rightarrow \quad 10 = 1 + C
\end{split}
\end{equation*}
\sphinxAtStartPar
Thus, \(C = 9\), and the solution is:
\begin{equation*}
\begin{split}
y = e^x + \frac{x^2}{2} + 9
\end{split}
\end{equation*}

\bigskip\hrule\bigskip



\section{Problem 6: Separable Differential Equation}
\label{\detokenize{prerequisites:problem-6-separable-differential-equation}}
\sphinxAtStartPar
Solve the differential equation:
\begin{equation*}
\begin{split}
\frac{dy}{dx} = \frac{1}{y}
\end{split}
\end{equation*}
\sphinxAtStartPar
with the initial condition \(y(1) = 2\).


\subsection{Solution:}
\label{\detokenize{prerequisites:id7}}
\sphinxAtStartPar
This equation is separable, so we can rewrite it as:
\begin{equation*}
\begin{split}
y \, dy = dx
\end{split}
\end{equation*}
\sphinxAtStartPar
Next, integrate both sides:
\begin{equation*}
\begin{split}
\int y \, dy = \int dx
\end{split}
\end{equation*}
\sphinxAtStartPar
This gives:
\begin{equation*}
\begin{split}
\frac{y^2}{2} = x + C
\end{split}
\end{equation*}
\sphinxAtStartPar
Now, multiply by 2 to simplify:
\begin{equation*}
\begin{split}
y^2 = 2x + 2C
\end{split}
\end{equation*}
\sphinxAtStartPar
Using the initial condition \(y(1) = 2\), substitute \(x = 1\) and \(y = 2\):
\begin{equation*}
\begin{split}
2^2 = 2(1) + 2C \quad \Rightarrow \quad 4 = 2 + 2C
\end{split}
\end{equation*}
\sphinxAtStartPar
Solving for \(C\), we get:
\begin{equation*}
\begin{split}
C = 1
\end{split}
\end{equation*}
\sphinxAtStartPar
Thus, the solution is:
\begin{equation*}
\begin{split}
y^2 = 2x + 2
\end{split}
\end{equation*}
\sphinxAtStartPar
Taking the square root of both sides:
\begin{equation*}
\begin{split}
y = \pm \sqrt{2x + 2}
\end{split}
\end{equation*}
\sphinxAtStartPar
For the positive root, the solution is:
\begin{equation*}
\begin{split}
y = \sqrt{2x + 2}
\end{split}
\end{equation*}

\bigskip\hrule\bigskip



\section{Problem 7: Linear and Non\sphinxhyphen{}Linear Differential Equations}
\label{\detokenize{prerequisites:problem-7-linear-and-non-linear-differential-equations}}
\sphinxAtStartPar
Consider the following two differential equations:
\begin{enumerate}
\sphinxsetlistlabels{\arabic}{enumi}{enumii}{}{.}%
\item {} 
\sphinxAtStartPar
\sphinxstylestrong{Linear Differential Equation}:\\
Solve the linear first\sphinxhyphen{}order differential equation:
\begin{equation*}
\begin{split}
   \frac{dy}{dx} + y = e^x
   \end{split}
\end{equation*}
\item {} 
\sphinxAtStartPar
\sphinxstylestrong{Non\sphinxhyphen{}Linear Differential Equation}:\\
Solve the non\sphinxhyphen{}linear differential equation:
\begin{equation*}
\begin{split}
   \frac{dy}{dx} = y^2 - x
   \end{split}
\end{equation*}
\end{enumerate}


\subsection{Solution:}
\label{\detokenize{prerequisites:id8}}

\subsubsection{1. \sphinxstylestrong{Linear Differential Equation}}
\label{\detokenize{prerequisites:linear-differential-equation}}
\sphinxAtStartPar
The equation is of the form:
\begin{equation*}
\begin{split}
\frac{dy}{dx} + P(x) y = Q(x)
\end{split}
\end{equation*}
\sphinxAtStartPar
where \(P(x) = 1\) and \(Q(x) = e^x\).

\sphinxAtStartPar
\sphinxstylestrong{Step 1}: Find the integrating factor \(\mu(x)\):
\begin{equation*}
\begin{split}
\mu(x) = e^{\int P(x) \, dx} = e^{\int 1 \, dx} = e^x
\end{split}
\end{equation*}
\sphinxAtStartPar
\sphinxstylestrong{Step 2}: Multiply both sides of the differential equation by the integrating factor \(\mu(x)\):
\begin{equation*}
\begin{split}
e^x \frac{dy}{dx} + e^x y = e^{2x}
\end{split}
\end{equation*}
\sphinxAtStartPar
\sphinxstylestrong{Step 3}: The left\sphinxhyphen{}hand side is the derivative of \(e^x y\), so we can rewrite the equation as:
\begin{equation*}
\begin{split}
\frac{d}{dx} \left( e^x y \right) = e^{2x}
\end{split}
\end{equation*}
\sphinxAtStartPar
\sphinxstylestrong{Step 4}: Integrate both sides:
\begin{equation*}
\begin{split}
e^x y = \int e^{2x} \, dx = \frac{e^{2x}}{2} + C
\end{split}
\end{equation*}
\sphinxAtStartPar
\sphinxstylestrong{Step 5}: Solve for \(y\):
\begin{equation*}
\begin{split}
y = \frac{e^x}{2} + Ce^{-x}
\end{split}
\end{equation*}
\sphinxAtStartPar
Thus, the general solution is:
\begin{equation*}
\begin{split}
y = \frac{e^x}{2} + Ce^{-x}
\end{split}
\end{equation*}

\bigskip\hrule\bigskip



\subsubsection{2. \sphinxstylestrong{Non\sphinxhyphen{}Linear Differential Equation}}
\label{\detokenize{prerequisites:non-linear-differential-equation}}
\sphinxAtStartPar
The given equation is:
\begin{equation*}
\begin{split}
\frac{dy}{dx} = y^2 - x
\end{split}
\end{equation*}
\sphinxAtStartPar
This is a \sphinxstylestrong{non\sphinxhyphen{}linear} differential equation and does not have a straightforward analytical solution. However, it can be solved using numerical methods, such as \sphinxstylestrong{Euler’s method}, or qualitative techniques like \sphinxstylestrong{direction fields} to analyze the behavior of the solutions.

\begin{sphinxadmonition}{note}{Note:}
\sphinxAtStartPar
While there is no simple analytic solution to present here, numerical solutions can be computed based on initial conditions.

\sphinxAtStartPar
We use the \sphinxstylestrong{Runge\sphinxhyphen{}Kutta method} to solve the differential equation:
\end{sphinxadmonition}
\begin{equation*}
\begin{split}
\frac{dy}{dx} = y^2 - x
\end{split}
\end{equation*}

\paragraph{Python Code:}
\label{\detokenize{prerequisites:python-code}}
\begin{sphinxuseclass}{cell}\begin{sphinxVerbatimInput}

\begin{sphinxuseclass}{cell_input}
\begin{sphinxVerbatim}[commandchars=\\\{\}]
\PYG{k+kn}{import} \PYG{n+nn}{numpy} \PYG{k}{as} \PYG{n+nn}{np}
\PYG{k+kn}{import} \PYG{n+nn}{matplotlib}\PYG{n+nn}{.}\PYG{n+nn}{pyplot} \PYG{k}{as} \PYG{n+nn}{plt}

\PYG{c+c1}{\PYGZsh{} Function representing the ODE dy/dx = y\PYGZca{}2 \PYGZhy{} x}
\PYG{k}{def} \PYG{n+nf}{dydx}\PYG{p}{(}\PYG{n}{x}\PYG{p}{,} \PYG{n}{y}\PYG{p}{)}\PYG{p}{:}
    \PYG{k}{return} \PYG{n}{y}\PYG{o}{*}\PYG{o}{*}\PYG{l+m+mi}{2} \PYG{o}{\PYGZhy{}} \PYG{n}{x}

\PYG{c+c1}{\PYGZsh{} Fourth\PYGZhy{}order Runge\PYGZhy{}Kutta method for numerical solution of ODE}
\PYG{k}{def} \PYG{n+nf}{runge\PYGZus{}kutta\PYGZus{}method}\PYG{p}{(}\PYG{n}{dydx}\PYG{p}{,} \PYG{n}{x0}\PYG{p}{,} \PYG{n}{y0}\PYG{p}{,} \PYG{n}{x\PYGZus{}end}\PYG{p}{,} \PYG{n}{h}\PYG{p}{)}\PYG{p}{:}
    \PYG{n}{x\PYGZus{}values} \PYG{o}{=} \PYG{n}{np}\PYG{o}{.}\PYG{n}{arange}\PYG{p}{(}\PYG{n}{x0}\PYG{p}{,} \PYG{n}{x\PYGZus{}end} \PYG{o}{+} \PYG{n}{h}\PYG{p}{,} \PYG{n}{h}\PYG{p}{)}
    \PYG{n}{y\PYGZus{}values} \PYG{o}{=} \PYG{n}{np}\PYG{o}{.}\PYG{n}{zeros}\PYG{p}{(}\PYG{n+nb}{len}\PYG{p}{(}\PYG{n}{x\PYGZus{}values}\PYG{p}{)}\PYG{p}{)}
    \PYG{n}{y\PYGZus{}values}\PYG{p}{[}\PYG{l+m+mi}{0}\PYG{p}{]} \PYG{o}{=} \PYG{n}{y0}
    
    \PYG{k}{for} \PYG{n}{i} \PYG{o+ow}{in} \PYG{n+nb}{range}\PYG{p}{(}\PYG{l+m+mi}{1}\PYG{p}{,} \PYG{n+nb}{len}\PYG{p}{(}\PYG{n}{x\PYGZus{}values}\PYG{p}{)}\PYG{p}{)}\PYG{p}{:}
        \PYG{n}{x} \PYG{o}{=} \PYG{n}{x\PYGZus{}values}\PYG{p}{[}\PYG{n}{i}\PYG{o}{\PYGZhy{}}\PYG{l+m+mi}{1}\PYG{p}{]}
        \PYG{n}{y} \PYG{o}{=} \PYG{n}{y\PYGZus{}values}\PYG{p}{[}\PYG{n}{i}\PYG{o}{\PYGZhy{}}\PYG{l+m+mi}{1}\PYG{p}{]}
        
        \PYG{n}{k1} \PYG{o}{=} \PYG{n}{h} \PYG{o}{*} \PYG{n}{dydx}\PYG{p}{(}\PYG{n}{x}\PYG{p}{,} \PYG{n}{y}\PYG{p}{)}
        \PYG{n}{k2} \PYG{o}{=} \PYG{n}{h} \PYG{o}{*} \PYG{n}{dydx}\PYG{p}{(}\PYG{n}{x} \PYG{o}{+} \PYG{n}{h}\PYG{o}{/}\PYG{l+m+mi}{2}\PYG{p}{,} \PYG{n}{y} \PYG{o}{+} \PYG{n}{k1}\PYG{o}{/}\PYG{l+m+mi}{2}\PYG{p}{)}
        \PYG{n}{k3} \PYG{o}{=} \PYG{n}{h} \PYG{o}{*} \PYG{n}{dydx}\PYG{p}{(}\PYG{n}{x} \PYG{o}{+} \PYG{n}{h}\PYG{o}{/}\PYG{l+m+mi}{2}\PYG{p}{,} \PYG{n}{y} \PYG{o}{+} \PYG{n}{k2}\PYG{o}{/}\PYG{l+m+mi}{2}\PYG{p}{)}
        \PYG{n}{k4} \PYG{o}{=} \PYG{n}{h} \PYG{o}{*} \PYG{n}{dydx}\PYG{p}{(}\PYG{n}{x} \PYG{o}{+} \PYG{n}{h}\PYG{p}{,} \PYG{n}{y} \PYG{o}{+} \PYG{n}{k3}\PYG{p}{)}
        
        \PYG{n}{y\PYGZus{}values}\PYG{p}{[}\PYG{n}{i}\PYG{p}{]} \PYG{o}{=} \PYG{n}{y} \PYG{o}{+} \PYG{p}{(}\PYG{n}{k1} \PYG{o}{+} \PYG{l+m+mi}{2}\PYG{o}{*}\PYG{n}{k2} \PYG{o}{+} \PYG{l+m+mi}{2}\PYG{o}{*}\PYG{n}{k3} \PYG{o}{+} \PYG{n}{k4}\PYG{p}{)} \PYG{o}{/} \PYG{l+m+mi}{6}
    
    \PYG{k}{return} \PYG{n}{x\PYGZus{}values}\PYG{p}{,} \PYG{n}{y\PYGZus{}values}

\PYG{c+c1}{\PYGZsh{} Adjust the domain to avoid large values causing overflow}
\PYG{n}{x0} \PYG{o}{=} \PYG{l+m+mi}{0}  \PYG{c+c1}{\PYGZsh{} Initial x}
\PYG{n}{y0} \PYG{o}{=} \PYG{l+m+mi}{1}  \PYG{c+c1}{\PYGZsh{} Initial y}
\PYG{n}{x\PYGZus{}end\PYGZus{}new} \PYG{o}{=} \PYG{l+m+mi}{1}  \PYG{c+c1}{\PYGZsh{} Limit the domain to avoid instability}
\PYG{n}{h} \PYG{o}{=} \PYG{l+m+mf}{0.01}  \PYG{c+c1}{\PYGZsh{} Step size}

\PYG{c+c1}{\PYGZsh{} Solve the ODE using Runge\PYGZhy{}Kutta method with a smaller domain}
\PYG{n}{x\PYGZus{}values\PYGZus{}rk\PYGZus{}new}\PYG{p}{,} \PYG{n}{y\PYGZus{}values\PYGZus{}rk\PYGZus{}new} \PYG{o}{=} \PYG{n}{runge\PYGZus{}kutta\PYGZus{}method}\PYG{p}{(}\PYG{n}{dydx}\PYG{p}{,} \PYG{n}{x0}\PYG{p}{,} \PYG{n}{y0}\PYG{p}{,} \PYG{n}{x\PYGZus{}end\PYGZus{}new}\PYG{p}{,} \PYG{n}{h}\PYG{p}{)}

\PYG{c+c1}{\PYGZsh{} Plot the solution with the new domain}
\PYG{n}{plt}\PYG{o}{.}\PYG{n}{plot}\PYG{p}{(}\PYG{n}{x\PYGZus{}values\PYGZus{}rk\PYGZus{}new}\PYG{p}{,} \PYG{n}{y\PYGZus{}values\PYGZus{}rk\PYGZus{}new}\PYG{p}{,} \PYG{n}{label}\PYG{o}{=}\PYG{l+s+s2}{\PYGZdq{}}\PYG{l+s+s2}{Numerical Solution (Runge\PYGZhy{}Kutta Method)}\PYG{l+s+s2}{\PYGZdq{}}\PYG{p}{,} \PYG{n}{color}\PYG{o}{=}\PYG{l+s+s2}{\PYGZdq{}}\PYG{l+s+s2}{orange}\PYG{l+s+s2}{\PYGZdq{}}\PYG{p}{)}
\PYG{n}{plt}\PYG{o}{.}\PYG{n}{xlabel}\PYG{p}{(}\PYG{l+s+s1}{\PYGZsq{}}\PYG{l+s+s1}{x}\PYG{l+s+s1}{\PYGZsq{}}\PYG{p}{)}
\PYG{n}{plt}\PYG{o}{.}\PYG{n}{ylabel}\PYG{p}{(}\PYG{l+s+s1}{\PYGZsq{}}\PYG{l+s+s1}{y}\PYG{l+s+s1}{\PYGZsq{}}\PYG{p}{)}
\PYG{n}{plt}\PYG{o}{.}\PYG{n}{title}\PYG{p}{(}\PYG{l+s+s2}{\PYGZdq{}}\PYG{l+s+s2}{Solution of \PYGZdl{}dy/dx = y\PYGZca{}2 \PYGZhy{} x\PYGZdl{} using Runge\PYGZhy{}Kutta Method}\PYG{l+s+s2}{\PYGZdq{}}\PYG{p}{)}
\PYG{n}{plt}\PYG{o}{.}\PYG{n}{legend}\PYG{p}{(}\PYG{p}{)}
\PYG{n}{plt}\PYG{o}{.}\PYG{n}{grid}\PYG{p}{(}\PYG{k+kc}{True}\PYG{p}{)}
\PYG{n}{plt}\PYG{o}{.}\PYG{n}{show}\PYG{p}{(}\PYG{p}{)}
\end{sphinxVerbatim}

\end{sphinxuseclass}\end{sphinxVerbatimInput}
\begin{sphinxVerbatimOutput}

\begin{sphinxuseclass}{cell_output}
\noindent\sphinxincludegraphics{{6d32a946f3fb858c048e4b90de296c287fc689ac1c4c2557923b012e961bf031}.png}

\end{sphinxuseclass}\end{sphinxVerbatimOutput}

\end{sphinxuseclass}

\bigskip\hrule\bigskip



\section{Problem 8: Boundary Value Problem with Initial Conditions}
\label{\detokenize{prerequisites:problem-8-boundary-value-problem-with-initial-conditions}}
\sphinxAtStartPar
Solve the following second\sphinxhyphen{}order boundary value problem (BVP):
\begin{equation*}
\begin{split}
\frac{d^2 y}{dx^2} = -k^2 y
\end{split}
\end{equation*}
\sphinxAtStartPar
with the boundary conditions:
\begin{equation*}
\begin{split}
y(0) = 0 \quad \text{and} \quad y(L) = 0
\end{split}
\end{equation*}
\sphinxAtStartPar
where \(k\) is a constant and \(L\) is the length of the domain.


\subsection{Solution:}
\label{\detokenize{prerequisites:id9}}

\subsubsection{Step 1: General Solution of the ODE}
\label{\detokenize{prerequisites:step-1-general-solution-of-the-ode}}
\sphinxAtStartPar
The given equation is a second\sphinxhyphen{}order homogeneous linear differential equation with constant coefficients:
\begin{equation*}
\begin{split}
\frac{d^2 y}{dx^2} = -k^2 y
\end{split}
\end{equation*}
\sphinxAtStartPar
The characteristic equation for this differential equation is:
\begin{equation*}
\begin{split}
r^2 + k^2 = 0
\end{split}
\end{equation*}
\sphinxAtStartPar
This gives the roots:
\begin{equation*}
\begin{split}
r = \pm ik
\end{split}
\end{equation*}
\sphinxAtStartPar
So, the general solution for the differential equation is:
\begin{equation*}
\begin{split}
y(x) = C_1 \sin(kx) + C_2 \cos(kx)
\end{split}
\end{equation*}

\subsubsection{Step 2: Apply the Boundary Conditions}
\label{\detokenize{prerequisites:step-2-apply-the-boundary-conditions}}\begin{enumerate}
\sphinxsetlistlabels{\arabic}{enumi}{enumii}{}{.}%
\item {} 
\sphinxAtStartPar
\sphinxstylestrong{First Boundary Condition: \(y(0) = 0\)}

\end{enumerate}

\sphinxAtStartPar
Substitute \(x = 0\) into the general solution:
\begin{equation*}
\begin{split}
y(0) = C_1 \sin(0) + C_2 \cos(0) = 0 \quad \Rightarrow \quad C_2 = 0
\end{split}
\end{equation*}
\sphinxAtStartPar
Thus, the solution simplifies to:
\begin{equation*}
\begin{split}
y(x) = C_1 \sin(kx)
\end{split}
\end{equation*}\begin{enumerate}
\sphinxsetlistlabels{\arabic}{enumi}{enumii}{}{.}%
\setcounter{enumi}{1}
\item {} 
\sphinxAtStartPar
\sphinxstylestrong{Second Boundary Condition: \(y(L) = 0\)}

\end{enumerate}

\sphinxAtStartPar
Now, substitute \(x = L\) into the simplified solution:
\begin{equation*}
\begin{split}
y(L) = C_1 \sin(kL) = 0
\end{split}
\end{equation*}
\sphinxAtStartPar
For non\sphinxhyphen{}trivial solutions (i.e., \(C_1 \neq 0\)), we require:
\begin{equation*}
\begin{split}
\sin(kL) = 0
\end{split}
\end{equation*}
\sphinxAtStartPar
This implies:
\begin{equation*}
\begin{split}
kL = n\pi \quad \Rightarrow \quad k = \frac{n\pi}{L}
\end{split}
\end{equation*}
\sphinxAtStartPar
where \(n\) is an integer (i.e., \(n = 1, 2, 3, \dots\)).


\subsubsection{Step 3: Final Solution}
\label{\detokenize{prerequisites:step-3-final-solution}}
\sphinxAtStartPar
The final solution for the boundary value problem is:
\begin{equation*}
\begin{split}
y(x) = C_1 \sin\left( \frac{n\pi}{L} x \right)
\end{split}
\end{equation*}
\sphinxAtStartPar
where \(C_1\) is an arbitrary constant and \(n\) is a positive integer.


\bigskip\hrule\bigskip



\section{Problem 9: Vector Field and Gradient}
\label{\detokenize{prerequisites:problem-9-vector-field-and-gradient}}
\sphinxAtStartPar
Given the scalar function:
\begin{equation*}
\begin{split}
\phi(x, y, z) = x^2 y + yz^2
\end{split}
\end{equation*}\begin{enumerate}
\sphinxsetlistlabels{\arabic}{enumi}{enumii}{}{.}%
\item {} 
\sphinxAtStartPar
Find the \sphinxstylestrong{gradient} of the scalar field \(\phi(x, y, z)\).

\item {} 
\sphinxAtStartPar
Plot the vector field representing the gradient in the \(xy\)\sphinxhyphen{}plane for \(z = 0\).

\end{enumerate}


\subsection{Solution:}
\label{\detokenize{prerequisites:id10}}

\subsubsection{1. \sphinxstylestrong{Gradient of the Scalar Field}}
\label{\detokenize{prerequisites:gradient-of-the-scalar-field}}
\sphinxAtStartPar
The gradient of a scalar field \(\phi(x, y, z)\) is given by:
\begin{equation*}
\begin{split}
\nabla \phi = \left( \frac{\partial \phi}{\partial x}, \frac{\partial \phi}{\partial y}, \frac{\partial \phi}{\partial z} \right)
\end{split}
\end{equation*}
\sphinxAtStartPar
For \(\phi(x, y, z) = x^2 y + yz^2\), the partial derivatives are:
\begin{equation*}
\begin{split}
\frac{\partial \phi}{\partial x} = 2xy
\end{split}
\end{equation*}\begin{equation*}
\begin{split}
\frac{\partial \phi}{\partial y} = x^2 + z^2
\end{split}
\end{equation*}\begin{equation*}
\begin{split}
\frac{\partial \phi}{\partial z} = 2yz
\end{split}
\end{equation*}
\sphinxAtStartPar
Thus, the gradient vector field is:
\begin{equation*}
\begin{split}
\nabla \phi = (2xy, x^2 + z^2, 2yz)
\end{split}
\end{equation*}

\subsubsection{2. \sphinxstylestrong{Plotting the Gradient Field in the \protect\(xy\protect\)\sphinxhyphen{}Plane}}
\label{\detokenize{prerequisites:plotting-the-gradient-field-in-the-xy-plane}}
\sphinxAtStartPar
Since we’re plotting in the \(xy\)\sphinxhyphen{}plane for \(z = 0\), the gradient simplifies to:
\begin{equation*}
\begin{split}
\nabla \phi(x, y, 0) = (2xy, x^2, 0)
\end{split}
\end{equation*}
\sphinxAtStartPar
Here is the Python code to plot the vector field in the \(xy\)\sphinxhyphen{}plane:

\begin{sphinxuseclass}{cell}\begin{sphinxVerbatimInput}

\begin{sphinxuseclass}{cell_input}
\begin{sphinxVerbatim}[commandchars=\\\{\}]
\PYG{k+kn}{import} \PYG{n+nn}{numpy} \PYG{k}{as} \PYG{n+nn}{np}
\PYG{k+kn}{import} \PYG{n+nn}{matplotlib}\PYG{n+nn}{.}\PYG{n+nn}{pyplot} \PYG{k}{as} \PYG{n+nn}{plt}

\PYG{c+c1}{\PYGZsh{} Define the scalar field gradient in the xy\PYGZhy{}plane (z=0)}
\PYG{k}{def} \PYG{n+nf}{grad\PYGZus{}phi}\PYG{p}{(}\PYG{n}{x}\PYG{p}{,} \PYG{n}{y}\PYG{p}{)}\PYG{p}{:}
    \PYG{n}{grad\PYGZus{}x} \PYG{o}{=} \PYG{l+m+mi}{2} \PYG{o}{*} \PYG{n}{x} \PYG{o}{*} \PYG{n}{y}
    \PYG{n}{grad\PYGZus{}y} \PYG{o}{=} \PYG{n}{x}\PYG{o}{*}\PYG{o}{*}\PYG{l+m+mi}{2}
    \PYG{k}{return} \PYG{n}{grad\PYGZus{}x}\PYG{p}{,} \PYG{n}{grad\PYGZus{}y}

\PYG{c+c1}{\PYGZsh{} Create a grid of points}
\PYG{n}{x\PYGZus{}vals} \PYG{o}{=} \PYG{n}{np}\PYG{o}{.}\PYG{n}{linspace}\PYG{p}{(}\PYG{o}{\PYGZhy{}}\PYG{l+m+mi}{2}\PYG{p}{,} \PYG{l+m+mi}{2}\PYG{p}{,} \PYG{l+m+mi}{20}\PYG{p}{)}
\PYG{n}{y\PYGZus{}vals} \PYG{o}{=} \PYG{n}{np}\PYG{o}{.}\PYG{n}{linspace}\PYG{p}{(}\PYG{o}{\PYGZhy{}}\PYG{l+m+mi}{2}\PYG{p}{,} \PYG{l+m+mi}{2}\PYG{p}{,} \PYG{l+m+mi}{20}\PYG{p}{)}
\PYG{n}{X}\PYG{p}{,} \PYG{n}{Y} \PYG{o}{=} \PYG{n}{np}\PYG{o}{.}\PYG{n}{meshgrid}\PYG{p}{(}\PYG{n}{x\PYGZus{}vals}\PYG{p}{,} \PYG{n}{y\PYGZus{}vals}\PYG{p}{)}

\PYG{c+c1}{\PYGZsh{} Compute the gradient at each point}
\PYG{n}{U}\PYG{p}{,} \PYG{n}{V} \PYG{o}{=} \PYG{n}{grad\PYGZus{}phi}\PYG{p}{(}\PYG{n}{X}\PYG{p}{,} \PYG{n}{Y}\PYG{p}{)}

\PYG{c+c1}{\PYGZsh{} Plot the vector field}
\PYG{n}{plt}\PYG{o}{.}\PYG{n}{figure}\PYG{p}{(}\PYG{n}{figsize}\PYG{o}{=}\PYG{p}{(}\PYG{l+m+mi}{6}\PYG{p}{,} \PYG{l+m+mi}{6}\PYG{p}{)}\PYG{p}{)}
\PYG{n}{plt}\PYG{o}{.}\PYG{n}{quiver}\PYG{p}{(}\PYG{n}{X}\PYG{p}{,} \PYG{n}{Y}\PYG{p}{,} \PYG{n}{U}\PYG{p}{,} \PYG{n}{V}\PYG{p}{)}
\PYG{n}{plt}\PYG{o}{.}\PYG{n}{xlabel}\PYG{p}{(}\PYG{l+s+s1}{\PYGZsq{}}\PYG{l+s+s1}{x}\PYG{l+s+s1}{\PYGZsq{}}\PYG{p}{)}
\PYG{n}{plt}\PYG{o}{.}\PYG{n}{ylabel}\PYG{p}{(}\PYG{l+s+s1}{\PYGZsq{}}\PYG{l+s+s1}{y}\PYG{l+s+s1}{\PYGZsq{}}\PYG{p}{)}
\PYG{n}{plt}\PYG{o}{.}\PYG{n}{title}\PYG{p}{(}\PYG{l+s+s1}{\PYGZsq{}}\PYG{l+s+s1}{Gradient of \PYGZdl{}}\PYG{l+s+s1}{\PYGZbs{}}\PYG{l+s+s1}{phi(x, y, z) = x\PYGZca{}2 y + yz\PYGZca{}2\PYGZdl{} in the xy\PYGZhy{}plane (z=0)}\PYG{l+s+s1}{\PYGZsq{}}\PYG{p}{)}
\PYG{n}{plt}\PYG{o}{.}\PYG{n}{grid}\PYG{p}{(}\PYG{k+kc}{True}\PYG{p}{)}
\PYG{n}{plt}\PYG{o}{.}\PYG{n}{show}\PYG{p}{(}\PYG{p}{)}
\end{sphinxVerbatim}

\end{sphinxuseclass}\end{sphinxVerbatimInput}
\begin{sphinxVerbatimOutput}

\begin{sphinxuseclass}{cell_output}
\noindent\sphinxincludegraphics{{9a22b04d4976f8a4bc9558934ae499fdeaa4a2f78d40eb7ff07acbd18ef78dad}.png}

\end{sphinxuseclass}\end{sphinxVerbatimOutput}

\end{sphinxuseclass}

\bigskip\hrule\bigskip



\section{Problem 10: Line, Surface, and Volume Integrals}
\label{\detokenize{prerequisites:problem-10-line-surface-and-volume-integrals}}

\subsection{Problem Statement}
\label{\detokenize{prerequisites:id11}}\begin{enumerate}
\sphinxsetlistlabels{\arabic}{enumi}{enumii}{}{.}%
\item {} 
\sphinxAtStartPar
\sphinxstylestrong{Line Integral}: Given the vector field \(\mathbf{F}(x, y) = (x^2, y^2)\), compute the line integral of \(\mathbf{F} \) along the curve \(C\), where \(C\) is the straight line path from point \((0, 0)\) to point \((1, 1)\).

\item {} 
\sphinxAtStartPar
\sphinxstylestrong{Surface Integral}: Compute the surface integral of the vector field \(\mathbf{F}(x, y, z) = (y, z, x)\) over the surface \(S\), where \(S\) is the portion of the plane \(z = 4 - x - y\) in the first octant.

\item {} 
\sphinxAtStartPar
\sphinxstylestrong{Volume Integral}: Compute the volume integral of the scalar field \(f(x, y, z) = x + y + z\) over the region \(V\), where \(V\) is the region bounded by the cylinder \(x^2 + y^2 = 1\) and the planes \(z = 0\) and \(z = 3\).

\end{enumerate}


\subsection{Solution}
\label{\detokenize{prerequisites:id12}}

\subsubsection{1. Line Integral}
\label{\detokenize{prerequisites:line-integral}}
\sphinxAtStartPar
We want to compute the line integral of the vector field \(\mathbf{F}(x, y) = (x^2, y^2)\) along the curve \(C\), a straight line from \((0, 0)\) to \((1, 1)\).

\sphinxAtStartPar
The line integral is given by:
\begin{equation*}
\begin{split}
\int_C \mathbf{F} \cdot d\mathbf{r}
\end{split}
\end{equation*}
\sphinxAtStartPar
Parameterizing the line as \(x(t) = t\) and \(y(t) = t\), the differential displacement vector is \(d\mathbf{r} = (1, 1) dt\).

\sphinxAtStartPar
The integral becomes:
\begin{equation*}
\begin{split}
\int_0^1 (t^2 + t^2) dt = \int_0^1 2t^2 dt
\end{split}
\end{equation*}
\sphinxAtStartPar
The result can be calculated as:

\begin{sphinxuseclass}{cell}\begin{sphinxVerbatimInput}

\begin{sphinxuseclass}{cell_input}
\begin{sphinxVerbatim}[commandchars=\\\{\}]
\PYG{k+kn}{import} \PYG{n+nn}{numpy} \PYG{k}{as} \PYG{n+nn}{np}
\PYG{k+kn}{from} \PYG{n+nn}{scipy}\PYG{n+nn}{.}\PYG{n+nn}{integrate} \PYG{k+kn}{import} \PYG{n}{quad}

\PYG{c+c1}{\PYGZsh{} Define the parameterized vector field F along the curve x(t) = t, y(t) = t}
\PYG{k}{def} \PYG{n+nf}{integrand\PYGZus{}line}\PYG{p}{(}\PYG{n}{t}\PYG{p}{)}\PYG{p}{:}
    \PYG{k}{return} \PYG{l+m+mi}{2} \PYG{o}{*} \PYG{n}{t}\PYG{o}{*}\PYG{o}{*}\PYG{l+m+mi}{2}  \PYG{c+c1}{\PYGZsh{} (x\PYGZca{}2 + y\PYGZca{}2) = 2t\PYGZca{}2}

\PYG{c+c1}{\PYGZsh{} Compute the line integral from t = 0 to t = 1}
\PYG{n}{line\PYGZus{}integral}\PYG{p}{,} \PYG{n}{\PYGZus{}} \PYG{o}{=} \PYG{n}{quad}\PYG{p}{(}\PYG{n}{integrand\PYGZus{}line}\PYG{p}{,} \PYG{l+m+mi}{0}\PYG{p}{,} \PYG{l+m+mi}{1}\PYG{p}{)}
\PYG{n+nb}{print}\PYG{p}{(}\PYG{l+s+sa}{f}\PYG{l+s+s2}{\PYGZdq{}}\PYG{l+s+s2}{Line integral: }\PYG{l+s+si}{\PYGZob{}}\PYG{n}{line\PYGZus{}integral}\PYG{l+s+si}{:}\PYG{l+s+s2}{.2f}\PYG{l+s+si}{\PYGZcb{}}\PYG{l+s+s2}{\PYGZdq{}}\PYG{p}{)}
\end{sphinxVerbatim}

\end{sphinxuseclass}\end{sphinxVerbatimInput}
\begin{sphinxVerbatimOutput}

\begin{sphinxuseclass}{cell_output}
\begin{sphinxVerbatim}[commandchars=\\\{\}]
Line integral: 0.67
\end{sphinxVerbatim}

\end{sphinxuseclass}\end{sphinxVerbatimOutput}

\end{sphinxuseclass}
\sphinxAtStartPar
The line integral of the vector field ( \textbackslash{}mathbf\{F\}(x, y) = (x\textasciicircum{}2, y\textasciicircum{}2) ) along the curve is:
\begin{equation*}
\begin{split}
\int_C \mathbf{F} \cdot d\mathbf{r} = \frac{2}{3}
\end{split}
\end{equation*}

\subsubsection{2. Surface Integral}
\label{\detokenize{prerequisites:surface-integral}}
\sphinxAtStartPar
We are asked to compute the surface integral of the vector field ( \textbackslash{}mathbf\{F\}(x, y, z) = (y, z, x) ) over the surface ( S ), which is the portion of the plane ( z = 4 \sphinxhyphen{} x \sphinxhyphen{} y ) in the first octant.


\paragraph{Step 1: Surface Integral Formula}
\label{\detokenize{prerequisites:step-1-surface-integral-formula}}
\sphinxAtStartPar
The surface integral of a vector field ( \textbackslash{}mathbf\{F\} ) over a surface ( S ) is given by:
\begin{equation*}
\begin{split}
\iint_S \mathbf{F} \cdot \mathbf{n} \, dS
\end{split}
\end{equation*}
\sphinxAtStartPar
where ( \textbackslash{}mathbf\{n\} ) is the unit normal vector to the surface, and ( dS ) is the differential surface area element.


\paragraph{Step 2: Parameterization and Normal Vector}
\label{\detokenize{prerequisites:step-2-parameterization-and-normal-vector}}
\sphinxAtStartPar
The surface ( S ) is described by the plane equation ( z = 4 \sphinxhyphen{} x \sphinxhyphen{} y ). This can be parameterized as:
\begin{equation*}
\begin{split}
\mathbf{r}(x, y) = (x, y, 4 - x - y)
\end{split}
\end{equation*}
\sphinxAtStartPar
The partial derivatives of ( \textbackslash{}mathbf\{r\} ) with respect to ( x ) and ( y ) give us the tangent vectors:
\begin{equation*}
\begin{split}
\frac{\partial \mathbf{r}}{\partial x} = (1, 0, -1) \quad \text{and} \quad \frac{\partial \mathbf{r}}{\partial y} = (0, 1, -1)
\end{split}
\end{equation*}
\sphinxAtStartPar
The cross product of these tangent vectors gives us the normal vector to the surface:
\begin{equation*}
\begin{split}
\mathbf{n} = \frac{\partial \mathbf{r}}{\partial x} \times \frac{\partial \mathbf{r}}{\partial y} = (-1, -1, 1)
\end{split}
\end{equation*}

\paragraph{Step 3: Dot Product of ( \textbackslash{}mathbf\{F\} ) and ( \textbackslash{}mathbf\{n\} )}
\label{\detokenize{prerequisites:step-3-dot-product-of-mathbf-f-and-mathbf-n}}
\sphinxAtStartPar
The vector field is ( \textbackslash{}mathbf\{F\}(x, y, z) = (y, z, x) ), so we need to compute the dot product:
\begin{equation*}
\begin{split}
\mathbf{F} \cdot \mathbf{n} = (y, z, x) \cdot (-1, -1, 1) = -y - z + x
\end{split}
\end{equation*}
\sphinxAtStartPar
Substitute ( z = 4 \sphinxhyphen{} x \sphinxhyphen{} y ) into the dot product:
\begin{equation*}
\begin{split}
\mathbf{F} \cdot \mathbf{n} = -y - (4 - x - y) + x = 2x - 4
\end{split}
\end{equation*}

\paragraph{Step 4: Limits of Integration}
\label{\detokenize{prerequisites:step-4-limits-of-integration}}
\sphinxAtStartPar
The surface lies in the first octant, so the limits for ( x ) and ( y ) are:
\begin{itemize}
\item {} 
\sphinxAtStartPar
( x ) ranges from 0 to 4.

\item {} 
\sphinxAtStartPar
For a given ( x ), ( y ) ranges from 0 to ( 4 \sphinxhyphen{} x ).

\end{itemize}


\paragraph{Step 5: Set Up the Surface Integral}
\label{\detokenize{prerequisites:step-5-set-up-the-surface-integral}}
\sphinxAtStartPar
The surface integral becomes:
\begin{equation*}
\begin{split}
\iint_S (2x - 4) \, dS = \int_0^4 \int_0^{4-x} (2x - 4) \, dy \, dx
\end{split}
\end{equation*}

\paragraph{Step 6: Solve the Surface Integral Numerically}
\label{\detokenize{prerequisites:step-6-solve-the-surface-integral-numerically}}
\sphinxAtStartPar
We can use Python to compute the surface integral:

\begin{sphinxuseclass}{cell}\begin{sphinxVerbatimInput}

\begin{sphinxuseclass}{cell_input}
\begin{sphinxVerbatim}[commandchars=\\\{\}]
\PYG{k+kn}{from} \PYG{n+nn}{scipy}\PYG{n+nn}{.}\PYG{n+nn}{integrate} \PYG{k+kn}{import} \PYG{n}{dblquad}

\PYG{c+c1}{\PYGZsh{} Define the integrand for the surface integral (2x \PYGZhy{} 4)}
\PYG{k}{def} \PYG{n+nf}{integrand\PYGZus{}surface}\PYG{p}{(}\PYG{n}{y}\PYG{p}{,} \PYG{n}{x}\PYG{p}{)}\PYG{p}{:}
    \PYG{k}{return} \PYG{l+m+mi}{2} \PYG{o}{*} \PYG{n}{x} \PYG{o}{\PYGZhy{}} \PYG{l+m+mi}{4}

\PYG{c+c1}{\PYGZsh{} Limits for the region of integration}
\PYG{n}{x\PYGZus{}min}\PYG{p}{,} \PYG{n}{x\PYGZus{}max} \PYG{o}{=} \PYG{l+m+mi}{0}\PYG{p}{,} \PYG{l+m+mi}{4}
\PYG{n}{y\PYGZus{}min} \PYG{o}{=} \PYG{k}{lambda} \PYG{n}{x}\PYG{p}{:} \PYG{l+m+mi}{0}
\PYG{n}{y\PYGZus{}max} \PYG{o}{=} \PYG{k}{lambda} \PYG{n}{x}\PYG{p}{:} \PYG{l+m+mi}{4} \PYG{o}{\PYGZhy{}} \PYG{n}{x}

\PYG{c+c1}{\PYGZsh{} Compute the surface integral}
\PYG{n}{surface\PYGZus{}integral}\PYG{p}{,} \PYG{n}{\PYGZus{}} \PYG{o}{=} \PYG{n}{dblquad}\PYG{p}{(}\PYG{n}{integrand\PYGZus{}surface}\PYG{p}{,} \PYG{n}{x\PYGZus{}min}\PYG{p}{,} \PYG{n}{x\PYGZus{}max}\PYG{p}{,} \PYG{n}{y\PYGZus{}min}\PYG{p}{,} \PYG{n}{y\PYGZus{}max}\PYG{p}{)}
\PYG{n+nb}{print}\PYG{p}{(}\PYG{l+s+sa}{f}\PYG{l+s+s2}{\PYGZdq{}}\PYG{l+s+s2}{Surface integral: }\PYG{l+s+si}{\PYGZob{}}\PYG{n}{surface\PYGZus{}integral}\PYG{l+s+si}{:}\PYG{l+s+s2}{.2f}\PYG{l+s+si}{\PYGZcb{}}\PYG{l+s+s2}{\PYGZdq{}}\PYG{p}{)}
\end{sphinxVerbatim}

\end{sphinxuseclass}\end{sphinxVerbatimInput}
\begin{sphinxVerbatimOutput}

\begin{sphinxuseclass}{cell_output}
\begin{sphinxVerbatim}[commandchars=\\\{\}]
Surface integral: \PYGZhy{}10.67
\end{sphinxVerbatim}

\end{sphinxuseclass}\end{sphinxVerbatimOutput}

\end{sphinxuseclass}

\paragraph{Step 7: Final Answer}
\label{\detokenize{prerequisites:step-7-final-answer}}
\sphinxAtStartPar
After performing the integration, the result of the surface integral is:
\begin{equation*}
\begin{split}
\iint_S \mathbf{F} \cdot \mathbf{n} \, dS = -\frac{32}{3}
\end{split}
\end{equation*}

\subsubsection{3. Volume Integral}
\label{\detokenize{prerequisites:volume-integral}}
\sphinxAtStartPar
We are asked to compute the volume integral of the scalar field \( f(x, y, z) = x + y + z \) over the region \( V \), which is bounded by the cylinder \( x^2 + y^2 = 1 \) and the planes \( z = 0 \) and \( z = 3 \).


\paragraph{Step 1: Volume Integral Formula}
\label{\detokenize{prerequisites:step-1-volume-integral-formula}}
\sphinxAtStartPar
The volume integral of a scalar field \( f(x, y, z) \) over a volume \( V \) is given by:
\begin{equation*}
\begin{split}
\iiint_V f(x, y, z) \, dV
\end{split}
\end{equation*}

\paragraph{Step 2: Switch to Cylindrical Coordinates}
\label{\detokenize{prerequisites:step-2-switch-to-cylindrical-coordinates}}
\sphinxAtStartPar
Since the region \( V \) is a cylinder, it’s convenient to switch to cylindrical coordinates:
\begin{itemize}
\item {} 
\sphinxAtStartPar
\( x = r \cos\theta \)

\item {} 
\sphinxAtStartPar
\( y = r \sin\theta \)

\item {} 
\sphinxAtStartPar
\( z = z \)

\end{itemize}

\sphinxAtStartPar
The volume element in cylindrical coordinates is:
\begin{equation*}
\begin{split}
dV = r \, dr \, d\theta \, dz
\end{split}
\end{equation*}
\sphinxAtStartPar
The scalar field \( f(x, y, z) = x + y + z \) becomes:
\begin{equation*}
\begin{split}
f(r, \theta, z) = r \cos\theta + r \sin\theta + z
\end{split}
\end{equation*}

\paragraph{Step 3: Set Up the Volume Integral}
\label{\detokenize{prerequisites:step-3-set-up-the-volume-integral}}
\sphinxAtStartPar
The region \( V \) is bounded by:
\begin{itemize}
\item {} 
\sphinxAtStartPar
\( r \in [0, 1] \) (the radius of the cylinder),

\item {} 
\sphinxAtStartPar
\( \theta \in [0, 2\pi] \) (the full revolution around the cylinder),

\item {} 
\sphinxAtStartPar
\( z \in [0, 3] \) (the height of the cylinder).

\end{itemize}

\sphinxAtStartPar
The volume integral in cylindrical coordinates is:
\begin{equation*}
\begin{split}
\iiint_V (r \cos\theta + r \sin\theta + z) \, r \, dr \, d\theta \, dz
\end{split}
\end{equation*}

\paragraph{Step 4: Solve the Volume Integral Numerically}
\label{\detokenize{prerequisites:step-4-solve-the-volume-integral-numerically}}
\sphinxAtStartPar
We can use Python to compute the volume integral:

\begin{sphinxuseclass}{cell}\begin{sphinxVerbatimInput}

\begin{sphinxuseclass}{cell_input}
\begin{sphinxVerbatim}[commandchars=\\\{\}]
\PYG{k+kn}{from} \PYG{n+nn}{scipy}\PYG{n+nn}{.}\PYG{n+nn}{integrate} \PYG{k+kn}{import} \PYG{n}{tplquad}
\PYG{k+kn}{import} \PYG{n+nn}{numpy} \PYG{k}{as} \PYG{n+nn}{np}

\PYG{c+c1}{\PYGZsh{} Define the scalar field in cylindrical coordinates f(r, theta, z)}
\PYG{k}{def} \PYG{n+nf}{integrand\PYGZus{}volume}\PYG{p}{(}\PYG{n}{z}\PYG{p}{,} \PYG{n}{r}\PYG{p}{,} \PYG{n}{theta}\PYG{p}{)}\PYG{p}{:}
    \PYG{k}{return} \PYG{n}{r} \PYG{o}{*} \PYG{p}{(}\PYG{n}{r} \PYG{o}{*} \PYG{n}{np}\PYG{o}{.}\PYG{n}{cos}\PYG{p}{(}\PYG{n}{theta}\PYG{p}{)} \PYG{o}{+} \PYG{n}{r} \PYG{o}{*} \PYG{n}{np}\PYG{o}{.}\PYG{n}{sin}\PYG{p}{(}\PYG{n}{theta}\PYG{p}{)} \PYG{o}{+} \PYG{n}{z}\PYG{p}{)}

\PYG{c+c1}{\PYGZsh{} Limits for the volume integral}
\PYG{n}{r\PYGZus{}min}\PYG{p}{,} \PYG{n}{r\PYGZus{}max} \PYG{o}{=} \PYG{l+m+mi}{0}\PYG{p}{,} \PYG{l+m+mi}{1}
\PYG{n}{theta\PYGZus{}min}\PYG{p}{,} \PYG{n}{theta\PYGZus{}max} \PYG{o}{=} \PYG{l+m+mi}{0}\PYG{p}{,} \PYG{l+m+mi}{2} \PYG{o}{*} \PYG{n}{np}\PYG{o}{.}\PYG{n}{pi}
\PYG{n}{z\PYGZus{}min}\PYG{p}{,} \PYG{n}{z\PYGZus{}max} \PYG{o}{=} \PYG{l+m+mi}{0}\PYG{p}{,} \PYG{l+m+mi}{3}

\PYG{c+c1}{\PYGZsh{} Compute the volume integral}
\PYG{n}{volume\PYGZus{}integral}\PYG{p}{,} \PYG{n}{\PYGZus{}} \PYG{o}{=} \PYG{n}{tplquad}\PYG{p}{(}\PYG{n}{integrand\PYGZus{}volume}\PYG{p}{,} \PYG{n}{z\PYGZus{}min}\PYG{p}{,} \PYG{n}{z\PYGZus{}max}\PYG{p}{,} \PYG{n}{r\PYGZus{}min}\PYG{p}{,} \PYG{n}{r\PYGZus{}max}\PYG{p}{,} \PYG{n}{theta\PYGZus{}min}\PYG{p}{,} \PYG{n}{theta\PYGZus{}max}\PYG{p}{)}
\PYG{n+nb}{print}\PYG{p}{(}\PYG{l+s+sa}{f}\PYG{l+s+s2}{\PYGZdq{}}\PYG{l+s+s2}{Volume integral: }\PYG{l+s+si}{\PYGZob{}}\PYG{n}{volume\PYGZus{}integral}\PYG{l+s+si}{:}\PYG{l+s+s2}{.2f}\PYG{l+s+si}{\PYGZcb{}}\PYG{l+s+s2}{\PYGZdq{}}\PYG{p}{)}
\end{sphinxVerbatim}

\end{sphinxuseclass}\end{sphinxVerbatimInput}
\begin{sphinxVerbatimOutput}

\begin{sphinxuseclass}{cell_output}
\begin{sphinxVerbatim}[commandchars=\\\{\}]
Volume integral: 34.07
\end{sphinxVerbatim}

\end{sphinxuseclass}\end{sphinxVerbatimOutput}

\end{sphinxuseclass}

\paragraph{Step 5: Final Answer}
\label{\detokenize{prerequisites:step-5-final-answer}}
\sphinxAtStartPar
After performing the integration, the result of the volume integral is:
\begin{equation*}
\begin{split}
\iiint_V (x + y + z) \, dV = 9\pi
\end{split}
\end{equation*}

\bigskip\hrule\bigskip



\section{Problem 11: Application of the Divergence Theorem and Stokes’ Theorem in Fluid Flow}
\label{\detokenize{prerequisites:problem-11-application-of-the-divergence-theorem-and-stokes-theorem-in-fluid-flow}}

\subsection{Problem Statement:}
\label{\detokenize{prerequisites:id13}}\begin{enumerate}
\sphinxsetlistlabels{\arabic}{enumi}{enumii}{}{.}%
\item {} 
\sphinxAtStartPar
\sphinxstylestrong{Divergence Theorem}: Given the vector field \( \mathbf{F}(x, y, z) = (x^2, y^2, z^2) \), use the \sphinxstylestrong{Divergence Theorem} to compute the flux of \( \mathbf{F} \) through the surface of the cube bounded by the planes \( x = 0 \), \( x = 1 \), \( y = 0 \), \( y = 1 \), \( z = 0 \), and \( z = 1 \).

\item {} 
\sphinxAtStartPar
\sphinxstylestrong{Stokes’ Theorem}: Given the vector field \( \mathbf{F}(x, y, z) = (-y, x, 0) \), use \sphinxstylestrong{Stokes’ Theorem} to compute the circulation of \( \mathbf{F} \) around the boundary of the surface defined by the disk \( x^2 + y^2 \leq 1 \) in the plane \( z = 0 \).

\end{enumerate}


\subsection{Solution}
\label{\detokenize{prerequisites:id14}}

\subsubsection{1. Divergence Theorem:}
\label{\detokenize{prerequisites:divergence-theorem}}
\sphinxAtStartPar
The \sphinxstylestrong{Divergence Theorem} relates the flux of a vector field through a closed surface to the volume integral of the divergence of the vector field over the region enclosed by the surface. Mathematically, it states:
\begin{equation*}
\begin{split}
\iint_S \mathbf{F} \cdot d\mathbf{A} = \iiint_V (\nabla \cdot \mathbf{F}) \, dV
\end{split}
\end{equation*}
\sphinxAtStartPar
where \( \mathbf{F} \) is the vector field, \( d\mathbf{A} \) is the surface area element, and \( \nabla \cdot \mathbf{F} \) is the divergence of \( \mathbf{F} \).

\sphinxAtStartPar
\sphinxstylestrong{Step 1}: Compute the divergence of the vector field \( \mathbf{F}(x, y, z) = (x^2, y^2, z^2) \):
\begin{equation*}
\begin{split}
\nabla \cdot \mathbf{F} = \frac{\partial}{\partial x}(x^2) + \frac{\partial}{\partial y}(y^2) + \frac{\partial}{\partial z}(z^2) = 2x + 2y + 2z
\end{split}
\end{equation*}
\sphinxAtStartPar
\sphinxstylestrong{Step 2}: Set up the volume integral over the cube \( 0 \leq x \leq 1 \), \( 0 \leq y \leq 1 \), and \( 0 \leq z \leq 1 \):
\begin{equation*}
\begin{split}
\iiint_V (2x + 2y + 2z) \, dV = \int_0^1 \int_0^1 \int_0^1 (2x + 2y + 2z) \, dx \, dy \, dz
\end{split}
\end{equation*}
\sphinxAtStartPar
\sphinxstylestrong{Step 3}: Compute the integral:
\begin{enumerate}
\sphinxsetlistlabels{\arabic}{enumi}{enumii}{}{.}%
\item {} 
\sphinxAtStartPar
Integrating with respect to \( x \):

\end{enumerate}
\begin{equation*}
\begin{split}
\int_0^1 (2x + 2y + 2z) \, dx = \left[ x^2 + 2yx + 2zx \right]_0^1 = 1 + 2y + 2z
\end{split}
\end{equation*}\begin{enumerate}
\sphinxsetlistlabels{\arabic}{enumi}{enumii}{}{.}%
\setcounter{enumi}{1}
\item {} 
\sphinxAtStartPar
Integrating with respect to \( y \):

\end{enumerate}
\begin{equation*}
\begin{split}
\int_0^1 (1 + 2y + 2z) \, dy = \left[ y + y^2 + 2zy \right]_0^1 = 1 + 1 + 2z = 2 + 2z
\end{split}
\end{equation*}\begin{enumerate}
\sphinxsetlistlabels{\arabic}{enumi}{enumii}{}{.}%
\setcounter{enumi}{2}
\item {} 
\sphinxAtStartPar
Finally, integrate with respect to \( z \):

\end{enumerate}
\begin{equation*}
\begin{split}
\int_0^1 (2 + 2z) \, dz = \left[ 2z + z^2 \right]_0^1 = 2 + 1 = 3
\end{split}
\end{equation*}
\sphinxAtStartPar
Thus, the flux of \( \mathbf{F} \) through the surface of the cube is:
\begin{equation*}
\begin{split}
\iint_S \mathbf{F} \cdot d\mathbf{A} = 3
\end{split}
\end{equation*}

\paragraph{Python}
\label{\detokenize{prerequisites:python}}
\sphinxAtStartPar
We will create a grid of points in the cube and define the vector field \( \mathbf{F}(x, y, z) = (x^2, y^2, z^2) \) on this grid.

\begin{sphinxuseclass}{cell}\begin{sphinxVerbatimInput}

\begin{sphinxuseclass}{cell_input}
\begin{sphinxVerbatim}[commandchars=\\\{\}]
\PYG{k+kn}{import} \PYG{n+nn}{numpy} \PYG{k}{as} \PYG{n+nn}{np}

\PYG{c+c1}{\PYGZsh{} Create a grid of points in the cube}
\PYG{n}{x\PYGZus{}vals}\PYG{p}{,} \PYG{n}{y\PYGZus{}vals}\PYG{p}{,} \PYG{n}{z\PYGZus{}vals} \PYG{o}{=} \PYG{n}{np}\PYG{o}{.}\PYG{n}{linspace}\PYG{p}{(}\PYG{l+m+mi}{0}\PYG{p}{,} \PYG{l+m+mi}{1}\PYG{p}{,} \PYG{l+m+mi}{5}\PYG{p}{)}\PYG{p}{,} \PYG{n}{np}\PYG{o}{.}\PYG{n}{linspace}\PYG{p}{(}\PYG{l+m+mi}{0}\PYG{p}{,} \PYG{l+m+mi}{1}\PYG{p}{,} \PYG{l+m+mi}{5}\PYG{p}{)}\PYG{p}{,} \PYG{n}{np}\PYG{o}{.}\PYG{n}{linspace}\PYG{p}{(}\PYG{l+m+mi}{0}\PYG{p}{,} \PYG{l+m+mi}{1}\PYG{p}{,} \PYG{l+m+mi}{5}\PYG{p}{)}
\PYG{n}{X}\PYG{p}{,} \PYG{n}{Y}\PYG{p}{,} \PYG{n}{Z} \PYG{o}{=} \PYG{n}{np}\PYG{o}{.}\PYG{n}{meshgrid}\PYG{p}{(}\PYG{n}{x\PYGZus{}vals}\PYG{p}{,} \PYG{n}{y\PYGZus{}vals}\PYG{p}{,} \PYG{n}{z\PYGZus{}vals}\PYG{p}{)}


\PYG{c+c1}{\PYGZsh{} Define the vector field F(x, y, z) = (x\PYGZca{}2, y\PYGZca{}2, z\PYGZca{}2)}
\PYG{c+c1}{\PYGZsh{} Vector field components}
\PYG{n}{F\PYGZus{}x} \PYG{o}{=} \PYG{n}{X}\PYG{o}{*}\PYG{o}{*}\PYG{l+m+mi}{2}  \PYG{c+c1}{\PYGZsh{} x component of the vector field}
\PYG{n}{F\PYGZus{}y} \PYG{o}{=} \PYG{n}{Y}\PYG{o}{*}\PYG{o}{*}\PYG{l+m+mi}{2}  \PYG{c+c1}{\PYGZsh{} y component of the vector field}
\PYG{n}{F\PYGZus{}z} \PYG{o}{=} \PYG{n}{Z}\PYG{o}{*}\PYG{o}{*}\PYG{l+m+mi}{2}  \PYG{c+c1}{\PYGZsh{} z component of the vector field}

\PYG{c+c1}{\PYGZsh{} Plot the vector field using quiver}
\PYG{n}{fig} \PYG{o}{=} \PYG{n}{plt}\PYG{o}{.}\PYG{n}{figure}\PYG{p}{(}\PYG{n}{figsize}\PYG{o}{=}\PYG{p}{(}\PYG{l+m+mi}{8}\PYG{p}{,} \PYG{l+m+mi}{8}\PYG{p}{)}\PYG{p}{)}
\PYG{n}{ax} \PYG{o}{=} \PYG{n}{fig}\PYG{o}{.}\PYG{n}{add\PYGZus{}subplot}\PYG{p}{(}\PYG{l+m+mi}{111}\PYG{p}{,} \PYG{n}{projection}\PYG{o}{=}\PYG{l+s+s1}{\PYGZsq{}}\PYG{l+s+s1}{3d}\PYG{l+s+s1}{\PYGZsq{}}\PYG{p}{)}
\PYG{n}{ax}\PYG{o}{.}\PYG{n}{quiver}\PYG{p}{(}\PYG{n}{X}\PYG{p}{,} \PYG{n}{Y}\PYG{p}{,} \PYG{n}{Z}\PYG{p}{,} \PYG{n}{F\PYGZus{}x}\PYG{p}{,} \PYG{n}{F\PYGZus{}y}\PYG{p}{,} \PYG{n}{F\PYGZus{}z}\PYG{p}{,} \PYG{n}{length}\PYG{o}{=}\PYG{l+m+mf}{0.1}\PYG{p}{)}

\PYG{c+c1}{\PYGZsh{} Set labels}
\PYG{n}{ax}\PYG{o}{.}\PYG{n}{set\PYGZus{}xlabel}\PYG{p}{(}\PYG{l+s+s1}{\PYGZsq{}}\PYG{l+s+s1}{X axis}\PYG{l+s+s1}{\PYGZsq{}}\PYG{p}{)}
\PYG{n}{ax}\PYG{o}{.}\PYG{n}{set\PYGZus{}ylabel}\PYG{p}{(}\PYG{l+s+s1}{\PYGZsq{}}\PYG{l+s+s1}{Y axis}\PYG{l+s+s1}{\PYGZsq{}}\PYG{p}{)}
\PYG{n}{ax}\PYG{o}{.}\PYG{n}{set\PYGZus{}zlabel}\PYG{p}{(}\PYG{l+s+s1}{\PYGZsq{}}\PYG{l+s+s1}{Z axis}\PYG{l+s+s1}{\PYGZsq{}}\PYG{p}{)}
\PYG{n}{ax}\PYG{o}{.}\PYG{n}{set\PYGZus{}title}\PYG{p}{(}\PYG{l+s+s1}{\PYGZsq{}}\PYG{l+s+s1}{Vector Field \PYGZdl{}}\PYG{l+s+se}{\PYGZbs{}\PYGZbs{}}\PYG{l+s+s1}{mathbf}\PYG{l+s+si}{\PYGZob{}F\PYGZcb{}}\PYG{l+s+s1}{(x, y, z) = (x\PYGZca{}2, y\PYGZca{}2, z\PYGZca{}2)\PYGZdl{} in the Cube}\PYG{l+s+s1}{\PYGZsq{}}\PYG{p}{)}

\PYG{n}{plt}\PYG{o}{.}\PYG{n}{show}\PYG{p}{(}\PYG{p}{)}
\end{sphinxVerbatim}

\end{sphinxuseclass}\end{sphinxVerbatimInput}
\begin{sphinxVerbatimOutput}

\begin{sphinxuseclass}{cell_output}
\noindent\sphinxincludegraphics{{4fffd4dc6d5e3baa3aa4aecb96bf07a24c08c48c2ae0b364b1ad46c884019ecf}.png}

\end{sphinxuseclass}\end{sphinxVerbatimOutput}

\end{sphinxuseclass}
\sphinxAtStartPar
Plot of the divergence field at \(Z=0\)

\begin{sphinxuseclass}{cell}\begin{sphinxVerbatimInput}

\begin{sphinxuseclass}{cell_input}
\begin{sphinxVerbatim}[commandchars=\\\{\}]
\PYG{k+kn}{import} \PYG{n+nn}{numpy} \PYG{k}{as} \PYG{n+nn}{np}
\PYG{k+kn}{import} \PYG{n+nn}{matplotlib}\PYG{n+nn}{.}\PYG{n+nn}{pyplot} \PYG{k}{as} \PYG{n+nn}{plt}
\PYG{k+kn}{from} \PYG{n+nn}{mpl\PYGZus{}toolkits}\PYG{n+nn}{.}\PYG{n+nn}{mplot3d} \PYG{k+kn}{import} \PYG{n}{Axes3D}

\PYG{c+c1}{\PYGZsh{} Create a grid of points in the cube}
\PYG{n}{n} \PYG{o}{=} \PYG{l+m+mi}{50}  \PYG{c+c1}{\PYGZsh{} Increase the grid density for a smoother surface plot}
\PYG{n}{x\PYGZus{}vals}\PYG{p}{,} \PYG{n}{y\PYGZus{}vals} \PYG{o}{=} \PYG{n}{np}\PYG{o}{.}\PYG{n}{linspace}\PYG{p}{(}\PYG{l+m+mi}{0}\PYG{p}{,} \PYG{l+m+mi}{1}\PYG{p}{,} \PYG{n}{n}\PYG{p}{)}\PYG{p}{,} \PYG{n}{np}\PYG{o}{.}\PYG{n}{linspace}\PYG{p}{(}\PYG{l+m+mi}{0}\PYG{p}{,} \PYG{l+m+mi}{1}\PYG{p}{,} \PYG{n}{n}\PYG{p}{)}
\PYG{n}{X}\PYG{p}{,} \PYG{n}{Y} \PYG{o}{=} \PYG{n}{np}\PYG{o}{.}\PYG{n}{meshgrid}\PYG{p}{(}\PYG{n}{x\PYGZus{}vals}\PYG{p}{,} \PYG{n}{y\PYGZus{}vals}\PYG{p}{)}

\PYG{c+c1}{\PYGZsh{} Calculate the divergence at z = 0 plane}
\PYG{n}{Divergence\PYGZus{}z0} \PYG{o}{=} \PYG{l+m+mi}{2} \PYG{o}{*} \PYG{n}{X} \PYG{o}{+} \PYG{l+m+mi}{2} \PYG{o}{*} \PYG{n}{Y}  \PYG{c+c1}{\PYGZsh{} Since z = 0, the divergence simplifies to 2x + 2y}

\PYG{c+c1}{\PYGZsh{} Create a 3D surface plot for the divergence on z = 0 plane}
\PYG{n}{fig} \PYG{o}{=} \PYG{n}{plt}\PYG{o}{.}\PYG{n}{figure}\PYG{p}{(}\PYG{n}{figsize}\PYG{o}{=}\PYG{p}{(}\PYG{l+m+mi}{8}\PYG{p}{,} \PYG{l+m+mi}{6}\PYG{p}{)}\PYG{p}{)}
\PYG{n}{ax} \PYG{o}{=} \PYG{n}{fig}\PYG{o}{.}\PYG{n}{add\PYGZus{}subplot}\PYG{p}{(}\PYG{l+m+mi}{111}\PYG{p}{,} \PYG{n}{projection}\PYG{o}{=}\PYG{l+s+s1}{\PYGZsq{}}\PYG{l+s+s1}{3d}\PYG{l+s+s1}{\PYGZsq{}}\PYG{p}{)}

\PYG{c+c1}{\PYGZsh{} Plot the surface}
\PYG{n}{surface\PYGZus{}plot} \PYG{o}{=} \PYG{n}{ax}\PYG{o}{.}\PYG{n}{plot\PYGZus{}surface}\PYG{p}{(}\PYG{n}{X}\PYG{p}{,} \PYG{n}{Y}\PYG{p}{,} \PYG{n}{Divergence\PYGZus{}z0}\PYG{p}{,} \PYG{n}{cmap}\PYG{o}{=}\PYG{l+s+s1}{\PYGZsq{}}\PYG{l+s+s1}{viridis}\PYG{l+s+s1}{\PYGZsq{}}\PYG{p}{,} \PYG{n}{edgecolor}\PYG{o}{=}\PYG{l+s+s1}{\PYGZsq{}}\PYG{l+s+s1}{none}\PYG{l+s+s1}{\PYGZsq{}}\PYG{p}{)}

\PYG{c+c1}{\PYGZsh{} Add a color bar}
\PYG{n}{fig}\PYG{o}{.}\PYG{n}{colorbar}\PYG{p}{(}\PYG{n}{surface\PYGZus{}plot}\PYG{p}{,} \PYG{n}{ax}\PYG{o}{=}\PYG{n}{ax}\PYG{p}{)}

\PYG{c+c1}{\PYGZsh{} Set labels}
\PYG{n}{ax}\PYG{o}{.}\PYG{n}{set\PYGZus{}xlabel}\PYG{p}{(}\PYG{l+s+s1}{\PYGZsq{}}\PYG{l+s+s1}{X axis}\PYG{l+s+s1}{\PYGZsq{}}\PYG{p}{)}
\PYG{n}{ax}\PYG{o}{.}\PYG{n}{set\PYGZus{}ylabel}\PYG{p}{(}\PYG{l+s+s1}{\PYGZsq{}}\PYG{l+s+s1}{Y axis}\PYG{l+s+s1}{\PYGZsq{}}\PYG{p}{)}
\PYG{n}{ax}\PYG{o}{.}\PYG{n}{set\PYGZus{}zlabel}\PYG{p}{(}\PYG{l+s+s1}{\PYGZsq{}}\PYG{l+s+s1}{Divergence}\PYG{l+s+s1}{\PYGZsq{}}\PYG{p}{)}
\PYG{n}{ax}\PYG{o}{.}\PYG{n}{set\PYGZus{}title}\PYG{p}{(}\PYG{l+s+s1}{\PYGZsq{}}\PYG{l+s+s1}{Divergence \PYGZdl{}}\PYG{l+s+se}{\PYGZbs{}\PYGZbs{}}\PYG{l+s+s1}{nabla }\PYG{l+s+se}{\PYGZbs{}\PYGZbs{}}\PYG{l+s+s1}{cdot }\PYG{l+s+se}{\PYGZbs{}\PYGZbs{}}\PYG{l+s+s1}{mathbf}\PYG{l+s+si}{\PYGZob{}F\PYGZcb{}}\PYG{l+s+s1}{ = 2x + 2y\PYGZdl{} on the \PYGZdl{}z = 0\PYGZdl{} plane}\PYG{l+s+s1}{\PYGZsq{}}\PYG{p}{)}

\PYG{n}{plt}\PYG{o}{.}\PYG{n}{show}\PYG{p}{(}\PYG{p}{)}
\end{sphinxVerbatim}

\end{sphinxuseclass}\end{sphinxVerbatimInput}
\begin{sphinxVerbatimOutput}

\begin{sphinxuseclass}{cell_output}
\noindent\sphinxincludegraphics{{b5edefacd9b4b1ff5e7e5b7976a56fe1e7fdf44aa419f5a9e3b69575075452d2}.png}

\end{sphinxuseclass}\end{sphinxVerbatimOutput}

\end{sphinxuseclass}
\sphinxAtStartPar
Volume integral for computing the divergence

\begin{sphinxuseclass}{cell}\begin{sphinxVerbatimInput}

\begin{sphinxuseclass}{cell_input}
\begin{sphinxVerbatim}[commandchars=\\\{\}]
\PYG{k+kn}{from} \PYG{n+nn}{scipy}\PYG{n+nn}{.}\PYG{n+nn}{integrate} \PYG{k+kn}{import} \PYG{n}{tplquad}

\PYG{c+c1}{\PYGZsh{} Define the divergence of the vector field}
\PYG{k}{def} \PYG{n+nf}{divergence}\PYG{p}{(}\PYG{n}{x}\PYG{p}{,} \PYG{n}{y}\PYG{p}{,} \PYG{n}{z}\PYG{p}{)}\PYG{p}{:}
    \PYG{k}{return} \PYG{l+m+mi}{2} \PYG{o}{*} \PYG{n}{x} \PYG{o}{+} \PYG{l+m+mi}{2} \PYG{o}{*} \PYG{n}{y} \PYG{o}{+} \PYG{l+m+mi}{2} \PYG{o}{*} \PYG{n}{z}

\PYG{c+c1}{\PYGZsh{} Limits of the cube: x, y, z in [0, 1]}
\PYG{n}{volume\PYGZus{}integral}\PYG{p}{,} \PYG{n}{\PYGZus{}} \PYG{o}{=} \PYG{n}{tplquad}\PYG{p}{(}\PYG{n}{divergence}\PYG{p}{,} \PYG{l+m+mi}{0}\PYG{p}{,} \PYG{l+m+mi}{1}\PYG{p}{,} \PYG{k}{lambda} \PYG{n}{x}\PYG{p}{:} \PYG{l+m+mi}{0}\PYG{p}{,} \PYG{k}{lambda} \PYG{n}{x}\PYG{p}{:} \PYG{l+m+mi}{1}\PYG{p}{,} \PYG{k}{lambda} \PYG{n}{x}\PYG{p}{,} \PYG{n}{y}\PYG{p}{:} \PYG{l+m+mi}{0}\PYG{p}{,} \PYG{k}{lambda} \PYG{n}{x}\PYG{p}{,} \PYG{n}{y}\PYG{p}{:} \PYG{l+m+mi}{1}\PYG{p}{)}
\PYG{n+nb}{print}\PYG{p}{(}\PYG{l+s+sa}{f}\PYG{l+s+s2}{\PYGZdq{}}\PYG{l+s+s2}{Volume integral of the divergence: }\PYG{l+s+si}{\PYGZob{}}\PYG{n}{volume\PYGZus{}integral}\PYG{l+s+si}{:}\PYG{l+s+s2}{.2f}\PYG{l+s+si}{\PYGZcb{}}\PYG{l+s+s2}{\PYGZdq{}}\PYG{p}{)}
\end{sphinxVerbatim}

\end{sphinxuseclass}\end{sphinxVerbatimInput}
\begin{sphinxVerbatimOutput}

\begin{sphinxuseclass}{cell_output}
\begin{sphinxVerbatim}[commandchars=\\\{\}]
Volume integral of the divergence: 3.00
\end{sphinxVerbatim}

\end{sphinxuseclass}\end{sphinxVerbatimOutput}

\end{sphinxuseclass}

\subsubsection{2. Stokes’ Theorem}
\label{\detokenize{prerequisites:stokes-theorem}}
\sphinxAtStartPar
\sphinxstylestrong{Stokes’ Theorem} relates the circulation of a vector field around a closed curve to the surface integral of the curl of the vector field over the surface bounded by the curve. Mathematically, it states:
\begin{equation*}
\begin{split}
\oint_C \mathbf{F} \cdot d\mathbf{r} = \iint_S (\nabla \times \mathbf{F}) \cdot d\mathbf{A}
\end{split}
\end{equation*}
\sphinxAtStartPar
where \( C \) is the boundary curve, and \( S \) is the surface bounded by \( C \).

\sphinxAtStartPar
\sphinxstylestrong{Step 1}: Compute the curl of the vector field \( \mathbf{F}(x, y, z) = (-y, x, 0) \):
\begin{equation*}
\begin{split}
\nabla \times \mathbf{F} = \left( \frac{\partial}{\partial y}(0) - \frac{\partial}{\partial z}(x), \frac{\partial}{\partial z}(-y) - \frac{\partial}{\partial x}(0), \frac{\partial}{\partial x}(x) - \frac{\partial}{\partial y}(-y) \right) = (0, 0, 2)
\end{split}
\end{equation*}
\sphinxAtStartPar
The curl of \( \mathbf{F} \) is \( (0, 0, 2) \).

\sphinxAtStartPar
\sphinxstylestrong{Step 2}: The surface is the disk \( x^2 + y^2 \leq 1 \) in the plane \( z = 0 \). The area element for a flat disk is \( dA = dx \, dy \).

\sphinxAtStartPar
Thus, the surface integral becomes:
\begin{equation*}
\begin{split}
\iint_S (0, 0, 2) \cdot (0, 0, 1) \, dA = 2 \iint_S dA
\end{split}
\end{equation*}
\sphinxAtStartPar
\sphinxstylestrong{Step 3}: The area of the disk is:
\begin{equation*}
\begin{split}
\iint_S dA = \pi (1^2) = \pi
\end{split}
\end{equation*}
\sphinxAtStartPar
Therefore, the circulation of \( \mathbf{F} \) around the boundary of the disk is:
\begin{equation*}
\begin{split}
\oint_C \mathbf{F} \cdot d\mathbf{r} = 2\pi
\end{split}
\end{equation*}

\subsubsection{Python}
\label{\detokenize{prerequisites:id15}}
\sphinxAtStartPar
Visualize the vector field

\begin{sphinxuseclass}{cell}\begin{sphinxVerbatimInput}

\begin{sphinxuseclass}{cell_input}
\begin{sphinxVerbatim}[commandchars=\\\{\}]
\PYG{k+kn}{import} \PYG{n+nn}{numpy} \PYG{k}{as} \PYG{n+nn}{np}
\PYG{k+kn}{import} \PYG{n+nn}{matplotlib}\PYG{n+nn}{.}\PYG{n+nn}{pyplot} \PYG{k}{as} \PYG{n+nn}{plt}

\PYG{c+c1}{\PYGZsh{} Define the grid for the disk x\PYGZca{}2 + y\PYGZca{}2 \PYGZlt{}= 1}
\PYG{n}{x\PYGZus{}vals} \PYG{o}{=} \PYG{n}{np}\PYG{o}{.}\PYG{n}{linspace}\PYG{p}{(}\PYG{o}{\PYGZhy{}}\PYG{l+m+mi}{1}\PYG{p}{,} \PYG{l+m+mi}{1}\PYG{p}{,} \PYG{l+m+mi}{20}\PYG{p}{)}
\PYG{n}{y\PYGZus{}vals} \PYG{o}{=} \PYG{n}{np}\PYG{o}{.}\PYG{n}{linspace}\PYG{p}{(}\PYG{o}{\PYGZhy{}}\PYG{l+m+mi}{1}\PYG{p}{,} \PYG{l+m+mi}{1}\PYG{p}{,} \PYG{l+m+mi}{20}\PYG{p}{)}
\PYG{n}{X}\PYG{p}{,} \PYG{n}{Y} \PYG{o}{=} \PYG{n}{np}\PYG{o}{.}\PYG{n}{meshgrid}\PYG{p}{(}\PYG{n}{x\PYGZus{}vals}\PYG{p}{,} \PYG{n}{y\PYGZus{}vals}\PYG{p}{)}

\PYG{c+c1}{\PYGZsh{} Mask the points that lie outside the disk (x\PYGZca{}2 + y\PYGZca{}2 \PYGZgt{} 1)}
\PYG{n}{mask} \PYG{o}{=} \PYG{n}{X}\PYG{o}{*}\PYG{o}{*}\PYG{l+m+mi}{2} \PYG{o}{+} \PYG{n}{Y}\PYG{o}{*}\PYG{o}{*}\PYG{l+m+mi}{2} \PYG{o}{\PYGZlt{}}\PYG{o}{=} \PYG{l+m+mi}{1}
\PYG{n}{X} \PYG{o}{=} \PYG{n}{X}\PYG{p}{[}\PYG{n}{mask}\PYG{p}{]}
\PYG{n}{Y} \PYG{o}{=} \PYG{n}{Y}\PYG{p}{[}\PYG{n}{mask}\PYG{p}{]}

\PYG{c+c1}{\PYGZsh{} Define the vector field F(x, y, z) = (\PYGZhy{}y, x, 0) at z = 0}
\PYG{n}{U} \PYG{o}{=} \PYG{o}{\PYGZhy{}}\PYG{n}{Y}  \PYG{c+c1}{\PYGZsh{} x\PYGZhy{}component of the vector field}
\PYG{n}{V} \PYG{o}{=} \PYG{n}{X}   \PYG{c+c1}{\PYGZsh{} y\PYGZhy{}component of the vector field}

\PYG{c+c1}{\PYGZsh{} Plot the vector field}
\PYG{n}{plt}\PYG{o}{.}\PYG{n}{figure}\PYG{p}{(}\PYG{n}{figsize}\PYG{o}{=}\PYG{p}{(}\PYG{l+m+mi}{6}\PYG{p}{,}\PYG{l+m+mi}{6}\PYG{p}{)}\PYG{p}{)}
\PYG{n}{plt}\PYG{o}{.}\PYG{n}{quiver}\PYG{p}{(}\PYG{n}{X}\PYG{p}{,} \PYG{n}{Y}\PYG{p}{,} \PYG{n}{U}\PYG{p}{,} \PYG{n}{V}\PYG{p}{)}
\PYG{n}{plt}\PYG{o}{.}\PYG{n}{title}\PYG{p}{(}\PYG{l+s+s2}{\PYGZdq{}}\PYG{l+s+s2}{Vector Field \PYGZdl{}}\PYG{l+s+se}{\PYGZbs{}\PYGZbs{}}\PYG{l+s+s2}{mathbf}\PYG{l+s+si}{\PYGZob{}F\PYGZcb{}}\PYG{l+s+s2}{(x, y, z) = (\PYGZhy{}y, x, 0)\PYGZdl{} in the plane \PYGZdl{}z = 0\PYGZdl{}}\PYG{l+s+s2}{\PYGZdq{}}\PYG{p}{)}
\PYG{n}{plt}\PYG{o}{.}\PYG{n}{xlabel}\PYG{p}{(}\PYG{l+s+s1}{\PYGZsq{}}\PYG{l+s+s1}{x}\PYG{l+s+s1}{\PYGZsq{}}\PYG{p}{)}
\PYG{n}{plt}\PYG{o}{.}\PYG{n}{ylabel}\PYG{p}{(}\PYG{l+s+s1}{\PYGZsq{}}\PYG{l+s+s1}{y}\PYG{l+s+s1}{\PYGZsq{}}\PYG{p}{)}
\PYG{n}{plt}\PYG{o}{.}\PYG{n}{grid}\PYG{p}{(}\PYG{k+kc}{True}\PYG{p}{)}
\PYG{n}{plt}\PYG{o}{.}\PYG{n}{axis}\PYG{p}{(}\PYG{l+s+s1}{\PYGZsq{}}\PYG{l+s+s1}{equal}\PYG{l+s+s1}{\PYGZsq{}}\PYG{p}{)}
\PYG{n}{plt}\PYG{o}{.}\PYG{n}{show}\PYG{p}{(}\PYG{p}{)}
\end{sphinxVerbatim}

\end{sphinxuseclass}\end{sphinxVerbatimInput}
\begin{sphinxVerbatimOutput}

\begin{sphinxuseclass}{cell_output}
\noindent\sphinxincludegraphics{{868f11922f683fc9c4ebefaaaf3048edb5d226932b890b8ceece1c6a5f4bb5f4}.png}

\end{sphinxuseclass}\end{sphinxVerbatimOutput}

\end{sphinxuseclass}
\sphinxAtStartPar
Plot the Boundary Curve:  The boundary curve of the disk is the unit circle

\begin{sphinxuseclass}{cell}\begin{sphinxVerbatimInput}

\begin{sphinxuseclass}{cell_input}
\begin{sphinxVerbatim}[commandchars=\\\{\}]
\PYG{c+c1}{\PYGZsh{} Define the boundary curve (circle) x\PYGZca{}2 + y\PYGZca{}2 = 1}
\PYG{n}{theta} \PYG{o}{=} \PYG{n}{np}\PYG{o}{.}\PYG{n}{linspace}\PYG{p}{(}\PYG{l+m+mi}{0}\PYG{p}{,} \PYG{l+m+mi}{2}\PYG{o}{*}\PYG{n}{np}\PYG{o}{.}\PYG{n}{pi}\PYG{p}{,} \PYG{l+m+mi}{100}\PYG{p}{)}
\PYG{n}{x\PYGZus{}circle} \PYG{o}{=} \PYG{n}{np}\PYG{o}{.}\PYG{n}{cos}\PYG{p}{(}\PYG{n}{theta}\PYG{p}{)}
\PYG{n}{y\PYGZus{}circle} \PYG{o}{=} \PYG{n}{np}\PYG{o}{.}\PYG{n}{sin}\PYG{p}{(}\PYG{n}{theta}\PYG{p}{)}

\PYG{c+c1}{\PYGZsh{} Plot the vector field and the boundary curve}
\PYG{n}{plt}\PYG{o}{.}\PYG{n}{figure}\PYG{p}{(}\PYG{n}{figsize}\PYG{o}{=}\PYG{p}{(}\PYG{l+m+mi}{6}\PYG{p}{,}\PYG{l+m+mi}{6}\PYG{p}{)}\PYG{p}{)}
\PYG{n}{plt}\PYG{o}{.}\PYG{n}{quiver}\PYG{p}{(}\PYG{n}{X}\PYG{p}{,} \PYG{n}{Y}\PYG{p}{,} \PYG{n}{U}\PYG{p}{,} \PYG{n}{V}\PYG{p}{,} \PYG{n}{color}\PYG{o}{=}\PYG{l+s+s1}{\PYGZsq{}}\PYG{l+s+s1}{blue}\PYG{l+s+s1}{\PYGZsq{}}\PYG{p}{,} \PYG{n}{alpha}\PYG{o}{=}\PYG{l+m+mf}{0.6}\PYG{p}{)}
\PYG{n}{plt}\PYG{o}{.}\PYG{n}{plot}\PYG{p}{(}\PYG{n}{x\PYGZus{}circle}\PYG{p}{,} \PYG{n}{y\PYGZus{}circle}\PYG{p}{,} \PYG{n}{color}\PYG{o}{=}\PYG{l+s+s1}{\PYGZsq{}}\PYG{l+s+s1}{red}\PYG{l+s+s1}{\PYGZsq{}}\PYG{p}{,} \PYG{n}{linewidth}\PYG{o}{=}\PYG{l+m+mi}{2}\PYG{p}{,} \PYG{n}{label}\PYG{o}{=}\PYG{l+s+s1}{\PYGZsq{}}\PYG{l+s+s1}{Boundary Curve}\PYG{l+s+s1}{\PYGZsq{}}\PYG{p}{)}
\PYG{n}{plt}\PYG{o}{.}\PYG{n}{title}\PYG{p}{(}\PYG{l+s+s2}{\PYGZdq{}}\PYG{l+s+s2}{Vector Field and Boundary Curve \PYGZdl{}x\PYGZca{}2 + y\PYGZca{}2 = 1\PYGZdl{}}\PYG{l+s+s2}{\PYGZdq{}}\PYG{p}{)}
\PYG{n}{plt}\PYG{o}{.}\PYG{n}{xlabel}\PYG{p}{(}\PYG{l+s+s1}{\PYGZsq{}}\PYG{l+s+s1}{x}\PYG{l+s+s1}{\PYGZsq{}}\PYG{p}{)}
\PYG{n}{plt}\PYG{o}{.}\PYG{n}{ylabel}\PYG{p}{(}\PYG{l+s+s1}{\PYGZsq{}}\PYG{l+s+s1}{y}\PYG{l+s+s1}{\PYGZsq{}}\PYG{p}{)}
\PYG{n}{plt}\PYG{o}{.}\PYG{n}{legend}\PYG{p}{(}\PYG{p}{)}
\PYG{n}{plt}\PYG{o}{.}\PYG{n}{grid}\PYG{p}{(}\PYG{k+kc}{True}\PYG{p}{)}
\PYG{n}{plt}\PYG{o}{.}\PYG{n}{axis}\PYG{p}{(}\PYG{l+s+s1}{\PYGZsq{}}\PYG{l+s+s1}{equal}\PYG{l+s+s1}{\PYGZsq{}}\PYG{p}{)}
\PYG{n}{plt}\PYG{o}{.}\PYG{n}{show}\PYG{p}{(}\PYG{p}{)}
\end{sphinxVerbatim}

\end{sphinxuseclass}\end{sphinxVerbatimInput}
\begin{sphinxVerbatimOutput}

\begin{sphinxuseclass}{cell_output}
\noindent\sphinxincludegraphics{{88ad401078b41b6d9f49647a05659ea7aaaaa72fbc3d1bd3d392ede8393f9c5a}.png}

\end{sphinxuseclass}\end{sphinxVerbatimOutput}

\end{sphinxuseclass}
\sphinxAtStartPar
Calculate the Curl

\begin{sphinxuseclass}{cell}\begin{sphinxVerbatimInput}

\begin{sphinxuseclass}{cell_input}
\begin{sphinxVerbatim}[commandchars=\\\{\}]
\PYG{c+c1}{\PYGZsh{} Area of the disk}
\PYG{n}{area\PYGZus{}of\PYGZus{}disk} \PYG{o}{=} \PYG{n}{np}\PYG{o}{.}\PYG{n}{pi} \PYG{o}{*} \PYG{l+m+mi}{1}\PYG{o}{*}\PYG{o}{*}\PYG{l+m+mi}{2}  \PYG{c+c1}{\PYGZsh{} radius = 1}
\PYG{n}{curl\PYGZus{}z} \PYG{o}{=} \PYG{l+m+mi}{2}  \PYG{c+c1}{\PYGZsh{} curl in the z direction}

\PYG{c+c1}{\PYGZsh{} Surface integral of the curl}
\PYG{n}{surface\PYGZus{}integral} \PYG{o}{=} \PYG{n}{curl\PYGZus{}z} \PYG{o}{*} \PYG{n}{area\PYGZus{}of\PYGZus{}disk}
\PYG{n+nb}{print}\PYG{p}{(}\PYG{l+s+sa}{f}\PYG{l+s+s2}{\PYGZdq{}}\PYG{l+s+s2}{Surface integral of the curl: }\PYG{l+s+si}{\PYGZob{}}\PYG{n}{surface\PYGZus{}integral}\PYG{l+s+si}{:}\PYG{l+s+s2}{.2f}\PYG{l+s+si}{\PYGZcb{}}\PYG{l+s+s2}{\PYGZdq{}}\PYG{p}{)}
\end{sphinxVerbatim}

\end{sphinxuseclass}\end{sphinxVerbatimInput}
\begin{sphinxVerbatimOutput}

\begin{sphinxuseclass}{cell_output}
\begin{sphinxVerbatim}[commandchars=\\\{\}]
Surface integral of the curl: 6.28
\end{sphinxVerbatim}

\end{sphinxuseclass}\end{sphinxVerbatimOutput}

\end{sphinxuseclass}

\section{Problem 12: Motion on an Inclined Plane}
\label{\detokenize{prerequisites:problem-12-motion-on-an-inclined-plane}}
\sphinxAtStartPar
A block of mass \( m = 5 \, \text{kg} \) is placed on a frictionless inclined plane, which makes an angle of \( \theta = 30^\circ \) with the horizontal. The block is released from rest at the top of the incline.
\begin{enumerate}
\sphinxsetlistlabels{\arabic}{enumi}{enumii}{}{.}%
\item {} 
\sphinxAtStartPar
What is the acceleration of the block along the incline?

\item {} 
\sphinxAtStartPar
How fast is the block moving after 2 seconds?

\item {} 
\sphinxAtStartPar
How far down the incline has the block moved after 2 seconds?

\end{enumerate}


\subsection{Solution}
\label{\detokenize{prerequisites:id16}}

\subsubsection{1. Acceleration Along the Incline}
\label{\detokenize{prerequisites:acceleration-along-the-incline}}
\sphinxAtStartPar
Using Newton’s second law of motion, the net force acting on the block along the incline is the component of the gravitational force that acts parallel to the incline. Since the plane is frictionless, the only force along the incline is due to gravity:
\begin{equation*}
\begin{split}
F_{\parallel} = mg \sin \theta
\end{split}
\end{equation*}
\sphinxAtStartPar
According to Newton’s second law, \( F = ma \), so the acceleration along the incline is:
\begin{equation*}
\begin{split}
a = \frac{F_{\parallel}}{m} = \frac{mg \sin \theta}{m} = g \sin \theta
\end{split}
\end{equation*}
\sphinxAtStartPar
Substituting the values:
\begin{equation*}
\begin{split}
a = (9.8 \, \text{m/s}^2) \sin 30^\circ = (9.8 \, \text{m/s}^2) \times 0.5 = 4.9 \, \text{m/s}^2
\end{split}
\end{equation*}
\sphinxAtStartPar
Thus, the acceleration of the block along the incline is:
\begin{equation*}
\begin{split}
a = 4.9 \, \text{m/s}^2
\end{split}
\end{equation*}

\subsubsection{2. Speed After 2 Seconds}
\label{\detokenize{prerequisites:speed-after-2-seconds}}
\sphinxAtStartPar
The block is released from rest, so its initial velocity is \( v_0 = 0 \, \text{m/s} \). Using the equation of motion:
\begin{equation*}
\begin{split}
v = v_0 + at
\end{split}
\end{equation*}
\sphinxAtStartPar
Substitute the known values:
\begin{equation*}
\begin{split}
v = 0 + (4.9 \, \text{m/s}^2) \times 2 \, \text{s} = 9.8 \, \text{m/s}
\end{split}
\end{equation*}
\sphinxAtStartPar
Thus, after 2 seconds, the speed of the block is:
\begin{equation*}
\begin{split}
v = 9.8 \, \text{m/s}
\end{split}
\end{equation*}

\subsubsection{3. Distance Traveled After 2 Seconds}
\label{\detokenize{prerequisites:distance-traveled-after-2-seconds}}
\sphinxAtStartPar
We can use the equation of motion:
\begin{equation*}
\begin{split}
s = v_0 t + \frac{1}{2} a t^2
\end{split}
\end{equation*}
\sphinxAtStartPar
Substitute the known values:
\begin{equation*}
\begin{split}
s = 0 \times 2 + \frac{1}{2} (4.9 \, \text{m/s}^2) (2 \, \text{s})^2 = \frac{1}{2} (4.9 \, \text{m/s}^2) \times 4 = 9.8 \, \text{m}
\end{split}
\end{equation*}
\sphinxAtStartPar
Thus, the block travels a distance of:
\begin{equation*}
\begin{split}
s = 9.8 \, \text{m}
\end{split}
\end{equation*}

\subsection{Summary of Results:}
\label{\detokenize{prerequisites:summary-of-results}}\begin{enumerate}
\sphinxsetlistlabels{\arabic}{enumi}{enumii}{}{.}%
\item {} 
\sphinxAtStartPar
The acceleration along the incline is \( 4.9 \, \text{m/s}^2 \).

\item {} 
\sphinxAtStartPar
The speed after 2 seconds is \( 9.8 \, \text{m/s} \).

\item {} 
\sphinxAtStartPar
The distance traveled after 2 seconds is \( 9.8 \, \text{m} \).

\end{enumerate}
\begin{align*}\!\begin{aligned}
\begin{tikzpicture}\\
% Inclined plane
\draw[thick] (0,0) -- (4,2);\\
% Block
\draw[fill=blue!30] (2,1) rectangle ++(0.5,0.5);\\
% Forces
\draw[->, red, thick] (2.25,1.25) -- ++(0,-2) node[below] {$mg$};
\draw[->, green, thick] (2.25,1.25) -- ++(1,0.5) node[right] {$N$};
\draw[->, blue, thick] (2.25,1.25) -- ++(1.5,-0.75) node[right] {$mg \sin \theta$};\\
% Angle label
\draw (0.7,0) arc[start angle=0,end angle=26.6,radius=0.7] node[midway, right] {$\theta$};\\
% Ground line
\draw[dashed] (0,0) -- (4,0);\\
\end{tikzpicture}\\
\end{aligned}\end{align*}

\bigskip\hrule\bigskip



\section{Problem 13: Object in Uniform Circular Motion with Angular Velocity}
\label{\detokenize{prerequisites:problem-13-object-in-uniform-circular-motion-with-angular-velocity}}
\sphinxAtStartPar
A car of mass \(m = 1000 \, \text{kg}\) is moving in a circular track of radius \(r = 50 \, \text{m}\) with a constant speed of \(v = 20 \, \text{m/s}\).
\begin{enumerate}
\sphinxsetlistlabels{\arabic}{enumi}{enumii}{}{.}%
\item {} 
\sphinxAtStartPar
What is the magnitude of the centripetal force acting on the car?

\item {} 
\sphinxAtStartPar
What is the angular velocity of the car?

\item {} 
\sphinxAtStartPar
If the coefficient of static friction between the tires and the track is \(\mu_s = 0.6\), what is the maximum speed the car can have before it starts to skid?

\end{enumerate}


\subsection{Solution}
\label{\detokenize{prerequisites:id17}}

\subsubsection{1. Magnitude of the Centripetal Force}
\label{\detokenize{prerequisites:magnitude-of-the-centripetal-force}}
\sphinxAtStartPar
In uniform circular motion, the centripetal force is the force that keeps an object moving in a circular path and is directed towards the center of the circle. The formula for centripetal force is:
\begin{equation*}
\begin{split}
F_c = \frac{mv^2}{r}
\end{split}
\end{equation*}
\sphinxAtStartPar
where:
\begin{itemize}
\item {} 
\sphinxAtStartPar
\(m = 1000 \, \text{kg}\) is the mass of the car,

\item {} 
\sphinxAtStartPar
\(v = 20 \, \text{m/s}\) is the speed of the car,

\item {} 
\sphinxAtStartPar
\(r = 50 \, \text{m}\) is the radius of the circular path.

\end{itemize}

\sphinxAtStartPar
Substituting the known values:
\begin{equation*}
\begin{split}
F_c = \frac{(1000 \, \text{kg}) \times (20 \, \text{m/s})^2}{50 \, \text{m}} = \frac{1000 \times 400}{50} = 8000 \, \text{N}
\end{split}
\end{equation*}
\sphinxAtStartPar
Thus, the magnitude of the centripetal force is:
\begin{equation*}
\begin{split}
F_c = 8000 \, \text{N}
\end{split}
\end{equation*}

\subsubsection{2. Angular Velocity of the Car}
\label{\detokenize{prerequisites:angular-velocity-of-the-car}}
\sphinxAtStartPar
The \sphinxstylestrong{angular velocity} \( \omega \) of an object in circular motion is related to its linear velocity \(v\) by the following relationship:
\begin{equation*}
\begin{split}
\omega = \frac{v}{r}
\end{split}
\end{equation*}
\sphinxAtStartPar
where:
\begin{itemize}
\item {} 
\sphinxAtStartPar
\(v = 20 \, \text{m/s}\) is the linear velocity of the car,

\item {} 
\sphinxAtStartPar
\(r = 50 \, \text{m}\) is the radius of the circular path.

\end{itemize}

\sphinxAtStartPar
Substituting the known values:
\begin{equation*}
\begin{split}
\omega = \frac{20 \, \text{m/s}}{50 \, \text{m}} = 0.4 \, \text{rad/s}
\end{split}
\end{equation*}
\sphinxAtStartPar
Thus, the angular velocity of the car is:
\begin{equation*}
\begin{split}
\omega = 0.4 \, \text{rad/s}
\end{split}
\end{equation*}

\subsubsection{3. Maximum Speed Before Skidding}
\label{\detokenize{prerequisites:maximum-speed-before-skidding}}
\sphinxAtStartPar
The car will start to skid when the frictional force between the tires and the track is no longer sufficient to provide the necessary centripetal force. The maximum frictional force (which acts as the centripetal force) is given by:
\begin{equation*}
\begin{split}
F_f = \mu_s F_n
\end{split}
\end{equation*}
\sphinxAtStartPar
where:
\begin{itemize}
\item {} 
\sphinxAtStartPar
\(\mu_s = 0.6\) is the coefficient of static friction,

\item {} 
\sphinxAtStartPar
\(F_n = mg\) is the normal force, where \(g = 9.8 \, \text{m/s}^2\) is the acceleration due to gravity.

\end{itemize}

\sphinxAtStartPar
Substituting the known values:
\begin{equation*}
\begin{split}
F_f = \mu_s mg = 0.6 \times 1000 \, \text{kg} \times 9.8 \, \text{m/s}^2 = 5880 \, \text{N}
\end{split}
\end{equation*}
\sphinxAtStartPar
The maximum centripetal force that friction can provide is \(F_f = 5880 \, \text{N}\). This force is also equal to the centripetal force required for circular motion, so we can set it equal to the formula for centripetal force and solve for the maximum speed \(v_\text{max}\):
\begin{equation*}
\begin{split}
F_f = \frac{m v_\text{max}^2}{r}
\end{split}
\end{equation*}
\sphinxAtStartPar
Substitute the values for \(F_f\), \(m\), and \(r\):
\begin{equation*}
\begin{split}
5880 \, \text{N} = \frac{1000 \, \text{kg} \times v_\text{max}^2}{50 \, \text{m}}
\end{split}
\end{equation*}
\sphinxAtStartPar
Solving for \(v_\text{max}\):
\begin{equation*}
\begin{split}
v_\text{max}^2 = \frac{5880 \times 50}{1000} = 294
\end{split}
\end{equation*}\begin{equation*}
\begin{split}
v_\text{max} = \sqrt{294} \approx 17.15 \, \text{m/s}
\end{split}
\end{equation*}
\sphinxAtStartPar
Thus, the maximum speed the car can have before skidding is:
\begin{equation*}
\begin{split}
v_\text{max} \approx 17.15 \, \text{m/s}
\end{split}
\end{equation*}

\subsection{Summary of Results}
\label{\detokenize{prerequisites:id18}}\begin{enumerate}
\sphinxsetlistlabels{\arabic}{enumi}{enumii}{}{.}%
\item {} 
\sphinxAtStartPar
The magnitude of the centripetal force acting on the car is \(8000 \, \text{N}\).

\item {} 
\sphinxAtStartPar
The angular velocity of the car is \(0.4 \, \text{rad/s}\).

\item {} 
\sphinxAtStartPar
The maximum speed the car can have before it starts to skid is approximately \(17.15 \, \text{m/s}\).

\end{enumerate}


\bigskip\hrule\bigskip



\section{Problem 14: Elastic Collision and Energy Conservation}
\label{\detokenize{prerequisites:problem-14-elastic-collision-and-energy-conservation}}
\sphinxAtStartPar
A ball of mass \(m_1 = 2 \, \text{kg}\) is moving with a velocity of \(v_1 = 5 \, \text{m/s}\) towards another stationary ball of mass \(m_2 = 3 \, \text{kg}\). The collision between the two balls is perfectly elastic.
\begin{enumerate}
\sphinxsetlistlabels{\arabic}{enumi}{enumii}{}{.}%
\item {} 
\sphinxAtStartPar
What are the velocities of the two balls after the collision?

\item {} 
\sphinxAtStartPar
Verify that both momentum and kinetic energy are conserved in the system.

\end{enumerate}


\subsection{Solution}
\label{\detokenize{prerequisites:id19}}

\subsubsection{1. Final Velocities of the Balls After the Collision}
\label{\detokenize{prerequisites:final-velocities-of-the-balls-after-the-collision}}
\sphinxAtStartPar
In a perfectly elastic collision, both \sphinxstylestrong{momentum} and \sphinxstylestrong{kinetic energy} are conserved.


\paragraph{Conservation of Momentum}
\label{\detokenize{prerequisites:conservation-of-momentum}}
\sphinxAtStartPar
The total momentum before the collision must equal the total momentum after the collision. The equation for conservation of momentum is:
\begin{equation*}
\begin{split}
m_1 v_1 + m_2 v_2 = m_1 v_1' + m_2 v_2'
\end{split}
\end{equation*}
\sphinxAtStartPar
Where:
\begin{itemize}
\item {} 
\sphinxAtStartPar
\(v_1' \) is the velocity of ball 1 after the collision,

\item {} 
\sphinxAtStartPar
\(v_2' \) is the velocity of ball 2 after the collision,

\item {} 
\sphinxAtStartPar
\(v_2 = 0\) since ball 2 is initially at rest.

\end{itemize}

\sphinxAtStartPar
Thus, the momentum conservation equation simplifies to:
\begin{equation*}
\begin{split}
m_1 v_1 = m_1 v_1' + m_2 v_2'
\end{split}
\end{equation*}
\sphinxAtStartPar
Substitute the known values for \(m_1\), \(v_1\), and \(m_2\):
\begin{equation*}
\begin{split}
2 \, \text{kg} \times 5 \, \text{m/s} = 2 \, \text{kg} \times v_1' + 3 \, \text{kg} \times v_2'
\end{split}
\end{equation*}\begin{equation*}
\begin{split}
10 \, \text{kg} \cdot \text{m/s} = 2 v_1' + 3 v_2'
\end{split}
\end{equation*}

\paragraph{Conservation of Kinetic Energy}
\label{\detokenize{prerequisites:conservation-of-kinetic-energy}}
\sphinxAtStartPar
The total kinetic energy before the collision must equal the total kinetic energy after the collision. The equation for conservation of kinetic energy is:
\begin{equation*}
\begin{split}
\frac{1}{2} m_1 v_1^2 + \frac{1}{2} m_2 v_2^2 = \frac{1}{2} m_1 v_1'^2 + \frac{1}{2} m_2 v_2'^2
\end{split}
\end{equation*}
\sphinxAtStartPar
Since \(v_2 = 0\), this simplifies to:
\begin{equation*}
\begin{split}
\frac{1}{2} m_1 v_1^2 = \frac{1}{2} m_1 v_1'^2 + \frac{1}{2} m_2 v_2'^2
\end{split}
\end{equation*}
\sphinxAtStartPar
Substitute the known values:
\begin{equation*}
\begin{split}
\frac{1}{2} \times 2 \, \text{kg} \times (5 \, \text{m/s})^2 = \frac{1}{2} \times 2 \, \text{kg} \times v_1'^2 + \frac{1}{2} \times 3 \, \text{kg} \times v_2'^2
\end{split}
\end{equation*}\begin{equation*}
\begin{split}
25 \, \text{J} = v_1'^2 + 1.5 v_2'^2
\end{split}
\end{equation*}
\sphinxAtStartPar
Now, we have two equations:
\begin{enumerate}
\sphinxsetlistlabels{\arabic}{enumi}{enumii}{}{.}%
\item {} 
\sphinxAtStartPar
\( 10 = 2 v_1' + 3 v_2' \)

\item {} 
\sphinxAtStartPar
\( 25 = v_1'^2 + 1.5 v_2'^2 \)

\end{enumerate}

\sphinxAtStartPar
We can solve this system of equations to find \(v_1'\) and \(v_2'\).


\paragraph{Solving the System of Equations}
\label{\detokenize{prerequisites:solving-the-system-of-equations}}
\sphinxAtStartPar
From the momentum equation:
\begin{equation*}
\begin{split}
v_1' = 5 - 1.5 v_2'
\end{split}
\end{equation*}
\sphinxAtStartPar
Substitute this into the energy equation:
\begin{equation*}
\begin{split}
25 = (5 - 1.5 v_2')^2 + 1.5 v_2'^2
\end{split}
\end{equation*}
\sphinxAtStartPar
Expand and simplify:
\begin{equation*}
\begin{split}
25 = (25 - 15 v_2' + 2.25 v_2'^2) + 1.5 v_2'^2
\end{split}
\end{equation*}\begin{equation*}
\begin{split}
25 = 25 - 15 v_2' + 3.75 v_2'^2
\end{split}
\end{equation*}
\sphinxAtStartPar
Simplifying further:
\begin{equation*}
\begin{split}
0 = -15 v_2' + 3.75 v_2'^2
\end{split}
\end{equation*}
\sphinxAtStartPar
Factor out \(v_2'\):
\begin{equation*}
\begin{split}
v_2' (3.75 v_2' - 15) = 0
\end{split}
\end{equation*}
\sphinxAtStartPar
Thus, \(v_2' = 0\) or \(v_2' = 4 \, \text{m/s}\).

\sphinxAtStartPar
If \(v_2' = 0\), the collision would not have been elastic, so \(v_2' = 4 \, \text{m/s}\).

\sphinxAtStartPar
Now, substitute \(v_2' = 4 \, \text{m/s}\) into the momentum equation:
\begin{equation*}
\begin{split}
v_1' = 5 - 1.5 \times 4 = 5 - 6 = -1 \, \text{m/s}
\end{split}
\end{equation*}
\sphinxAtStartPar
Thus, the velocities after the collision are:
\begin{itemize}
\item {} 
\sphinxAtStartPar
\(v_1' = -1 \, \text{m/s}\) (ball 1 reverses direction),

\item {} 
\sphinxAtStartPar
\(v_2' = 4 \, \text{m/s}\) (ball 2 moves forward).

\end{itemize}


\subsubsection{2. Verifying Conservation of Momentum and Energy}
\label{\detokenize{prerequisites:verifying-conservation-of-momentum-and-energy}}

\paragraph{Momentum Conservation}
\label{\detokenize{prerequisites:momentum-conservation}}
\sphinxAtStartPar
Before the collision:
\begin{equation*}
\begin{split}
p_\text{initial} = m_1 v_1 + m_2 v_2 = 2 \times 5 + 3 \times 0 = 10 \, \text{kg} \cdot \text{m/s}
\end{split}
\end{equation*}
\sphinxAtStartPar
After the collision:
\begin{equation*}
\begin{split}
p_\text{final} = m_1 v_1' + m_2 v_2' = 2 \times (-1) + 3 \times 4 = -2 + 12 = 10 \, \text{kg} \cdot \text{m/s}
\end{split}
\end{equation*}
\sphinxAtStartPar
Momentum is conserved, as \(p_\text{initial} = p_\text{final}\).


\paragraph{Energy Conservation}
\label{\detokenize{prerequisites:energy-conservation}}
\sphinxAtStartPar
Before the collision:
\begin{equation*}
\begin{split}
KE_\text{initial} = \frac{1}{2} m_1 v_1^2 = \frac{1}{2} \times 2 \times (5)^2 = 25 \, \text{J}
\end{split}
\end{equation*}
\sphinxAtStartPar
After the collision:
\begin{equation*}
\begin{split}
KE_\text{final} = \frac{1}{2} m_1 v_1'^2 + \frac{1}{2} m_2 v_2'^2 = \frac{1}{2} \times 2 \times (-1)^2 + \frac{1}{2} \times 3 \times (4)^2
\end{split}
\end{equation*}\begin{equation*}
\begin{split}
KE_\text{final} = 1 + 24 = 25 \, \text{J}
\end{split}
\end{equation*}
\sphinxAtStartPar
Energy is conserved, as \(KE_\text{initial} = KE_\text{final}\).


\subsection{Summary of Results}
\label{\detokenize{prerequisites:id20}}\begin{enumerate}
\sphinxsetlistlabels{\arabic}{enumi}{enumii}{}{.}%
\item {} 
\sphinxAtStartPar
After the collision, ball 1 has a velocity of \(v_1' = -1 \, \text{m/s}\) and ball 2 has a velocity of \(v_2' = 4 \, \text{m/s}\).

\item {} 
\sphinxAtStartPar
Both momentum and kinetic energy are conserved in this elastic collision.

\end{enumerate}


\bigskip\hrule\bigskip



\section{Problem 15: Object Falling from a Height (Energy Conservation)}
\label{\detokenize{prerequisites:problem-15-object-falling-from-a-height-energy-conservation}}
\sphinxAtStartPar
A rock of mass \(m = 2 \, \text{kg}\) is dropped from a height of \(h = 10 \, \text{m}\) above the ground. Assume no air resistance.
\begin{enumerate}
\sphinxsetlistlabels{\arabic}{enumi}{enumii}{}{.}%
\item {} 
\sphinxAtStartPar
What is the speed of the rock just before it hits the ground?

\item {} 
\sphinxAtStartPar
Verify that the potential energy at the start equals the kinetic energy just before impact.

\item {} 
\sphinxAtStartPar
How long does it take for the rock to hit the ground?

\end{enumerate}


\subsection{Solution}
\label{\detokenize{prerequisites:id21}}
\sphinxAtStartPar
This problem involves the \sphinxstylestrong{conversion of gravitational potential energy into kinetic energy} as the rock falls. We will also verify that energy is conserved in the system.


\subsubsection{1. Speed of the Rock Just Before Impact}
\label{\detokenize{prerequisites:speed-of-the-rock-just-before-impact}}
\sphinxAtStartPar
To find the speed of the rock just before it hits the ground, we use the \sphinxstylestrong{conservation of mechanical energy}. Initially, the rock has gravitational potential energy, and as it falls, this energy is converted into kinetic energy.

\sphinxAtStartPar
The total mechanical energy of the rock at the top (just before it is dropped) is:
\begin{equation*}
\begin{split}
E_\text{total} = PE_\text{initial} = mgh
\end{split}
\end{equation*}
\sphinxAtStartPar
Where:
\begin{itemize}
\item {} 
\sphinxAtStartPar
\(m = 2 \, \text{kg}\) is the mass of the rock,

\item {} 
\sphinxAtStartPar
\(g = 9.8 \, \text{m/s}^2\) is the acceleration due to gravity,

\item {} 
\sphinxAtStartPar
\(h = 10 \, \text{m}\) is the height from which the rock is dropped.

\end{itemize}

\sphinxAtStartPar
Substituting the known values:
\begin{equation*}
\begin{split}
PE_\text{initial} = 2 \, \text{kg} \times 9.8 \, \text{m/s}^2 \times 10 \, \text{m} = 196 \, \text{J}
\end{split}
\end{equation*}
\sphinxAtStartPar
At the moment just before the rock hits the ground, all the potential energy has been converted into kinetic energy. The total mechanical energy of the rock just before impact is:
\begin{equation*}
\begin{split}
E_\text{total} = KE_\text{final} = \frac{1}{2} m v^2
\end{split}
\end{equation*}
\sphinxAtStartPar
Since energy is conserved, \(PE_\text{initial} = KE_\text{final}\). So:
\begin{equation*}
\begin{split}
mgh = \frac{1}{2} m v^2
\end{split}
\end{equation*}
\sphinxAtStartPar
Cancel the mass \(m\) on both sides of the equation:
\begin{equation*}
\begin{split}
gh = \frac{1}{2} v^2
\end{split}
\end{equation*}
\sphinxAtStartPar
Solving for \(v\):
\begin{equation*}
\begin{split}
v^2 = 2gh
\end{split}
\end{equation*}
\sphinxAtStartPar
Substitute the values of \(g\) and \(h\):
\begin{equation*}
\begin{split}
v^2 = 2 \times 9.8 \, \text{m/s}^2 \times 10 \, \text{m} = 196 \, \text{m}^2/\text{s}^2
\end{split}
\end{equation*}
\sphinxAtStartPar
Taking the square root:
\begin{equation*}
\begin{split}
v = \sqrt{196} = 14 \, \text{m/s}
\end{split}
\end{equation*}
\sphinxAtStartPar
Thus, the speed of the rock just before it hits the ground is:
\begin{equation*}
\begin{split}
v = 14 \, \text{m/s}
\end{split}
\end{equation*}

\subsubsection{2. Verify Energy Conservation}
\label{\detokenize{prerequisites:verify-energy-conservation}}
\sphinxAtStartPar
Initially, the rock has only \sphinxstylestrong{potential energy}:
\begin{equation*}
\begin{split}
PE_\text{initial} = mgh = 196 \, \text{J}
\end{split}
\end{equation*}
\sphinxAtStartPar
Just before the rock hits the ground, all this potential energy is converted into \sphinxstylestrong{kinetic energy}. The final kinetic energy is:
\begin{equation*}
\begin{split}
KE_\text{final} = \frac{1}{2} m v^2 = \frac{1}{2} \times 2 \, \text{kg} \times (14 \, \text{m/s})^2 = \frac{1}{2} \times 2 \times 196 = 196 \, \text{J}
\end{split}
\end{equation*}
\sphinxAtStartPar
Since \(PE_\text{initial} = KE_\text{final}\), energy is conserved in the system.


\subsubsection{3. Time to Hit the Ground}
\label{\detokenize{prerequisites:time-to-hit-the-ground}}
\sphinxAtStartPar
To find how long it takes for the rock to hit the ground, we use the equation for vertical motion under constant acceleration:
\begin{equation*}
\begin{split}
h = \frac{1}{2} g t^2
\end{split}
\end{equation*}
\sphinxAtStartPar
Solving for \(t\):
\begin{equation*}
\begin{split}
t^2 = \frac{2h}{g}
\end{split}
\end{equation*}
\sphinxAtStartPar
Substitute the values for \(h\) and \(g\):
\begin{equation*}
\begin{split}
t^2 = \frac{2 \times 10 \, \text{m}}{9.8 \, \text{m/s}^2} \approx 2.04 \, \text{s}^2
\end{split}
\end{equation*}
\sphinxAtStartPar
Taking the square root:
\begin{equation*}
\begin{split}
t = \sqrt{2.04} \approx 1.43 \, \text{s}
\end{split}
\end{equation*}
\sphinxAtStartPar
Thus, it takes approximately:
\begin{equation*}
\begin{split}
t \approx 1.43 \, \text{s}
\end{split}
\end{equation*}
\sphinxAtStartPar
for the rock to hit the ground.


\subsection{Summary of Results}
\label{\detokenize{prerequisites:id22}}\begin{enumerate}
\sphinxsetlistlabels{\arabic}{enumi}{enumii}{}{.}%
\item {} 
\sphinxAtStartPar
The speed of the rock just before it hits the ground is \(14 \, \text{m/s}\).

\item {} 
\sphinxAtStartPar
The potential energy at the start equals the kinetic energy just before impact, confirming that energy is conserved.

\item {} 
\sphinxAtStartPar
It takes approximately \(1.43 \, \text{seconds}\) for the rock to hit the ground.

\end{enumerate}

\sphinxstepscope


\chapter{Bibliography}
\label{\detokenize{bibliography:bibliography}}\label{\detokenize{bibliography::doc}}
\begin{sphinxthebibliography}{MYOH12}
\bibitem[CC18]{bibliography:id2}
\sphinxAtStartPar
Yunus A. Cengel and John M. Cimbala. \sphinxstyleemphasis{Fluid Mechanics: Fundamentals and Applications}. McGraw\sphinxhyphen{}Hill, 2018.
\bibitem[FMP11]{bibliography:id6}
\sphinxAtStartPar
Robert W. Fox, Alan T. McDonald, and Philip J. Pritchard. \sphinxstyleemphasis{Introduction to Fluid Mechanics}. Wiley, 2011.
\bibitem[MYOH12]{bibliography:id4}
\sphinxAtStartPar
Bruce R. Munson, Donald F. Young, Theodore H. Okiishi, and Wade W. Huebsch. \sphinxstyleemphasis{Fundamentals of Fluid Mechanics}. Wiley, 2012.
\bibitem[Sha03]{bibliography:id5}
\sphinxAtStartPar
Irving H. Shames. \sphinxstyleemphasis{Mechanics of Fluids}. McGraw\sphinxhyphen{}Hill, 2003.
\bibitem[Whi16]{bibliography:id3}
\sphinxAtStartPar
Frank M. White. \sphinxstyleemphasis{Fluid Mechanics}. McGraw\sphinxhyphen{}Hill, 2016.
\end{sphinxthebibliography}







\renewcommand{\indexname}{Index}
\printindex
\end{document}